\chapter{Aljabar linear} \label{linalg:appendix}

%%%%%%%%%%%%%%%%%%%%%%%%%%%%%%%%%%%%%%%%%%%%%%%%%%%%%%%%%%%%%%%%%%%%%%%%%%%%%%

\section{Vektor, pemetaan, dan matriks}
\label{vecsandmaps:section}

%mbxINTROSUBSECTION

\sectionnotes{2 kuliah}

Dalam kehidupan nyata, paling sering terdapat lebih dari satu variabel.
Kita ingin menata penanganan banyak variabel secara konsisten,
dan khususnya menata penanganan persamaan linear dan pemetaan
linear, karena keduanya sangat berguna dan cukup mudah ditangani.
Para matematikawan bergurau bahwa
\myquote{bagi seorang insinyur, setiap masalah itu linear, dan segala sesuatu adalah
matriks.}
Dan memang, mereka (para insinyur) tidak keliru.
Cukup sering, menyelesaikan suatu masalah rekayasa berarti menentukan
masalah linear berdimensi hingga yang tepat untuk diselesaikan, yang kemudian
diselesaikan melalui sejumlah manipulasi matriks.
Yang paling penting, masalah linear adalah masalah yang kita ketahui cara menyelesaikannya,
dan kita memiliki banyak alat untuk menyelesaikannya.
Bagi insinyur, matematikawan, fisikawan, dan siapa pun dalam bidang teknis,
mempelajari aljabar linear mutlak sangat penting.

Sebagai motivasi, andaikan kita ingin menyelesaikan
\begin{equation*}
\begin{aligned}
& x-y = 2 , \\
& 2x+y = 4 ,
\end{aligned}
\end{equation*}
untuk $x$ dan $y$.
Artinya, kita menginginkan bilangan $x$ dan $y$ sedemikian sehingga kedua
persamaan tersebut dipenuhi.
Kita dapat memulai dengan menjumlahkan kedua persamaan untuk memperoleh
\begin{equation*}
x+2x-y+y = 2+4, \qquad \text{atau} \qquad 3x = 6 .
\end{equation*}
Dengan kata lain, $x=2$. Setelah memperolehnya, kita substitusikan $x=2$ ke dalam
persamaan pertama untuk mendapatkan $2-y=2$, sehingga $y=0$. Baiklah\@, itu mudah. Mengapa
persamaan linear begitu diributkan? Nah, coba lakukan ini jika Anda memiliki
5000 peubah\footnote{Salah satu kelemahan membuat segala sesuatu tampak seperti
masalah linear ialah jumlah variabel cenderung menjadi sangat besar.}.
Selain itu, kita mungkin memiliki persamaan semacam itu bukan hanya untuk bilangan,
tetapi juga untuk fungsi dan turunan fungsi dalam persamaan diferensial.
Jelas kita memerlukan cara yang sistematis untuk mengerjakannya.
Akibat yang menyenangkan dari membuat segala sesuatu sistematis dan lebih sederhana untuk dituliskan
adalah bahwa komputer menjadi lebih mudah melakukan pekerjaan itu untuk kita.
Komputer sangat mahir mengerjakan banyak tugas berulang
secara tepat, asalkan kita menemukan cara sistematis bagi komputer untuk
melaksanakan tugas-tugas tersebut.

\subsection{Vektor dan operasi pada vektor}

Pandang $n$ bilangan real sebagai suatu tupel-$n$:
\begin{equation*}
(x_1,x_2,\ldots,x_n).
\end{equation*}
Himpunan tupel-$n$ semacam itu adalah apa yang disebut
\emph{ruang berdimensi $n$}\index{ruang berdimensi $n$},
yang sering dinotasikan dengan ${\mathbb R}^n$.
Kadang-kadang kita menyebutnya
\emph{\myindex{ruang Euclid}} berdimensi $n$%
\footnote{Dinamai menurut matematikawan Yunani kuno
\href{https://en.wikipedia.org/wiki/Euclid}{Euclid dari Aleksandria}
(sekitar 300 SM), mungkin matematikawan yang paling terkenal; bahkan
kota-kota kecil sering memiliki Jalan Euclid atau Adimarga Euclid.}.
Dalam dua dimensi, ${\mathbb R}^2$ disebut
\emph{\myindex{bidang Kartesius}}%
\footnote{Dinamai menurut matematikawan Prancis
\href{https://en.wikipedia.org/wiki/Descartes}{Ren\'e Descartes}
(1596--1650). Bidang itu disebut \myquote{Kartesius} karena namanya dalam bahasa Latin adalah Renatus
Cartesius.}.
Setiap tupel-$n$ semacam itu merepresentasikan suatu titik dalam ruang berdimensi $n$.
Sebagai contoh, titik
$(1,2)$ dalam bidang ${\mathbb R}^2$
berada satu satuan ke kanan dan dua satuan ke atas dari
titik asal.

Ketika kita melakukan aljabar dengan tupel-$n$ bilangan ini, kita menyebutnya
\emph{vektor}\index{vektor}\footnote{%
Notasi yang lazim untuk membedakan vektor dari titik adalah menulis $(1,2)$
untuk titik dan $\langle 1,2 \rangle$ untuk vektor. Kita menulis keduanya sebagai
$(1,2)$.}. Para matematikawan suka membedakan
apa yang merupakan vektor dan apa yang merupakan titik dalam ruang atau bidang,
dan ternyata
pembedaan itu penting; namun, untuk keperluan aljabar linear,
kita dapat memandang segala sesuatu direpresentasikan oleh suatu vektor.
Salah satu cara memandang vektor, yang sangat berguna dalam kalkulus
dan persamaan diferensial, adalah sebagai anak panah. Vektor adalah objek yang memiliki
\emph{\myindex{arah}} dan \emph{besar}.
Sebagai contoh, vektor $(1,2)$
adalah anak panah dari titik asal ke titik $(1,2)$ dalam bidang.
Besarnya adalah panjang anak panah tersebut.
Lihat \figurevref{linalg-vecarrow:fig}.
Jika kita memandang vektor sebagai anak panah,
anak panahnya tidak harus selalu berawal di titik asal. Namun, jika kita
memindahkannya, anak panah itu harus selalu mempertahankan arah dan besar yang sama.

\begin{myfig}
\capstart
\inputpdft{linalg-vecarrow}{%
Diagram bidang x sub 1 x sub 2, dengan sebuah anak panah digambar dari titik
(0,0) ke titik (1,2).}
\caption{Vektor $(1,2)$ digambar sebagai anak panah dari titik asal ke titik
$(1,2)$.\label{linalg-vecarrow:fig}}
\end{myfig}

Karena vektor merupakan anak panah, ketika hendak memberi nama suatu vektor,
kita menggambar anak panah kecil di atasnya:
\begin{equation*}
\vec{x}
\end{equation*}
Notasi populer lainnya adalah $\mathbf{x}$ bercetak tebal, walaupun kita akan menggunakan anak panah
kecil. Menulis huruf tebal dalam buku mungkin mudah, tetapi tidak begitu
mudah menuliskannya dengan tangan di atas kertas atau papan tulis.
Para matematikawan bahkan sering tidak menulis anak panahnya. Seorang matematikawan akan
menulis $x$ dan mengingat bahwa $x$ adalah vektor, bukan bilangan.
Sama seperti Anda mengingat bahwa Bob adalah paman Anda dan tidak perlu
terus-menerus mengulang \myquote{Paman Bob}. Anda cukup mengatakan \myquote{Bob.}
Namun, dalam buku ini kita akan memanggil Bob \myquote{Paman Bob}
dan menuliskan vektor dengan anak panah kecil.

\emph{\myindex{Besar}} dapat dihitung menggunakan Teorema Pythagoras.
Vektor $(1,2)$ yang digambar dalam gambar memiliki besar $\sqrt{1^2+2^2} =
\sqrt{5}$. Besar tersebut dinotasikan dengan $\lVert \vec{x} \rVert$,
dan, dalam berapa pun jumlah dimensinya, dapat dihitung dengan cara yang sama:
\begin{equation*}
\lVert \vec{x} \rVert
=
\lVert (x_1,x_2,\ldots,x_n) \rVert
=
\sqrt{x_1^2+x_2^2+\cdots+x_n^2} .
\end{equation*}

Untuk alasan yang akan menjadi jelas dalam bagian berikutnya, kita sering
menuliskan vektor sebagai apa yang disebut
\emph{vektor kolom}\index{vektor kolom}:
\begin{equation*}
\vec{x} =
\begin{bmatrix}
x_{1} \\ x_2 \\ \vdots \\ x_n
\end{bmatrix} .
\end{equation*}
Jangan khawatir. Itu hanya cara lain untuk menuliskan hal yang sama.
Sebagai contoh, vektor $(1,2)$ dapat ditulis sebagai
\begin{equation*}
\begin{bmatrix}
1 \\ 2
\end{bmatrix} .
\end{equation*}

Fakta bahwa kita menuliskan anak panah di atas vektor memungkinkan kita menulis beberapa
vektor $\vec{x}_1$, $\vec{x}_2$, dan seterusnya, tanpa mengacaukannya dengan
komponen-komponen suatu vektor lain $\vec{x}$.

Lalu, di manakah \emph{aljabar} dalam \emph{aljabar linear}?
Nah, anak panah dapat dijumlahkan, dikurangkan,
dan dikalikan dengan bilangan.
Pertama-tama kita membahas \emph{penjumlahan}\index{menjumlahkan vektor}.
Jika kita memiliki dua anak panah, kita cukup
bergerak sepanjang yang satu, lalu sepanjang yang lain. Lihat
\figurevref{linalg-vecadd:fig}.

\begin{myfig}
\capstart
\inputpdft{linalg-vecadd}{%
Diagram bidang x sub 1 x sub 2. Sebuah anak panah bertitik-titik digambar dari
titik (0,0) ke titik (1,2). Sebuah anak panah putus-putus
digambar dari titik (1,2) ke titik (3, minus 1), yakni berawal dari
ujung anak panah bertitik-titik. Sebuah anak panah utuh digambar dari (0,0)
ke titik (3, minus 1), yakni berawal di titik asal dan berakhir di
ujung anak panah putus-putus.}
\caption{Menjumlahkan vektor $(1,2)$, yang digambar bertitik-titik, dan $(2,-3)$, yang digambar putus-putus.
Hasilnya, $(3,-1)$, digambar sebagai anak panah utuh.\label{linalg-vecadd:fig}}
\end{myfig}

Cukup mudah melihat pengaruhnya terhadap bilangan-bilangan yang merepresentasikan
vektor. Andaikan kita ingin menjumlahkan $(1,2)$ dengan $(2,-3)$ seperti pada gambar.
Kita bergerak sepanjang $(1,2)$
dan kemudian bergerak sepanjang $(2,-3)$.
Yang kita lakukan adalah bergerak satu satuan ke kanan, dua satuan ke atas, lalu
bergerak dua satuan ke kanan dan tiga satuan ke bawah (minus tiga). Ini
berarti kita berakhir di $\bigl(1+2,2+(-3)\bigr) = (3,-1)$.
Begitulah penjumlahan selalu bekerja:
\begin{equation*}
\begin{bmatrix}
x_{1} \\ x_2 \\ \vdots \\ x_n
\end{bmatrix} +
\begin{bmatrix}
y_{1} \\ y_2 \\ \vdots \\ y_n
\end{bmatrix} =
\begin{bmatrix}
x_1 + y_{1} \\ x_2+ y_2 \\ \vdots \\ x_n + y_n
\end{bmatrix} .
\end{equation*}

\emph{Pengurangan}\index{mengurangkan vektor} bekerja secara serupa.
Makna visual dari $\vec{x}- \vec{y}$ adalah bahwa
mula-mula kita bergerak sepanjang $\vec{x}$, kemudian bergerak
mundur sepanjang $\vec{y}$.
Lihat \figurevref{linalg-vecsub:fig}.
Ini seperti menjumlahkan
$\vec{x}+ (- \vec{y})$, dengan $-\vec{y}$
adalah anak panah yang kita peroleh dengan menghapus kepala anak panah
dari satu sisi dan menggambarnya di sisi lain.
Artinya, kita membalik
arahnya. Ditinjau dari bilangannya, kita cukup bergerak mundur baik
secara horizontal maupun vertikal,
sehingga kita menegatifkan kedua bilangan. Misalnya, jika $\vec{y}$ adalah $(-2,1)$,
maka $-\vec{y}$ adalah $(2,-1)$.

\begin{myfig}
\capstart
\inputpdft{linalg-vecsub}{%
Diagram bidang x sub 1 x sub 2. Sebuah anak panah bertitik-titik digambar dari
titik (0,0) ke titik (1,2). Sebuah anak panah putus-putus
digambar dari titik (3,1) ke titik (1,2), yakni berawal dari
titik (3,1) dan berakhir di ujung anak panah bertitik-titik.
Sebuah anak panah utuh digambar dari (0,0)
ke titik (3,1), yakni berawal di titik asal dan berakhir di
ekor anak panah putus-putus.}
\caption{Pengurangan: vektor $(1,2)$, yang digambar bertitik-titik, dikurangi $(-2,1)$,
yang digambar putus-putus. Hasilnya,
$(3,1)$, digambar sebagai anak panah utuh.\label{linalg-vecsub:fig}}
\end{myfig}

Hal intuitif lain yang dapat dilakukan terhadap vektor adalah
\emph{menskalakannya}\index{menskalakan vektor}.
Kita merepresentasikan ini dengan perkalian suatu bilangan dengan vektor.
Karena itu, ketika ingin membedakan antara vektor dan bilangan, kita
menyebut bilangan sebagai \emph{skalar}\index{skalar}.
Sebagai contoh,
andaikan kita ingin bergerak tiga kali lebih jauh. Jika vektornya $(1,2)$,
bergerak 3 kali lebih jauh berarti bergerak 3 satuan ke kanan dan 6 satuan ke atas,
sehingga kita memperoleh vektor $(3,6)$.
Kita cukup mengalikan setiap bilangan dalam vektor dengan 3.
Jika $\alpha$ adalah suatu bilangan, maka
\begin{equation*}
\alpha
\begin{bmatrix}
x_{1} \\ x_2 \\ \vdots \\ x_n
\end{bmatrix} =
\begin{bmatrix}
\alpha x_{1} \\ \alpha x_2 \\ \vdots \\ \alpha x_n
\end{bmatrix} .
\end{equation*}
Penskalaan (dengan bilangan positif) mengalikan besarnya
dan tidak mengubah arahnya.
Besar $(1,2)$
adalah $\sqrt{5}$. Besar 3 kali $(1,2)$, yakni besar $(3,6)$, adalah
$3\sqrt{5}$.

Jika kita mengalikan suatu vektor dengan skalar negatif,
vektor tersebut tidak hanya diskalakan, tetapi juga berbalik arah.
Mengalikan $(1,2)$ dengan $-3$ berarti kita harus bergerak 3 kali lebih jauh tetapi dalam
arah yang berlawanan, jadi 3 satuan ke kiri dan 6 satuan ke bawah, atau dengan kata
lain, $(-3,-6)$. Seperti telah disebutkan di atas, $-\vec{y}$ adalah kebalikan dari
$\vec{y}$, dan ini sama dengan $(-1)\vec{y}$.

Dalam \figurevref{linalg-vecscale:fig}, Anda dapat melihat beberapa contoh
makna visual penskalaan suatu vektor.

\begin{myfig}
\capstart
\inputpdft{linalg-vecscale}{%
Tiga diagram.
Pertama, vektor x digambar sebagai anak panah.
Berikutnya, vektor 2 x digambar
sebagai anak panah yang dua kali lebih panjang, dengan dua versi acuan x digambar
di sampingnya sebagai pembanding skala. Terakhir, vektor minus 1.5 x digambar;
vektor ini mengarah ke arah yang berlawanan dan panjangnya 1.5 kali panjang
x. Sekali lagi, dua versi acuan x digambar di sampingnya sebagai pembanding skala.}
\caption{Suatu vektor $\vec{x}$, vektor $2\vec{x}$ (arah sama,
besar dua kali lipat), dan vektor $-1.5\vec{x}$ (arah berlawanan,
besar 1.5 kali lipat).\label{linalg-vecscale:fig}}
\end{myfig}

Kita menggabungkan semua operasi ini untuk mengerjakan ungkapan yang lebih rumit.
Mari kita hitung sebuah contoh kecil:
\begin{equation*}
3
\begin{bmatrix}
1 \\ 2
\end{bmatrix}
+
2
\begin{bmatrix}
-4 \\ -1
\end{bmatrix}
-
3
\begin{bmatrix}
-2 \\ 2
\end{bmatrix}
=
\begin{bmatrix}
3(1)+2(-4)-3(-2) \\ 3(2)+2(-1)-3(2)
\end{bmatrix}
=
\begin{bmatrix}
1 \\ -2
\end{bmatrix}
.
\end{equation*}

\medskip

Kita mengatakan bahwa vektor memiliki arah dan besar. Besar mudah
direpresentasikan---hanyalah suatu bilangan. \emph{\myindex{Arah}} biasanya
diberikan oleh vektor dengan besar satu. Vektor semacam itu kita sebut
\emph{\myindex{vektor satuan}}. Artinya, $\vec{u}$ adalah vektor satuan ketika
$\lVert \vec{u} \rVert = 1$. Sebagai contoh, vektor $(1,0)$,
$(\nicefrac{1}{\sqrt{2}},\nicefrac{1}{\sqrt{2}})$, dan $(0,-1)$ semuanya merupakan
vektor satuan.

Untuk merepresentasikan arah suatu vektor $\vec{x}$, kita perlu mencari
vektor satuan dalam arah yang sama. Untuk melakukannya, kita cukup menskalakan ulang
$\vec{x}$ dengan kebalikan besarnya:
$\frac{1}{\lVert \vec{x} \rVert} \vec{x}$, atau secara lebih ringkas,
$\frac{\vec{x}}{\lVert \vec{x} \rVert}$.

Sebagai contoh, vektor satuan searah dengan $(1,2)$ adalah
\begin{equation*}
\frac{1}{\sqrt{1^2+2^2}} (1,2)
=
\left( \frac{1}{\sqrt{5}}, \frac{2}{\sqrt{5}} \right) .
\end{equation*}

\subsection{Pemetaan linear dan matriks}

Suatu \emph{\myindex{fungsi bernilai vektor}}
$F$ adalah aturan yang menerima vektor $\vec{x}$ dan menghasilkan vektor lain
$\vec{y}$. Sebagai contoh, $F$ dapat berupa penskalaan yang menggandakan ukuran
vektor:
\begin{equation*}
F(\vec{x}) = 2 \vec{x} .
\end{equation*}
Jika diterapkan, misalnya, pada $(1,3)$, kita memperoleh
\begin{equation*}
F
\left( \begin{bmatrix} 1 \\ 3 \end{bmatrix} \right)
=
2
\begin{bmatrix} 1 \\ 3 \end{bmatrix}
=
\begin{bmatrix} 2 \\ 6 \end{bmatrix} .
\end{equation*}
Jika $F$ adalah pemetaan yang membawa vektor dalam
${\mathbb R}^2$ ke
${\mathbb R}^2$ (seperti di atas), kita menulis
\begin{equation*}
F \colon {\mathbb R}^2 \to {\mathbb R}^2 .
\end{equation*}
Kata \emph{fungsi} dan \emph{pemetaan} digunakan secara kurang lebih saling menggantikan,
meskipun lebih sering \emph{pemetaan} digunakan ketika membicarakan
fungsi bernilai vektor, dan kata \emph{fungsi} sering digunakan ketika
fungsi tersebut bernilai skalar.

Mahasiswa matematika pemula (dan banyak matematikawan berpengalaman)
yang melihat ungkapan seperti
\begin{equation*}
f(3x+8y)
\end{equation*}
sangat ingin menulis
\begin{equation*}
3f(x)+8f(y) .
\end{equation*}
Lagi pula, siapa yang tidak pernah ingin menulis $\sqrt{x+y} = \sqrt{x} + \sqrt{y}$ atau
sesuatu semacam itu pada suatu saat dalam kehidupan matematisnya?
Bukankah hidup akan sederhana jika kita dapat melakukannya?
Tentu saja kita tidak selalu dapat melakukannya (misalnya, tidak untuk akar
kuadrat!).
Namun, terdapat banyak fungsi lain yang memungkinkan kita melakukan persis seperti di atas.
Fungsi-fungsi semacam itu disebut \emph{linear}.

Suatu pemetaan $F \colon {\mathbb R}^n \to {\mathbb R}^m$
disebut \emph{linear}\index{pemetaan linear} jika
\begin{equation*}
F(\vec{x}+\vec{y}) = F(\vec{x})+F(\vec{y}),
\end{equation*}
untuk sebarang vektor $\vec{x}$ dan $\vec{y}$,
dan juga
\begin{equation*}
F(\alpha \vec{x}) = \alpha F(\vec{x}) ,
\end{equation*}
untuk sebarang skalar $\alpha$.
Fungsi $F$ yang kita definisikan di atas, yang menggandakan ukuran semua vektor, bersifat linear. Mari
kita periksa:
\begin{equation*}
F(\vec{x}+\vec{y})
=
2(\vec{x}+\vec{y})
=
2\vec{x}+2\vec{y}
=
F(\vec{x})+F(\vec{y}) ,
\end{equation*}
dan juga
\begin{equation*}
F(\alpha \vec{x}) = 2 \alpha \vec{x} = \alpha 2 \vec{x} = \alpha F(\vec{x}) .
\end{equation*}

Kita juga menyebut fungsi linear sebagai
\emph{transformasi linear}\index{transformasi}.
Jika ingin benar-benar bergaya dan membuat teman-teman Anda terkesan, Anda dapat menyebutnya
\emph{operator linear}\index{operator}.
Ketika suatu pemetaan bersifat linear, kita sering tidak menuliskan tanda kurung. Kita cukup menulis
\begin{equation*}
F \vec{x}
\end{equation*}
sebagai pengganti $F(\vec{x})$. Kita melakukan ini karena linearitas berarti bahwa
pemetaan $F$
berperilaku seperti mengalikan $\vec{x}$ dengan \myquote{sesuatu.}
Sesuatu itu adalah matriks.

Suatu \emph{\myindex{matriks}}
adalah susunan bilangan berukuran $m
\times n$ ($m$ baris dan $n$ kolom). Suatu matriks
$3 \times 5$ adalah
\begin{equation*}
A =
\begin{bmatrix}
a_{11} & a_{12} & a_{13} & a_{14} & a_{15} \\
a_{21} & a_{22} & a_{23} & a_{24} & a_{25} \\
a_{31} & a_{32} & a_{33} & a_{34} & a_{35}
\end{bmatrix} .
\end{equation*}
Bilangan-bilangan $a_{ij}$ disebut \emph{elemen}\index{elemen suatu matriks}
atau \emph{entri}\index{entri suatu matriks}.

Vektor kolom hanyalah matriks berukuran $m \times 1$. Serupa dengan
vektor kolom, terdapat pula
\emph{\myindex{vektor baris}}, yaitu matriks berukuran $1 \times n$.
Jika kita memiliki matriks $n \times n$, kita mengatakan bahwa matriks itu adalah
\emph{\myindex{matriks persegi}}.

Bagaimana matriks $A$ berkaitan dengan suatu pemetaan linear?
Matriks memberi tahu kita ke mana
vektor-vektor khusus tertentu dipetakan. Kita memberikan nama kepada vektor-vektor tersebut.
\emph{\myindex{Vektor basis standar}} dari ${\mathbb R}^n$ adalah
\begin{equation*}
\vec{e}_1 =
\begin{bmatrix}
1 \\ 0 \\ 0 \\ \vdots \\ 0
\end{bmatrix} ,
\qquad
\vec{e}_2 =
\begin{bmatrix}
0 \\ 1 \\ 0 \\ \vdots \\ 0
\end{bmatrix} ,
\qquad
\vec{e}_3 =
\begin{bmatrix}
0 \\ 0 \\ 1 \\ \vdots \\ 0
\end{bmatrix} ,
\qquad
\cdots ,
\qquad
\vec{e}_n =
\begin{bmatrix}
0 \\ 0 \\ 0 \\ \vdots \\ 1
\end{bmatrix} .
\end{equation*}
Dalam ${\mathbb R}^3$, vektor-vektor ini adalah
\begin{equation*}
\vec{e}_1 =
\begin{bmatrix}
1 \\ 0 \\ 0
\end{bmatrix} ,
\qquad
\vec{e}_2 =
\begin{bmatrix}
0 \\ 1 \\ 0
\end{bmatrix} ,
\qquad
\vec{e}_3 =
\begin{bmatrix}
0 \\ 0 \\ 1
\end{bmatrix} .
\end{equation*}
Anda mungkin ingat dari kalkulus beberapa variabel bahwa vektor-vektor ini
kadang-kadang disebut $\vec{\imath}$, $\vec{\jmath}$, $\vec{k}$.
Alasan vektor-vektor ini disebut suatu \emph{\myindex{basis}} adalah karena setiap
vektor dapat ditulis sebagai \emph{\myindex{kombinasi linear}} tunggal dari vektor-vektor tersebut.
Sebagai contoh, dalam ${\mathbb R}^3$, vektor $(4,5,6)$ dapat ditulis sebagai
\begin{equation*}
4 \vec{e}_1 +
5 \vec{e}_2 +
6 \vec{e}_3
=
4
\begin{bmatrix}
1 \\ 0 \\ 0
\end{bmatrix}
+
5
\begin{bmatrix}
0 \\ 1 \\ 0
\end{bmatrix}
+
6
\begin{bmatrix}
0 \\ 0 \\ 1
\end{bmatrix}
=
\begin{bmatrix}
4 \\ 5 \\ 6
\end{bmatrix} .
\end{equation*}

Jadi, bagaimana suatu matriks merepresentasikan pemetaan linear?
Kolom-kolom matriks adalah vektor-vektor yang menjadi citra $\vec{e}_1$,
$\vec{e}_2$, dan seterusnya, oleh matriks tersebut sebagai suatu pemetaan
linear.
Sebagai contoh, perhatikan
\begin{equation*}
M =
\begin{bmatrix}
1 & 2 \\ 3 & 4
\end{bmatrix} .
\end{equation*}
Sebagai pemetaan linear, $M \colon {\mathbb R}^2 \to {\mathbb R}^2$ membawa
$\vec{e}_1 = \left[ \begin{smallmatrix} 1 \\ 0 \end{smallmatrix} \right]$ ke
$\left[ \begin{smallmatrix} 1 \\ 3 \end{smallmatrix} \right]$
dan
$\vec{e}_2 = \left[ \begin{smallmatrix} 0 \\ 1 \end{smallmatrix} \right]$ ke
$\left[ \begin{smallmatrix} 2 \\ 4 \end{smallmatrix} \right]$. Dengan kata
lain,
\begin{equation*}
M \vec{e}_1 =
\begin{bmatrix}
1 & 2 \\ 3 & 4
\end{bmatrix}
\begin{bmatrix}
1 \\ 0
\end{bmatrix}
=
\begin{bmatrix}
1 \\ 3
\end{bmatrix},
\qquad
\text{dan}
\qquad
M \vec{e}_2 =
\begin{bmatrix}
1 & 2 \\ 3 & 4
\end{bmatrix}
\begin{bmatrix}
0 \\ 1
\end{bmatrix}
=
\begin{bmatrix}
2 \\ 4
\end{bmatrix}.
\end{equation*}

Secara lebih umum, jika kita memiliki matriks $n \times m$, yaitu $n$ baris
dan $m$ kolom, maka pemetaan $A \colon {\mathbb R}^m \to {\mathbb R}^n$
membawa $\vec{e}_j$ ke kolom ke-$j$ dari $A$.
Sebagai contoh,
\begin{equation*}
A =
\begin{bmatrix}
a_{11} & a_{12} & a_{13} & a_{14} & a_{15} \\
a_{21} & a_{22} & a_{23} & a_{24} & a_{25} \\
a_{31} & a_{32} & a_{33} & a_{34} & a_{35}
\end{bmatrix}
\end{equation*}
merepresentasikan pemetaan dari ${\mathbb R}^5$ ke ${\mathbb R}^3$ yang melakukan
\begin{equation*}
A \vec{e}_1 =
\begin{bmatrix}
a_{11} \\ a_{21} \\ a_{31}
\end{bmatrix} ,
\qquad
A \vec{e}_2 =
\begin{bmatrix}
a_{12} \\ a_{22} \\ a_{32}
\end{bmatrix} ,
\qquad
A \vec{e}_3 =
\begin{bmatrix}
a_{13} \\ a_{23} \\ a_{33}
\end{bmatrix} ,
\qquad
A \vec{e}_4 =
\begin{bmatrix}
a_{14} \\ a_{24} \\ a_{34}
\end{bmatrix} ,
\qquad
A \vec{e}_5 =
\begin{bmatrix}
a_{15} \\ a_{25} \\ a_{35}
\end{bmatrix} .
\end{equation*}

Bagaimana dengan vektor lain $\vec{x}$ yang bukan vektor basis standar?
Ke mana vektor itu dipetakan? Kita menggunakan
linearitas. Pertama-tama, kita menuliskan vektor sebagai kombinasi linear dari vektor-vektor
basis standar:
\begin{equation*}
\vec{x} =
\begin{bmatrix}
x_1 \\ x_2 \\ x_3 \\ x_4 \\ x_5
\end{bmatrix}
=
x_1
\begin{bmatrix}
1 \\ 0 \\ 0 \\ 0 \\ 0
\end{bmatrix}
+
x_2
\begin{bmatrix}
0 \\ 1 \\ 0 \\ 0 \\ 0
\end{bmatrix}
+
x_3
\begin{bmatrix}
0 \\ 0 \\ 1 \\ 0 \\ 0
\end{bmatrix}
+
x_4
\begin{bmatrix}
0 \\ 0 \\ 0 \\ 1 \\ 0
\end{bmatrix}
+
x_5
\begin{bmatrix}
0 \\ 0 \\ 0 \\ 0 \\ 1
\end{bmatrix}
=
x_1 \vec{e}_1 +
x_2 \vec{e}_2 +
x_3 \vec{e}_3 +
x_4 \vec{e}_4 +
x_5 \vec{e}_5 .
\end{equation*}
Kemudian,
\begin{equation*}
A \vec{x}
=
A (
x_1 \vec{e}_1 +
x_2 \vec{e}_2 +
x_3 \vec{e}_3 +
x_4 \vec{e}_4 +
x_5 \vec{e}_5
)
=
x_1 A\vec{e}_1 +
x_2 A\vec{e}_2 +
x_3 A\vec{e}_3 +
x_4 A\vec{e}_4 +
x_5 A\vec{e}_5 .
\end{equation*}
Jika kita mengetahui ke mana $A$ membawa semua vektor basis, kita mengetahui ke mana ia membawa
semua vektor.

Andaikan $M$ adalah matriks $2 \times 2$ dari atas,
maka
\begin{equation*}
M
\begin{bmatrix}
-2 \\ 0.1
\end{bmatrix}
=
\begin{bmatrix}
1 & 2 \\
3 & 4
\end{bmatrix}
\begin{bmatrix}
-2 \\ 0.1
\end{bmatrix}
=
-2
\begin{bmatrix}
1 \\ 3
\end{bmatrix}
+
0.1
\begin{bmatrix}
2 \\ 4
\end{bmatrix}
=
\begin{bmatrix}
-1.8 \\ -5.6
\end{bmatrix} .
\end{equation*}

Setiap pemetaan linear dari ${\mathbb R}^m$ ke ${\mathbb R}^n$ dapat
direpresentasikan oleh matriks $n \times m$. Anda hanya perlu menentukan ke mana pemetaan itu
membawa vektor-vektor basis standar. Sebaliknya, setiap matriks $n \times m$
merepresentasikan pemetaan linear. Oleh karena itu, kita dapat memandang matriks sebagai
pemetaan linear, dan pemetaan linear sebagai matriks.

Atau benarkah demikian? Dalam buku ini kita terutama mempelajari operator diferensial
linear, dan operator diferensial linear adalah pemetaan linear, walaupun tidak bekerja pada
${\mathbb R}^n$, tetapi pada
ruang fungsi berdimensi tak hingga:
\begin{equation*}
L f = g .
\end{equation*}
Untuk suatu fungsi $f$, kita memperoleh fungsi $g$, dan $L$ linear dalam arti bahwa
\begin{equation*}
L ( f + h) = Lf + Lh, \qquad \text{dan} \qquad
L (\alpha f) = \alpha Lf .
\end{equation*}
untuk sebarang bilangan (skalar) $\alpha$ dan semua fungsi $f$ serta $h$.

Jadi, jawabannya sebenarnya tidak. Namun, jika kita memperhatikan vektor dalam
ruang berdimensi hingga ${\mathbb R}^n$, maka ya, setiap pemetaan linear adalah suatu
matriks.
Pada awal bagian ini kita menyebutkan bahwa kita dapat
\myquote{membuat segala sesuatu menjadi vektor.} Hal itu tidak sepenuhnya benar, tetapi
benar
secara hampiran. Ruang fungsi \myquote{berdimensi tak hingga} itu dapat
dihampiri oleh ruang berdimensi hingga, dan kemudian operator linear
hanyalah matriks. Jadi, secara hampiran, hal ini benar. Dan sejauh menyangkut perhitungan nyata
yang dapat kita lakukan pada komputer, kita memang hanya dapat bekerja dengan
berhingga banyak dimensi. Jika Anda meminta komputer atau kalkulator Anda
menggambar grafik suatu fungsi,
alat itu mengambil sampel fungsi pada berhingga banyak titik lalu
menghubungkan titik-titik tersebut\footnote{Jika Anda pernah menggunakan Matlab, Anda mungkin
memperhatikan bahwa untuk menggambar grafik fungsi, kita mengambil suatu vektor masukan, meminta Matlab
menghitung vektor nilai fungsi yang bersesuaian, lalu meminta
Matlab menggambar hasilnya.}.
Alat itu sebenarnya tidak memberikan tak hingga banyak nilai.
Cara Anda menggunakan komputer atau kalkulator sejauh ini
sudah merupakan suatu penghampiran tertentu dari ruang fungsi oleh suatu
ruang berdimensi hingga.

\medskip

Untuk mengakhiri bagian ini, kita perhatikan bagaimana $A \vec{x}$ dapat dituliskan dengan lebih ringkas.
Andaikan
\begin{equation*}
A =
\begin{bmatrix}
a_{11} & a_{12} & a_{13} \\
a_{21} & a_{22} & a_{23}
\end{bmatrix}
\qquad \text{dan} \qquad
\vec{x} =
\begin{bmatrix}
x_1 \\ x_2 \\ x_3
\end{bmatrix} .
\end{equation*}
Maka
\begin{equation*}
A \vec{x} =
\begin{bmatrix}
a_{11} & a_{12} & a_{13} \\
a_{21} & a_{22} & a_{23}
\end{bmatrix}
\begin{bmatrix}
x_1 \\ x_2 \\ x_3
\end{bmatrix}
=
\begin{bmatrix}
a_{11} x_1 + a_{12} x_2 + a_{13} x_3 \\
a_{21} x_1 + a_{22} x_2 + a_{23} x_3
\end{bmatrix}  .
\end{equation*}
Sebagai contoh,
\begin{equation*}
\begin{bmatrix}
1 & 2 \\ 3 & 4
\end{bmatrix}
\begin{bmatrix}
2 \\ -1
\end{bmatrix}
=
\begin{bmatrix}
1 \cdot 2 + 2 \cdot (-1) \\
3 \cdot 2 + 4 \cdot (-1)
\end{bmatrix}
=
\begin{bmatrix}
0 \\ 2
\end{bmatrix}  .
\end{equation*}
Artinya, Anda mengambil entri-entri dalam suatu baris matriks, mengalikannya dengan
entri-entri dalam vektor, menjumlahkan semuanya, dan hasilnya adalah entri yang bersesuaian
dalam vektor hasil.


\subsection{Latihan}

\begin{exercise}
Pada selembar kertas grafik, gambarlah vektor-vektor berikut:
\begin{tasks}(3)
\task
$\begin{bmatrix}
2 \\
5
\end{bmatrix}
$
\task
$\begin{bmatrix}
-2 \\
-4
\end{bmatrix}
$
\task
$(3,-4)$
\end{tasks}
\end{exercise}

\begin{exercise}
Pada selembar kertas grafik, gambarlah vektor $(1,2)$ yang berawal di (berpangkal di)
titik yang diberikan:
\begin{tasks}(3)
\task
berpangkal di $(0,0)$
\task
berpangkal di $(1,2)$
\task
berpangkal di $(0,-1)$
\end{tasks}
\end{exercise}

\begin{exercise}
Pada selembar kertas grafik, gambarlah operasi-operasi berikut.
Gambar dan beri label pada vektor-vektor yang terlibat dalam operasi
serta hasilnya:
\begin{tasks}(3)
\task
$\begin{bmatrix}
1 \\
-4
\end{bmatrix}
+
\begin{bmatrix}
2 \\
3
\end{bmatrix}$
\task
$\begin{bmatrix}
-3 \\
2
\end{bmatrix}
-
\begin{bmatrix}
1 \\
3
\end{bmatrix}$
\task
$3\begin{bmatrix}
2 \\
1
\end{bmatrix}$
\end{tasks}
\end{exercise}

\begin{exercise}
Hitunglah besar dari
\begin{tasks}(3)
\task
$\begin{bmatrix}
7 \\
2
\end{bmatrix}
$
\task
$\begin{bmatrix}
-2 \\
3 \\
1
\end{bmatrix}
$
\task
$(1,3,-4)$
\end{tasks}
\end{exercise}

\begin{exercise}
\pagebreak[2]
Hitunglah
\begin{tasks}(3)
\task
$\begin{bmatrix}
2 \\
3
\end{bmatrix}
+
\begin{bmatrix}
7 \\
-8
\end{bmatrix}
$
\task
$\begin{bmatrix}
-2 \\
3
\end{bmatrix}
-
\begin{bmatrix}
6 \\
-4
\end{bmatrix}
$
\task
$
-\begin{bmatrix}
-3 \\
2
\end{bmatrix}
$
\task
$
4\begin{bmatrix}
-1 \\
5
\end{bmatrix}
$
\task
$
5\begin{bmatrix}
1 \\
0
\end{bmatrix}
+
9
\begin{bmatrix}
0 \\
1
\end{bmatrix}
$
\task
$
3\begin{bmatrix}
1 \\
-8
\end{bmatrix}
-
2
\begin{bmatrix}
3 \\
-1
\end{bmatrix}
$
\end{tasks}
\end{exercise}

\begin{exercise}
Carilah vektor satuan yang searah dengan vektor yang diberikan
\begin{tasks}(3)
\task
$\begin{bmatrix}
1 \\
-3
\end{bmatrix}
$
\task
$\begin{bmatrix}
2 \\
1 \\
-1
\end{bmatrix}
$
\task
$(3,1,-2)$
\end{tasks}
\end{exercise}

\begin{exercise}
Jika $\vec{x} = (1,2)$ dan $\vec{y}$ dijumlahkan, kita memperoleh
$\vec{x}+\vec{y} = (0,2)$. Berapakah $\vec{y}$?
\end{exercise}

\begin{exercise}
Tuliskan $(1,2,3)$ sebagai kombinasi linear dari vektor-vektor basis standar
$\vec{e}_1$, $\vec{e}_2$, dan $\vec{e}_3$.
\end{exercise}

\begin{exercise}
Jika besar $\vec{x}$ adalah 4, berapakah besar dari
\begin{tasks}(6)
\task
$0\vec{x}$
\task
$3\vec{x}$
\task
$-\vec{x}$
\task
$-4\vec{x}$
\task
$\vec{x}+\vec{x}$
\task
$\vec{x}-\vec{x}$
\end{tasks}
\end{exercise}

\begin{exercise}
Misalkan suatu pemetaan linear $F \colon {\mathbb R}^2 \to {\mathbb R}^2$
membawa $(1,0)$ ke $(2,-1)$ dan membawa $(0,1)$ ke $(3,3)$.
Ke manakah pemetaan itu membawa
\begin{tasks}(3)
\task
$(1,1)$
\task
$(2,0)$
\task
$(2,-1)$
\end{tasks}
\end{exercise}

\begin{exercise}
Misalkan suatu pemetaan linear $F \colon {\mathbb R}^3 \to {\mathbb R}^2$
membawa $(1,0,0)$ ke $(2,1)$, membawa $(0,1,0)$ ke $(3,4)$, dan
membawa $(0,0,1)$ ke $(5,6)$. Tuliskan matriks yang merepresentasikan
pemetaan $F$.
\end{exercise}

\begin{exercise}
Misalkan suatu pemetaan $F \colon {\mathbb R}^2 \to \mathbb{R}^2$ membawa
$(1,0)$ ke $(1,2)$, $(0,1)$ ke $(3,4)$, dan $(1,1)$ ke $(0,-1)$.
Jelaskan mengapa $F$ tidak linear.
\end{exercise}

\begin{exercise}[menantang]
Misalkan ${\mathbb R}^3$ merepresentasikan ruang polinomial kuadrat
dalam $t$: Suatu titik $(a_0,a_1,a_2)$ dalam ${\mathbb R}^3$ merepresentasikan
polinomial $a_0 + a_1 t + a_2 t^2$.
Pandang turunan $\frac{d}{dt}$ sebagai pemetaan dari ${\mathbb R}^3$ ke
${\mathbb R}^3$,
dan perhatikan bahwa $\frac{d}{dt}$ linear.
Tuliskan $\frac{d}{dt}$ sebagai matriks $3 \times 3$.
\end{exercise}

\setcounter{exercise}{100}

\begin{exercise}
Hitunglah besar dari
\begin{tasks}(3)
\task
$\begin{bmatrix}
1 \\
3
\end{bmatrix}
$
\task
$\begin{bmatrix}
2 \\
3 \\
-1
\end{bmatrix}
$
\task
$(-2,1,-2)$
\end{tasks}
\end{exercise}
\exsol{%
a)~$\sqrt{10}$
\quad b)~$\sqrt{14}$
\quad c)~$3$
}

\begin{exercise}
Carilah vektor satuan yang searah dengan vektor yang diberikan
\begin{tasks}(3)
\task
$\begin{bmatrix}
-1 \\
1
\end{bmatrix}
$
\task
$\begin{bmatrix}
1 \\
-1 \\
2
\end{bmatrix}
$
\task
$(2,-5,2)$
\end{tasks}
\end{exercise}
\exsol{%
a)~
$\begin{bmatrix}
\frac{-1}{\sqrt{2}} \\
\frac{1}{\sqrt{2}}
\end{bmatrix}$
\quad
b)~$\begin{bmatrix}
\frac{1}{\sqrt{6}} \\
\frac{-1}{\sqrt{6}} \\
\frac{2}{\sqrt{6}}
\end{bmatrix}$
\quad
c)~$\left( \frac{2}{\sqrt{33}},\frac{-5}{\sqrt{33}},\frac{2}{\sqrt{33}} \right)$
}

\begin{exercise}
\pagebreak[2]
Hitunglah
\begin{tasks}(3)
\task
$\begin{bmatrix}
3 \\
1
\end{bmatrix}
+
\begin{bmatrix}
6 \\
-3
\end{bmatrix}
$
\task
$\begin{bmatrix}
-1 \\
2
\end{bmatrix}
-
\begin{bmatrix}
2 \\
-1
\end{bmatrix}
$
\task
$
-\begin{bmatrix}
-5 \\
3
\end{bmatrix}
$
\task
$
2\begin{bmatrix}
-2 \\
4
\end{bmatrix}
$
\task
$
3\begin{bmatrix}
1 \\
0
\end{bmatrix}
+
7
\begin{bmatrix}
0 \\
1
\end{bmatrix}
$
\task
$
2\begin{bmatrix}
2 \\
-3
\end{bmatrix}
-
6
\begin{bmatrix}
2 \\
-1
\end{bmatrix}
$
\end{tasks}
\end{exercise}
\exsol{%
a)~$\begin{bmatrix}
9 \\
-2
\end{bmatrix}$
\quad b)~$\begin{bmatrix}
-3 \\
3
\end{bmatrix}$
\quad c)~$\begin{bmatrix}
5 \\
-3
\end{bmatrix}$
\quad d)~$\begin{bmatrix}
-4 \\
8
\end{bmatrix}$
\quad e)~$\begin{bmatrix}
3 \\
7
\end{bmatrix}$
\quad f)~$\begin{bmatrix}
-8 \\
0
\end{bmatrix}$
}

\begin{exercise}
Jika besar $\vec{x}$ adalah 5, berapakah besar dari
\begin{tasks}(3)
\task
$4\vec{x}$
\task
$-2\vec{x}$
\task
$-4\vec{x}$
\end{tasks}
\end{exercise}
\exsol{%
a)~$20$
\quad b)~$10$
\quad c)~$20$
}

\begin{exercise}
Misalkan suatu pemetaan linear $F \colon {\mathbb R}^2 \to {\mathbb R}^2$
membawa $(1,0)$ ke $(1,-1)$ dan membawa $(0,1)$ ke $(2,0)$.
Ke manakah pemetaan itu membawa
\begin{tasks}(3)
\task
$(1,1)$
\task
$(0,2)$
\task
$(1,-1)$
\end{tasks}
\end{exercise}
\exsol{%
a)~$(3,-1)$ \quad b)~$(4,0)$ \quad c)~$(-1,-1)$
}

%%%%%%%%%%%%%%%%%%%%%%%%%%%%%%%%%%%%%%%%%%%%%%%%%%%%%%%%%%%%%%%%%%%%%%%%%%%%%%

\sectionnewpage
\section{Aljabar matriks}
\label{matalg:section}

%mbxINTROSUBSECTION

\sectionnotes{2--3 kuliah}

\subsection{Matriks satu-kali-satu}

Mari kita berikan motivasi bagi apa yang ingin kita capai dengan matriks.
Pemetaan linear bernilai real pada garis real, yaitu fungsi linear
yang menerima bilangan dan menghasilkan bilangan, hanyalah perkalian dengan suatu
bilangan. Perhatikan suatu pemetaan yang didefinisikan dengan mengalikan suatu
bilangan. Mari kita sebut bilangan ini $\alpha$. Pemetaan tersebut membawa $x$ ke
$\alpha x$. Kita dapat
\emph{menjumlahkan} pemetaan-pemetaan semacam itu:
Jika kita memiliki pemetaan lain $\beta$, maka
\begin{equation*}
\alpha x + \beta x = (\alpha + \beta) x .
\end{equation*}
Kita memperoleh pemetaan baru $\alpha+\beta$ yang mengalikan $x$ dengan, tentu saja,
$\alpha+\beta$. Jika $D$ adalah pemetaan yang menggandakan masukannya,
$Dx = 2x$, dan $T$ adalah pemetaan yang melipat-tigakan, $Tx = 3x$, maka
$D+T$ adalah pemetaan yang mengalikan dengan $5$, $(D+T)x = 5x$.

Secara serupa, kita dapat \emph{mengomposisikan} pemetaan semacam itu.
Artinya, kita dapat menerapkan yang satu lalu yang lain.
Kita mengambil $x$, memasukkannya ke
pemetaan pertama $\alpha$ untuk memperoleh $\alpha$ kali $x$, lalu memasukkan
$\alpha x$ ke pemetaan kedua $\beta$. Dengan kata lain,
\begin{equation*}
\beta ( \alpha x ) = (\beta \alpha) x .
\end{equation*}
Kita cukup mengalikan kedua bilangan tersebut. Dengan menggunakan pemetaan
penggandaan dan pelipat-tigaan kita, jika kita menggandakan lalu melipat-tigakan,
yakni $T(Dx)$, maka
kita memperoleh $3(2x) = 6x$. Komposisi $TD$ adalah pemetaan yang mengalikan
dengan $6$. Untuk matriks yang lebih besar, komposisi juga menjadi sejenis
perkalian.

\subsection{Penjumlahan matriks dan perkalian skalar}

Pemetaan yang mengalikan bilangan dengan bilangan hanyalah matriks $1 \times 1$.
Bilangan $\alpha$ di atas dapat ditulis sebagai matriks $[\alpha]$.
Barangkali kita ingin melakukan hal yang sama pada semua matriks seperti yang
kita lakukan pada matriks $1 \times 1$ pada awal bagian ini.
Pertama, mari kita jumlahkan matriks.
Jika kita memiliki matriks $A$ dan matriks $B$ yang berukuran sama,
misalnya $m \times n$, maka keduanya merupakan pemetaan dari
${\mathbb{R}}^n$ ke ${\mathbb{R}}^m$. Pemetaan $A+B$ juga harus merupakan pemetaan dari
${\mathbb{R}}^n$ ke ${\mathbb{R}}^m$, dan harus melakukan hal berikut terhadap
vektor:
\begin{equation*}
(A+B) \vec{x} = A\vec{x} + B \vec{x} .
\end{equation*}
Ternyata kita cukup menjumlahkan matriks-matriks itu elemen demi elemen: Jika
entri ke-$ij$ dari $A$ adalah $a_{ij}$, dan
entri ke-$ij$ dari $B$ adalah $b_{ij}$, maka
entri ke-$ij$ dari $A+B$ adalah $a_{ij} + b_{ij}$. Jika
\begin{equation*}
A =
\begin{bmatrix}
a_{11} & a_{12} & a_{13}  \\
a_{21} & a_{22} & a_{23}
\end{bmatrix}
\qquad \text{dan} \qquad
B =
\begin{bmatrix}
b_{11} & b_{12} & b_{13}  \\
b_{21} & b_{22} & b_{23}
\end{bmatrix} ,
\end{equation*}
maka
\begin{equation*}
A+B =
\begin{bmatrix}
a_{11} + b_{11} & a_{12} + b_{12} & a_{13} + b_{13}  \\
a_{21} + b_{21} & a_{22} + b_{22} & a_{23} + b_{23}
\end{bmatrix} .
\end{equation*}
Kita ilustrasikan dengan contoh yang lebih konkret:
\begin{equation*}
\begin{bmatrix}
1 & 2 \\
3 & 4 \\
5 & 6
\end{bmatrix}
+
\begin{bmatrix}
7 & 8 \\
9 & 10 \\
11 & -1
\end{bmatrix}
=
\begin{bmatrix}
1+7 & 2+8 \\
3+9 & 4+10 \\
5+11 & 6-1
\end{bmatrix}
=
\begin{bmatrix}
8 & 10 \\
12 & 14 \\
16 & 5
\end{bmatrix} .
\end{equation*}
Mari kita periksa bahwa ini memberikan hasil yang tepat pada suatu vektor. Kita menggunakan sebagian
aljabar vektor yang telah kita ketahui, dan
mengelompokkan ulang suku-sukunya:
\begin{equation*}
\begin{split}
\begin{bmatrix}
1 & 2 \\
3 & 4 \\
5 & 6
\end{bmatrix}
\begin{bmatrix}
2 \\ -1
\end{bmatrix}
+
\begin{bmatrix}
7 & 8 \\
9 & 10 \\
11 & -1
\end{bmatrix}
\begin{bmatrix}
2 \\ -1
\end{bmatrix}
& =
\left(
2
\begin{bmatrix}
1 \\
3 \\
5
\end{bmatrix}
-
\begin{bmatrix}
2 \\
4 \\
6
\end{bmatrix}
\right)
+
\left(
2
\begin{bmatrix}
7 \\
9 \\
11
\end{bmatrix}
-
\begin{bmatrix}
8 \\
10 \\
-1
\end{bmatrix}
\right)
\\
& =
2
\left(
\begin{bmatrix}
1 \\
3 \\
5
\end{bmatrix}
+
\begin{bmatrix}
7 \\
9 \\
11
\end{bmatrix}
\right)
-
\left(
\begin{bmatrix}
2 \\
4 \\
6
\end{bmatrix}
+
\begin{bmatrix}
8 \\
10 \\
-1
\end{bmatrix}
\right)
\\
& =
2
\begin{bmatrix}
1+7 \\
3+9 \\
5+11
\end{bmatrix}
-
\begin{bmatrix}
2+8 \\
4+10 \\
6-1
\end{bmatrix}
=
2
\begin{bmatrix}
8 \\
12 \\
16
\end{bmatrix}
-
\begin{bmatrix}
10 \\
14 \\
5
\end{bmatrix}
\\
& =
\begin{bmatrix}
8 & 10 \\
12 & 14 \\
16 & 5
\end{bmatrix}
\begin{bmatrix}
2 \\
-1
\end{bmatrix}
\quad
\left(
=
\begin{bmatrix}
2(8)- 10 \\
2(12) - 14 \\
2(16) - 5
\end{bmatrix}
=
\begin{bmatrix}
6 \\
10 \\
27
\end{bmatrix}
\right) .
\end{split}
\end{equation*}
Jika bilangan-bilangan tersebut kita ganti dengan huruf, hal itu akan menjadi sebuah bukti!
Perhatikan bahwa sebenarnya kita tidak harus menghitung
hasilnya untuk meyakinkan diri bahwa kedua ungkapan tersebut sama.

Jika ukuran matriks tidak cocok, penjumlahan tidak terdefinisi.
Jika $A$ berukuran $3 \times 2$ dan $B$ berukuran $2 \times 5$, kita tidak dapat menjumlahkan
kedua matriks. Kita tidak tahu apa yang mungkin dimaksud oleh penjumlahan itu.

\medskip

Berguna pula untuk memiliki matriks yang, ketika dijumlahkan dengan sebarang matriks lain,
tidak mengubah apa pun. Inilah matriks nol, matriks yang semua entrinya nol:
\begin{equation*}
\begin{bmatrix}
1 & 2 \\
3 & 4
\end{bmatrix}
+
\begin{bmatrix}
0 & 0 \\
0 & 0
\end{bmatrix}
=
\begin{bmatrix}
1 & 2 \\
3 & 4
\end{bmatrix} .
\end{equation*}
Kita sering menotasikan matriks nol
dengan $0$ tanpa menyebutkan ukurannya. Kita kemudian cukup menulis $A + 0$, dengan
mengasumsikan bahwa $0$ adalah matriks nol yang berukuran sama dengan $A$.

\medskip

Sebenarnya terdapat dua hal yang dapat kita gunakan untuk mengalikan matriks. Kita dapat mengalikan
matriks dengan skalar atau dengan matriks lain. Mula-mula mari kita
perhatikan perkalian dengan skalar.
Untuk suatu matriks $A$ dan skalar $\alpha$, kita ingin $\alpha A$ menjadi matriks
yang memenuhi
\begin{equation*}
(\alpha A) \vec{x} = \alpha (A \vec{x}) .
\end{equation*}
Artinya, hasilnya cukup diskalakan dengan $\alpha$. Jika Anda memikirkannya,
menskalakan setiap suku dalam $A$ dengan $\alpha$ menghasilkan tepat hal itu:
Jika
\begin{equation*}
A =
\begin{bmatrix}
a_{11} & a_{12} & a_{13}  \\
a_{21} & a_{22} & a_{23}
\end{bmatrix},
\qquad\text{maka} \qquad
\alpha A =
\begin{bmatrix}
\alpha a_{11} & \alpha a_{12} & \alpha a_{13}  \\
\alpha a_{21} & \alpha a_{22} & \alpha a_{23}
\end{bmatrix} .
\end{equation*}
Sebagai contoh,
\begin{equation*}
2
\begin{bmatrix}
1 & 2 & 3 \\
4 & 5 & 6
\end{bmatrix} =
\begin{bmatrix}
2 & 4 & 6 \\
8 & 10 & 12
\end{bmatrix} .
\end{equation*}

Mari kita daftarkan beberapa sifat penjumlahan matriks dan perkalian skalar.
Nyatakan matriks nol dengan $0$,
skalar dengan $\alpha$, $\beta$, dan matriks dengan $A$, $B$, $C$. Maka:
\begin{align*}
A + 0 & = A = 0 + A , \\
A + B & = B + A , \\
(A + B) + C & = A + (B + C) , \\
\alpha(A+B) & = \alpha A+\alpha B, \\
(\alpha+\beta)A & = \alpha A + \beta A.
\end{align*}
Aturan-aturan ini semestinya tampak sangat familier.

\subsection{Perkalian matriks}

Seperti disebutkan di atas, komposisi pemetaan linear juga merupakan suatu
perkalian matriks. Misalkan $A$ adalah matriks $m \times n$,
yakni $A$ membawa
${\mathbb R}^n$ ke
${\mathbb R}^m$,
dan $B$ adalah matriks $n \times p$, yakni $B$ membawa
${\mathbb R}^p$ ke
${\mathbb R}^n$. Komposisi $AB$ harus bekerja sebagai berikut
\begin{equation*}
AB\vec{x} = A(B\vec{x}) .
\end{equation*}
Pertama, suatu vektor $\vec{x}$ dalam ${\mathbb R}^p$ dibawa ke
vektor $B\vec{x}$ dalam
${\mathbb R}^n$. Kemudian pemetaan $A$ membawanya ke vektor $A(B\vec{x})$
dalam ${\mathbb R}^m$. Dengan kata lain, komposisi $AB$ harus berupa matriks $m
\times p$. Ditinjau dari ukurannya, kita harus memperoleh
\begin{equation*}
%mbxSTARTIGNORE
\text{``}
%mbxENDIGNORE
%mbxlatex \text{"}
\quad
[ m \times n ]
\,
[ n \times p ]
=
[ m \times p ] . \quad
%mbxSTARTIGNORE
\text{''}
%mbxENDIGNORE
%mbxlatex \text{"}
\end{equation*}
Perhatikan bahwa ukuran di tengah harus cocok.

Baiklah\@, sekarang kita mengetahui ukuran matriks yang seharusnya dapat kita kalikan
dan ukuran hasil kalinya.
Mari kita lihat cara benar-benar menghitung perkalian matriks.
Kita mulai dengan apa yang disebut
\emph{\myindex{hasil kali titik}} (atau \emph{\myindex{hasil kali dalam}}) dari dua vektor.
Biasanya, ini berupa vektor baris yang dikalikan
dengan vektor kolom berukuran sama. Hasil kali titik mengalikan
setiap pasangan entri dari vektor pertama dan kedua, lalu menjumlahkan
hasil kali tersebut. Hasil akhirnya adalah satu bilangan.
Sebagai contoh,
\begin{equation*}
\begin{bmatrix}
a_1 & a_2 & a_3
\end{bmatrix}
\cdot
\begin{bmatrix}
b_1 \\
b_2 \\
b_3
\end{bmatrix}
= a_1 b_1 + a_2 b_2 + a_3 b_3 .
\end{equation*}
Demikian pula untuk vektor yang lebih besar (atau lebih kecil).
Hasil kali titik sebenarnya merupakan hasil kali dua matriks: matriks $1 \times n$
dan matriks $n \times 1$ yang menghasilkan
matriks $1 \times 1$---suatu bilangan.

Dengan berbekal hasil kali titik, kita mendefinisikan
\emph{\myindex{hasil kali matriks}}\index{hasil kali matriks}.
Kita menyatakan baris ke-$i$ dari $A$ dengan $\operatorname{row}_i(A)$
dan kolom ke-$j$ dari $A$ dengan
$\operatorname{column}_j(A)$.
Untuk matriks $m \times n$ $A$ dan matriks $n \times p$ $B$,
kita dapat menghitung hasil kali $AB$: Matriks $AB$ adalah matriks $m \times p$
yang entri ke-$ij$-nya merupakan hasil kali titik
\begin{equation*}
\operatorname{row}_i(A) \cdot
\operatorname{column}_j(B) .
\end{equation*}
Sebagai contoh, jika diberikan matriks $2 \times 3$ dan matriks $3 \times 2$,
kita harus memperoleh matriks $2 \times 2$:
\begin{equation} \label{linalg:eqmatrixmulex}
\begin{bmatrix}
a_{11} & a_{12} & a_{13} \\
a_{21} & a_{22} & a_{23}
\end{bmatrix}
\begin{bmatrix}
b_{11} & b_{12} \\
b_{21} & b_{22} \\
b_{31} & b_{32}
\end{bmatrix}
=
\begin{bmatrix}
a_{11} b_{11} +
a_{12} b_{21} +
a_{13} b_{31} & ~~
a_{11} b_{12} +
a_{12} b_{22} +
a_{13} b_{32} \\
a_{21} b_{11} +
a_{22} b_{21} +
a_{23} b_{31} & ~~
a_{21} b_{12} +
a_{22} b_{22} +
a_{23} b_{32}
\end{bmatrix} ,
\end{equation}
atau dengan sejumlah bilangan:
\begin{equation*}
\begin{bmatrix}
1 & 2 & 3 \\
4 & 5 & 6
\end{bmatrix}
\begin{bmatrix}
-1 & 2 \\
-7 & 0 \\
1 & -1
\end{bmatrix}
=
\begin{bmatrix}
1\cdot (-1) + 2\cdot (-7) + 3 \cdot 1 & ~~
1\cdot 2 + 2\cdot 0 + 3 \cdot (-1) \\
4\cdot (-1) + 5\cdot (-7) + 6 \cdot 1 & ~~
4\cdot 2 + 5\cdot 0 + 6 \cdot (-1)
\end{bmatrix}
=
\begin{bmatrix}
-12 & -1 \\
-33 & 2
\end{bmatrix} .
\end{equation*}

Akibat yang berguna dari definisi tersebut ialah bahwa evaluasi
$A \vec{x}$ untuk suatu matriks $A$
dan suatu vektor (kolom) $\vec{x}$ juga merupakan perkalian matriks.
Itulah sebenarnya alasan kita memandang
vektor sebagai vektor kolom, atau matriks $n \times 1$.
Sebagai contoh,
\begin{equation*}
\begin{bmatrix}
1 & 2 \\ 3 & 4
\end{bmatrix}
\begin{bmatrix}
2 \\ -1
\end{bmatrix}
=
\begin{bmatrix}
1 \cdot 2 + 2 \cdot (-1) \\
3 \cdot 2 + 4 \cdot (-1)
\end{bmatrix}
=
\begin{bmatrix}
0 \\ 2
\end{bmatrix}  .
\end{equation*}
Jika Anda melihat bagian sebelumnya, itulah persis contoh terakhir
yang kita berikan.

Luangkan sejenak untuk mencermati perhitungan perkalian matriks $AB$
dan
definisi sebelumnya tentang $A\vec{y}$ sebagai suatu pemetaan.
Yang kita lakukan dalam perkalian matriks adalah menerapkan
pemetaan $A$ pada kolom-kolom $B$. Hal ini biasanya dituliskan sebagai berikut.
Andaikan kita menulis matriks $n \times p$
$B = [ \vec{b}_1 ~ \vec{b}_2 ~ \cdots ~ \vec{b}_p ]$, dengan
$\vec{b}_1, \vec{b}_2, \ldots, \vec{b}_p$ adalah kolom-kolom $B$. Maka
untuk matriks $m \times n$ $A$,
\begin{equation*}
AB =
A [ \vec{b}_1 ~ \vec{b}_2 ~ \cdots ~ \vec{b}_p ]
=
[ A\vec{b}_1 ~ A\vec{b}_2 ~ \cdots ~ A\vec{b}_p ] .
\end{equation*}
Kolom-kolom matriks $m \times p$ $AB$
adalah
vektor-vektor $A\vec{b}_1, A\vec{b}_2, \ldots, A\vec{b}_p$.
Sebagai contoh, dalam \eqref{linalg:eqmatrixmulex},
kolom-kolom dari
\begin{equation*}
\begin{bmatrix}
a_{11} & a_{12} & a_{13} \\
a_{21} & a_{22} & a_{23}
\end{bmatrix}
\begin{bmatrix}
b_{11} & b_{12} \\
b_{21} & b_{22} \\
b_{31} & b_{32}
\end{bmatrix}
\end{equation*}
adalah
\begin{equation*}
\begin{bmatrix}
a_{11} & a_{12} & a_{13} \\
a_{21} & a_{22} & a_{23}
\end{bmatrix}
\begin{bmatrix}
b_{11} \\
b_{21} \\
b_{31}
\end{bmatrix}
\qquad
\text{dan}
\qquad
\begin{bmatrix}
a_{11} & a_{12} & a_{13} \\
a_{21} & a_{22} & a_{23}
\end{bmatrix}
\begin{bmatrix}
b_{12} \\
b_{22} \\
b_{32}
\end{bmatrix} .
\end{equation*}
Ini merupakan cara yang sangat berguna untuk memahami perkalian matriks.
Cara ini juga semestinya mempermudah mengingat cara melakukan perkalian matriks.

\subsection{Beberapa aturan aljabar matriks}

Untuk perkalian, kita menginginkan analog dari 1. Artinya,
kita menginginkan matriks yang membiarkan segala sesuatu tetap seperti semula.
Analog ini adalah apa yang disebut
\emph{\myindex{matriks identitas}}.
Matriks identitas adalah matriks persegi dengan 1 pada
diagonal utama dan nol di semua tempat lain. Matriks ini biasanya dinotasikan dengan $I$.
Untuk setiap ukuran, kita memiliki matriks identitas yang berbeda, sehingga kadang-kadang kita menyatakan
ukurannya sebagai subskrip. Sebagai contoh, $I_3$ adalah matriks identitas
$3 \times 3$
\begin{equation*}
I = I_3 =
\begin{bmatrix}
1 & 0 & 0 \\
0 & 1 & 0 \\
0 & 0 & 1
\end{bmatrix} .
\end{equation*}
Mari kita lihat cara matriks itu bekerja pada contoh yang lebih kecil,
\begin{equation*}
\begin{bmatrix}
a_{11} & a_{12} \\
a_{21} & a_{22}
\end{bmatrix}
\begin{bmatrix}
1 & 0 \\
0 & 1
\end{bmatrix} =
\begin{bmatrix}
a_{11} \cdot 1 + a_{12} \cdot 0
& ~~
a_{11} \cdot 0 + a_{12} \cdot 1
\\
a_{21} \cdot 1 + a_{22} \cdot 0
& ~~
a_{21} \cdot 0 + a_{22} \cdot 1
\end{bmatrix}
=
\begin{bmatrix}
a_{11} & a_{12} \\
a_{21} & a_{22}
\end{bmatrix} .
\end{equation*}
Perkalian dengan identitas dari sebelah kiri tampak serupa, dan juga
tidak mengubah apa pun.

\medskip

Kita memiliki aturan-aturan berikut untuk perkalian matriks. Andaikan
$A$, $B$, $C$ adalah matriks dengan ukuran yang tepat sehingga semua yang berikut
bermakna. Misalkan $\alpha$ menyatakan suatu skalar (bilangan). Maka
\begin{align*}
A(BC) & = (AB)C & & \text{(\myindex{hukum asosiatif})} , \\
A(B+C) & = AB + AC & & \text{(\myindex{hukum distributif})} , \\
(B+C)A & = BA + CA & & \text{(hukum distributif)} , \\
\alpha(AB) & = (\alpha A)B = A(\alpha B) , & &  \\
IA & = A = AI & & \text{(identitas)}.
\end{align*}

\begin{example}
Mari kita demonstrasikan beberapa aturan ini. Sebagai contoh, hukum asosiatif:
\begin{equation*}
\underbrace{
\begin{bmatrix}
-3 & 3 \\ 2 & -2
\end{bmatrix}
}_A
\biggl(
\underbrace{
\begin{bmatrix}
4 & 4 \\ 1 & -3
\end{bmatrix}
}_B
\underbrace{
\begin{bmatrix}
-1 & 4 \\ 5 & 2
\end{bmatrix}
}_C
\biggr)
=
\underbrace{
\begin{bmatrix}
-3 & 3 \\ 2 & -2
\end{bmatrix}
}_A
\underbrace{
\begin{bmatrix}
16 & 24 \\ -16 & -2
\end{bmatrix}
}_{BC}
=
\underbrace{
\begin{bmatrix}
-96 & -78 \\ 64 & 52
\end{bmatrix}
}_{A(BC)} ,
\end{equation*}
dan
\begin{equation*}
\biggl(
\underbrace{
\begin{bmatrix}
-3 & 3 \\ 2 & -2
\end{bmatrix}
}_A
\underbrace{
\begin{bmatrix}
4 & 4 \\ 1 & -3
\end{bmatrix}
}_B
\biggr)
\underbrace{
\begin{bmatrix}
-1 & 4 \\ 5 & 2
\end{bmatrix}
}_C
=
\underbrace{
\begin{bmatrix}
-9 & -21 \\ 6 & 14
\end{bmatrix}
}_{AB}
\underbrace{
\begin{bmatrix}
-1 & 4 \\ 5 & 2
\end{bmatrix}
}_C
=
\underbrace{
\begin{bmatrix}
-96 & -78 \\ 64 & 52
\end{bmatrix}
}_{(AB)C} .
\end{equation*}
Atau bagaimana dengan perkalian oleh skalar:
\begin{equation*}
10
\biggl(
\underbrace{
\begin{bmatrix}
-3 & 3 \\ 2 & -2
\end{bmatrix}
}_A
\underbrace{
\begin{bmatrix}
4 & 4 \\ 1 & -3
\end{bmatrix}
}_B
\biggr)
=
10
\underbrace{
\begin{bmatrix}
-9 & -21 \\ 6 & 14
\end{bmatrix}
}_{A B}
=
\underbrace{
\begin{bmatrix}
-90 & -210 \\ 60 & 140
\end{bmatrix}
}_{10 (AB)} ,
\end{equation*}
\begin{equation*}
\biggl(
10
\underbrace{
\begin{bmatrix}
-3 & 3 \\ 2 & -2
\end{bmatrix}
}_A
\biggr)
\underbrace{
\begin{bmatrix}
4 & 4 \\ 1 & -3
\end{bmatrix}
}_B
=
\underbrace{
\begin{bmatrix}
-30 & 30 \\ 20 & -20
\end{bmatrix}
}_{10 A}
\underbrace{
\begin{bmatrix}
4 & 4 \\ 1 & -3
\end{bmatrix}
}_B
=
\underbrace{
\begin{bmatrix}
-90 & -210 \\ 60 & 140
\end{bmatrix}
}_{(10 A)B} ,
\end{equation*}
dan
\begin{equation*}
\underbrace{
\begin{bmatrix}
-3 & 3 \\ 2 & -2
\end{bmatrix}
}_A
\biggl(
10
\underbrace{
\begin{bmatrix}
4 & 4 \\ 1 & -3
\end{bmatrix}
}_B
\biggr)
=
\underbrace{
\begin{bmatrix}
-3 & 3 \\ 2 & -2
\end{bmatrix}
}_{A}
\underbrace{
\begin{bmatrix}
40 & 40 \\ 10 & -30
\end{bmatrix}
}_{10B}
=
\underbrace{
\begin{bmatrix}
-90 & -210 \\ 60 & 140
\end{bmatrix}
}_{A(10B)} .
\end{equation*}
\end{example}

Satu aturan perkalian yang telah Anda gunakan sejak sekolah dasar untuk bilangan
sangat mencolok tidak ada untuk matriks.
Artinya, perkalian matriks
tidak komutatif. Pertama, hanya karena $AB$ bermakna, mungkin saja
$BA$ bahkan tidak terdefinisi. Sebagai contoh, jika $A$ berukuran $2 \times 3$, dan
$B$ berukuran $3 \times 4$, kita dapat mengalikan $AB$ tetapi tidak $BA$.

Bahkan jika $AB$ dan $BA$ keduanya terdefinisi, hal itu tidak berarti keduanya sama.
Sebagai contoh, ambillah
$A = \left[ \begin{smallmatrix} 1 & 1 \\ 1 & 1 \end{smallmatrix} \right]$
dan
$B = \left[ \begin{smallmatrix} 1 & 0 \\ 0 & 2 \end{smallmatrix} \right]$:
\begin{equation*}
AB =
\begin{bmatrix} 1 & 1 \\ 1 & 1 \end{bmatrix}
\begin{bmatrix} 1 & 0 \\ 0 & 2 \end{bmatrix}
=
\begin{bmatrix} 1 & 2 \\ 1 & 2 \end{bmatrix}
\qquad
\neq
\qquad
\begin{bmatrix} 1 & 1 \\ 2 & 2 \end{bmatrix}
=
\begin{bmatrix} 1 & 0 \\ 0 & 2 \end{bmatrix}
\begin{bmatrix} 1 & 1 \\ 1 & 1 \end{bmatrix}
=
BA .
\end{equation*}

\subsection{Invers}

Beberapa aturan aljabar lain yang Anda ketahui untuk bilangan tidak sepenuhnya
berlaku pada matriks:
\begin{enumerate}[(i)]
\item $AB = AC$ tidak selalu mengakibatkan $B=C$, bahkan jika $A$ bukan 0.
\item $AB = 0$ tidak selalu berarti bahwa $A=0$ atau $B=0$.
\end{enumerate}
Sebagai contoh:
\begin{equation*}
\begin{bmatrix} 0 & 1 \\ 0 & 0 \end{bmatrix}
\begin{bmatrix} 0 & 1 \\ 0 & 0 \end{bmatrix}
=
\begin{bmatrix} 0 & 0 \\ 0 & 0 \end{bmatrix}
=
\begin{bmatrix} 0 & 1 \\ 0 & 0 \end{bmatrix}
\begin{bmatrix} 0 & 2 \\ 0 & 0 \end{bmatrix} .
\end{equation*}

Agar aturan-aturan ini berlaku, kita tidak sekadar memerlukan salah satu matriks tidak nol,
kita perlu \myquote{membagi} dengan
suatu matriks. Di sinilah \emph{\myindex{invers matriks}} berperan.
Andaikan $A$ dan $B$ adalah matriks $n \times n$ sedemikian sehingga
\begin{equation*}
AB = I = BA .
\end{equation*}
Kita menyebut $B$ sebagai invers dari $A$ dan menotasikan $B$ dengan $A^{-1}$.
Barangkali tidak mengejutkan bahwa ${(A^{-1})}^{-1} = A$, sebab jika invers $A$
adalah $B$, maka invers $B$ adalah $A$.
Jika invers $A$ ada, kita mengatakan bahwa $A$
\emph{invertibel\index{matriks invertibel}}.
Jika $A$ tidak invertibel, kita mengatakan bahwa $A$
\emph{singular\index{matriks singular}}.

Jika $A = [a]$ adalah matriks $1 \times 1$, maka $A^{-1}$ adalah $a^{-1} =
\frac{1}{a}$. Dari situlah notasinya berasal. Perhitungannya jauh tidak sesederhana ini ketika $A$
lebih besar.

Perumusan aturan pencoretan yang tepat adalah:
\begin{center}
\emph{Jika $A$ invertibel,
maka
$AB = AC$ mengakibatkan $B=C$.}
\end{center}
Perhitungannya sama dengan yang Anda lakukan dalam aljabar biasa dengan bilangan,
tetapi Anda harus
berhati-hati agar tidak pernah mempertukarkan urutan matriks:
\begin{align*}
AB & = AC , \\
A^{-1}AB & = A^{-1}AC , \\
IB & = IC , \\
B & = C .
\end{align*}
Demikian pula untuk pencoretan dari sebelah kanan:
\begin{center}
\emph{Jika $A$ invertibel,
maka $BA = CA$ mengakibatkan $B=C$.}
\end{center}
Aturan tersebut antara lain menyatakan bahwa
invers suatu matriks tunggal jika ada: Jika $AB = I = AC$, maka $A$
invertibel dan $B=C$.

Nanti kita akan melihat cara menghitung invers suatu matriks
secara umum. Untuk sekarang,
perhatikan bahwa terdapat rumus sederhana untuk invers
matriks $2 \times 2$
\begin{equation*}
\begin{bmatrix}
a & b \\
c & d
\end{bmatrix}^{-1}
=
\frac{1}{ad-bc}
\begin{bmatrix}
d & -b \\
-c & a
\end{bmatrix} .
\end{equation*}

Sebagai contoh:
\begin{equation*}
\begin{bmatrix}
1 & 1 \\
2 & 4
\end{bmatrix}^{-1}
=
\frac{1}{1\cdot 4-1 \cdot 2}
\begin{bmatrix}
4 & -1 \\
-2 & 1
\end{bmatrix}
=
\begin{bmatrix}
2 & \nicefrac{-1}{2} \\
-1 & \nicefrac{1}{2}
\end{bmatrix} .
\end{equation*}
Mari kita coba:
\begin{equation*}
\begin{bmatrix}
1 & 1 \\
2 & 4
\end{bmatrix}
\begin{bmatrix}
2 & \nicefrac{-1}{2} \\
-1 & \nicefrac{1}{2}
\end{bmatrix}
=
\begin{bmatrix}
1 & 0 \\
0 & 1
\end{bmatrix}
\qquad
\text{dan}
\qquad
\begin{bmatrix}
2 & \nicefrac{-1}{2} \\
-1 & \nicefrac{1}{2}
\end{bmatrix}
\begin{bmatrix}
1 & 1 \\
2 & 4
\end{bmatrix}
=
\begin{bmatrix}
1 & 0 \\
0 & 1
\end{bmatrix} .
\end{equation*}
Sebagaimana kita tidak dapat membagi dengan setiap bilangan, tidak setiap matriks
invertibel. Namun, dalam kasus matriks, kita dapat memiliki matriks
singular yang tidak nol. Sebagai contoh,
\begin{equation*}
\begin{bmatrix}
1 & 1 \\
2 & 2
\end{bmatrix}
\end{equation*}
adalah matriks singular. Namun, bukankah kita baru saja memberikan rumus untuk invers?
Mari kita coba:
\begin{equation*}
\begin{bmatrix}
1 & 1 \\
2 & 2
\end{bmatrix}^{-1}
=
\frac{1}{1\cdot 2-1 \cdot 2}
\begin{bmatrix}
2 & -1 \\
-2 & 1
\end{bmatrix}
=
%mbxSTARTIGNORE
\text{\Huge ?}
%mbxENDIGNORE
%mbxlatex \text{???}
\end{equation*}
Kita mengalami sedikit masalah:
Kita mencoba membagi dengan nol.

Jadi, matriks $2 \times 2$ $A$ invertibel apabila
\begin{equation*}
ad - bc \neq 0
\end{equation*}
dan jika tidak, matriks tersebut singular. Ungkapan $ad-bc$ disebut
\emph{determinan}, dan kita akan membahasnya
secara lebih cermat dalam bagian selanjutnya.
Terdapat ungkapan serupa untuk matriks persegi
berukuran sebarang.

\subsection{Matriks diagonal}

Jenis matriks persegi yang sederhana (dan secara mengejutkan berguna) adalah apa yang disebut
\emph{\myindex{matriks diagonal}}. Matriks ini memiliki entri yang semuanya nol
kecuali entri pada diagonal utama dari kiri atas ke kanan bawah. Sebagai
contoh, matriks diagonal $4 \times 4$ berbentuk
\begin{equation*}
\begin{bmatrix}
d_1 & 0 & 0 & 0 \\
0 & d_2 & 0 & 0 \\
0 & 0 & d_3 & 0 \\
0 & 0 & 0 & d_4
\end{bmatrix} .
\end{equation*}
Matriks semacam itu memiliki sifat-sifat yang baik ketika kita mengalikan dengannya. Jika kita mengalikan
matriks tersebut dengan suatu vektor, matriks itu mengalikan entri ke-$k$ dengan $d_k$. Sebagai
contoh,
\begin{equation*}
\begin{bmatrix}
1 & 0 & 0 \\
0 & 2 & 0 \\
0 & 0 & 3
\end{bmatrix}
\begin{bmatrix}
4 \\ 5 \\ 6
\end{bmatrix}
=
\begin{bmatrix}
1 \cdot 4 \\ 2 \cdot 5 \\ 3 \cdot 6
\end{bmatrix}
=
\begin{bmatrix}
4 \\ 10 \\ 18
\end{bmatrix} .
\end{equation*}
Secara serupa, ketika mengalikan matriks lain dari sebelah kiri, matriks diagonal mengalikan
baris ke-$k$ dengan $d_k$. Sebagai contoh,
\begin{equation*}
\begin{bmatrix}
2 & 0 & 0 \\
0 & 3 & 0 \\
0 & 0 & -1
\end{bmatrix}
\begin{bmatrix}
1 & 1 & 1 \\
1 & 1 & 1 \\
1 & 1 & 1
\end{bmatrix}
=
\begin{bmatrix}
2 & 2 & 2 \\
3 & 3 & 3 \\
-1 & -1 & -1
\end{bmatrix} .
\end{equation*}
Di sisi lain, jika dikalikan dari sebelah kanan, matriks diagonal mengalikan kolom-kolom:
\begin{equation*}
\begin{bmatrix}
1 & 1 & 1 \\
1 & 1 & 1 \\
1 & 1 & 1
\end{bmatrix}
\begin{bmatrix}
2 & 0 & 0 \\
0 & 3 & 0 \\
0 & 0 & -1
\end{bmatrix}
=
\begin{bmatrix}
2 & 3 & -1 \\
2 & 3 & -1 \\
2 & 3 & -1
\end{bmatrix} .
\end{equation*}
Dan sangat mudah mengalikan dua matriks diagonal---kita
mengalikan entri-entrinya:
\begin{equation*}
\begin{bmatrix}
1 & 0 & 0 \\
0 & 2 & 0 \\
0 & 0 & 3
\end{bmatrix}
\begin{bmatrix}
2 & 0 & 0 \\
0 & 3 & 0 \\
0 & 0 & -1
\end{bmatrix}
=
\begin{bmatrix}
1 \cdot 2 & 0 & 0 \\
0 & 2 \cdot 3 & 0 \\
0 & 0 & 3 \cdot (-1)
\end{bmatrix}
=
\begin{bmatrix}
2 & 0 & 0 \\
0 & 6 & 0 \\
0 & 0 & -3
\end{bmatrix} .
\end{equation*}

Untuk alasan terakhir ini, matriks diagonal mudah diinverskan---Anda cukup menginverskan setiap elemen diagonal:
\begin{equation*}
\begin{bmatrix}
d_1 & 0 & 0 \\
0 & d_2 & 0 \\
0 & 0 & d_3
\end{bmatrix}^{-1}
=
\begin{bmatrix}
d_1^{-1} & 0 & 0 \\
0 & d_2^{-1} & 0 \\
0 & 0 & d_3^{-1}
\end{bmatrix} .
\end{equation*}
Mari kita periksa sebuah contoh
\begin{equation*}
\underbrace{
\begin{bmatrix}
2 & 0 & 0 \\
0 & 3 & 0 \\
0 & 0 & 4
\end{bmatrix}^{-1}
}_{A^{-1}}
\underbrace{
\begin{bmatrix}
2 & 0 & 0 \\
0 & 3 & 0 \\
0 & 0 & 4
\end{bmatrix}
}_{A}
=
\underbrace{
\begin{bmatrix}
\frac{1}{2} & 0 & 0 \\
0 & \frac{1}{3} & 0 \\
0 & 0 & \frac{1}{4}
\end{bmatrix}
}_{A^{-1}}
\underbrace{
\begin{bmatrix}
2 & 0 & 0 \\
0 & 3 & 0 \\
0 & 0 & 4
\end{bmatrix}
}_{A}
=
\underbrace{
\begin{bmatrix}
1 & 0 & 0 \\
0 & 1 & 0 \\
0 & 0 & 1
\end{bmatrix}
}_{I} .
\end{equation*}
Tidak mengherankan bahwa cara kita menyelesaikan banyak masalah dalam aljabar linear
(dan dalam persamaan diferensial) adalah mencoba mereduksi masalah ke
kasus matriks diagonal.

\subsection{Transpos}

Vektor tidak harus selalu berupa vektor kolom.
Itu hanyalah suatu konvensi.
Dari waktu ke waktu,
menukar baris dan kolom diperlukan.
Operasi yang menukar baris dan kolom adalah apa yang disebut
\emph{\myindex{transpos}}.
Transpos dari $A$ dinotasikan dengan $A^T$. Contoh:
\begin{equation*}
\begin{bmatrix}
1 & 2 & 3 \\
4 & 5 & 6
\end{bmatrix}^T =
\begin{bmatrix}
1 & 4 \\
2 & 5 \\
3 & 6
\end{bmatrix} .
\end{equation*}
Transpos membawa matriks $m \times n$ menjadi matriks $n \times m$.

Fitur utama transpos adalah bahwa jika hasil kali $AB$ bermakna,
maka $B^TA^T$ juga bermakna, setidaknya ditinjau dari ukurannya.
Kenyataannya, kita memperoleh tepat transpos dari $AB$. Artinya:
\begin{equation*}
{(AB)}^T = B^TA^T .
\end{equation*}
Sebagai contoh,
\begin{equation*}
{\left(
\begin{bmatrix}
1 & 2 & 3 \\
4 & 5 & 6
\end{bmatrix}
\begin{bmatrix}
0 & 1 \\
1 & 0 \\
2 & -2
\end{bmatrix}
\right)}^T =
\begin{bmatrix}
0 & 1 & 2 \\
1 & 0 & -2
\end{bmatrix}
\begin{bmatrix}
1 & 4 \\
2 & 5 \\
3 & 6
\end{bmatrix} .
\end{equation*}
Diserahkan kepada pembaca untuk memverifikasi bahwa menghitung hasil kali matriks di
sebelah kiri lalu mentransposkannya sama dengan menghitung hasil kali matriks di
sebelah kanan.

Jika kita memiliki vektor kolom $\vec{x}$ yang dikenai matriks $A$
dan kita mentransposkan hasilnya,
maka vektor baris $\vec{x}^T$ diterapkan pada $A^T$ dari sebelah kiri:
\begin{equation*}
{(A\vec{x})}^T = \vec{x}^TA^T .
\end{equation*}

Tempat lain yang membuat transpos berguna adalah ketika kita ingin menerapkan hasil kali
titik\footnote{Sebagai catatan sampingan, matematikawan
menulis $\vec{y}^T\vec{x}$ dan fisikawan
menulis $\vec{x}^T\vec{y}$. Ssst\ldots jangan beri tahu siapa pun, tetapi para fisikawan
mungkin benar dalam hal ini.}
pada dua vektor kolom:
\begin{equation*}
\vec{x} \cdot \vec{y} = \vec{y}^T \vec{x} .
\end{equation*}
Begitulah hasil kali titik sering dituliskan dalam perangkat lunak.

Kita mengatakan suatu matriks $A$ \emph{simetris}\index{matriks simetris}
jika $A = A^T$. Sebagai contoh,
\begin{equation*}
\begin{bmatrix}
1 & 2 & 3 \\
2 & 4 & 5 \\
3 & 5 & 6
\end{bmatrix}
\end{equation*}
adalah matriks simetris. Perhatikan bahwa matriks simetris selalu
persegi, yakni $n \times n$. Matriks simetris
memiliki banyak sifat yang baik\footnote{Walaupun sejauh ini kita belum mempelajari
cukup banyak tentang matriks untuk benar-benar menghargainya.},
dan cukup sering muncul dalam penerapan.

\subsection{Latihan}

\begin{exercise}
Jumlahkan matriks-matriks berikut
\begin{tasks}(2)
\task
$\begin{bmatrix}
-1 & 2 & 2 \\
5 & 8 & -1
\end{bmatrix}
+
\begin{bmatrix}
3 & 2 & 3 \\
8 & 3 & 5
\end{bmatrix}$
\task
$\begin{bmatrix}
1 & 2 & 4 \\
2 & 3 & 1 \\
0 & 5 & 1
\end{bmatrix}
+
\begin{bmatrix}
2 & -8 & -3 \\
3 & 1 & 0 \\
6 & -4 & 1
\end{bmatrix}$
\end{tasks}
\end{exercise}

\begin{exercise}
Hitunglah
\begin{tasks}(2)
\task
$3\begin{bmatrix}
0 & 3 \\
-2 & 2
\end{bmatrix}
+
6
\begin{bmatrix}
1 & 5 \\
-1 & 5
\end{bmatrix}$
\task
$2\begin{bmatrix}
-3 & 1 \\
2 & 2
\end{bmatrix}
-
3
\begin{bmatrix}
2 & -1 \\
3 & 2
\end{bmatrix}$
\end{tasks}
\end{exercise}

\begin{exercise}
Kalikan matriks-matriks berikut
\begin{tasks}(2)
\task
$\begin{bmatrix}
-1 & 2 \\
3 & 1 \\
5 & 8
\end{bmatrix}
\begin{bmatrix}
3 & -1 & 3 & 1 \\
8 & 3 & 2 & -3
\end{bmatrix}$
\task
$\begin{bmatrix}
1 & 2 & 3 \\
3 & 1 & 1 \\
1 & 0 & 3
\end{bmatrix}
\begin{bmatrix}
2 & 3 & 1 & 7 \\
1 & 2 & 3 & -1 \\
1 & -1 & 3 & 0
\end{bmatrix}$
\task
$\begin{bmatrix}
4 & 1 & 6 & 3 \\
5 & 6 & 5 & 0 \\
4 & 6 & 6 & 0
\end{bmatrix}
\begin{bmatrix}
2 & 5 \\
1 & 2 \\
3 & 5 \\
5 & 6
\end{bmatrix}$
\task
$\begin{bmatrix}
1 & 1 & 4 \\
0 & 5 & 1
\end{bmatrix}
\begin{bmatrix}
2 & 2 \\
1 & 0 \\
6 & 4
\end{bmatrix}$
\end{tasks}
\end{exercise}

\begin{exercise}
Hitunglah invers matriks-matriks yang diberikan
\begin{tasks}(4)
\task
$\begin{bmatrix}
-3
\end{bmatrix}$
\task
$\begin{bmatrix}
0 & -1 \\
1 & 0
\end{bmatrix}$
\task
$\begin{bmatrix}
1 & 4 \\
1 & 3
\end{bmatrix}$
\task
$\begin{bmatrix}
2 & 2 \\
1 & 4
\end{bmatrix}$
\end{tasks}
\end{exercise}

\begin{exercise}
Hitunglah invers matriks-matriks yang diberikan
\begin{tasks}(3)
\task
$\begin{bmatrix}
-2 & 0 \\
0 & 1
\end{bmatrix}$
\task
$\begin{bmatrix}
3 & 0 & 0 \\
0 & -2 & 0 \\
0 & 0 & 1
\end{bmatrix}$
\task
$\begin{bmatrix}
1 & 0 & 0 & 0 \\
0 & -1 & 0 & 0 \\
0 & 0 & 0.01 & 0 \\
0 & 0 & 0 & -5
\end{bmatrix}$
\end{tasks}
\end{exercise}

\setcounter{exercise}{100}

\begin{exercise}
Jumlahkan matriks-matriks berikut
\begin{tasks}(2)
\task
$\begin{bmatrix}
2 & 1 & 0 \\
1 & 1 & -1
\end{bmatrix}
+
\begin{bmatrix}
5 & 3 & 4 \\
1 & 2 & 5
\end{bmatrix}$
\task
$\begin{bmatrix}
6 & -2 & 3 \\
7 & 3 & 3 \\
8 & -1 & 2
\end{bmatrix}
+
\begin{bmatrix}
-1 & -1 & -3 \\
6 & 7 & 3 \\
-9 & 4 & -1
\end{bmatrix}$
\end{tasks}
\end{exercise}
\exsol{%
a)~$\begin{bmatrix}
7 & 4 & 4 \\
2 & 3 & 4
\end{bmatrix}$
\quad
b)~$\begin{bmatrix}
5 & -3 & 0 \\
13 & 10 & 6 \\
-1 & 3 & 1
\end{bmatrix}$
}

\begin{exercise}
Hitunglah
\begin{tasks}(2)
\task
$2\begin{bmatrix}
1 & 2 \\
3 & 4
\end{bmatrix}
+
3
\begin{bmatrix}
-1 & 3 \\
1 & 2
\end{bmatrix}$
\task
$3\begin{bmatrix}
2 & -1 \\
1 & 3
\end{bmatrix}
-
2
\begin{bmatrix}
2 & 1 \\
-1 & 2
\end{bmatrix}$
\end{tasks}
\end{exercise}
\exsol{%
a)~$\begin{bmatrix}
-1 & 13 \\
9 & 14
\end{bmatrix}$
\quad
b)~$\begin{bmatrix}
2 & -5 \\
5 & 5
\end{bmatrix}$
}

\begin{exercise}
\pagebreak[2]
Kalikan matriks-matriks berikut
\begin{tasks}(2)
\task
$\begin{bmatrix}
2 & 1 & 4 \\
3 & 4 & 4
\end{bmatrix}
\begin{bmatrix}
2 & 4 \\
6 & 3 \\
3 & 5
\end{bmatrix}$
\task
$\begin{bmatrix}
0 & 3 & 3 \\
2 & -2 & 1 \\
3 & 5 & -2
\end{bmatrix}
\begin{bmatrix}
6 & 6 & 2 \\
4 & 6 & 0 \\
2 & 0 & 4
\end{bmatrix}$
\task
$\begin{bmatrix}
3 & 4 & 1 \\
2 & -1 & 0 \\
4 & -1 & 5
\end{bmatrix}
\begin{bmatrix}
0 & 2 & 5 & 0 \\
2 & 0 & 5 & 2 \\
3 & 6 & 1 & 6
\end{bmatrix}$
\task
$\begin{bmatrix}
-2 & -2 \\
5 & 3 \\
2 & 1
\end{bmatrix}
\begin{bmatrix}
0 & 3 \\
1 & 3
\end{bmatrix}$
\end{tasks}
\end{exercise}
\exsol{%
a)~$\begin{bmatrix}
22 & 31 \\
42 & 44
\end{bmatrix}$
\quad
b)~$\begin{bmatrix}
18 & 18 & 12 \\
6 & 0 & 8 \\
34 & 48 & -2
\end{bmatrix}$
\quad
c)~$\begin{bmatrix}
11 & 12 & 36 & 14 \\
-2 & 4 & 5 & -2 \\
13 & 38 & 20 & 28
\end{bmatrix}$
\quad
d)~$\begin{bmatrix}
-2 & -12 \\
3 & 24 \\
1 & 9
\end{bmatrix}$
}

\begin{exercise}
Hitunglah invers matriks-matriks yang diberikan
\begin{tasks}(4)
\task
$\begin{bmatrix}
2
\end{bmatrix}$
\task
$\begin{bmatrix}
0 & 1 \\
1 & 0
\end{bmatrix}$
\task
$\begin{bmatrix}
1 & 2 \\
3 & 5
\end{bmatrix}$
\task
$\begin{bmatrix}
4 & 2 \\
4 & 4
\end{bmatrix}$
\end{tasks}
\end{exercise}
\exsol{%
a)
$\begin{bmatrix}
\nicefrac{1}{2}
\end{bmatrix}$
\quad b)
$\begin{bmatrix}
0 & 1 \\
1 & 0
\end{bmatrix}$
\quad c)
$\begin{bmatrix}
-5 & 2 \\
3 & -1
\end{bmatrix}$
\quad d)
$\begin{bmatrix}
\nicefrac{1}{2} & \nicefrac{-1}{4} \\
\nicefrac{-1}{2} & \nicefrac{1}{2}
\end{bmatrix}$
}

\begin{exercise}
Hitunglah invers matriks-matriks yang diberikan
\begin{tasks}(3)
\task
$\begin{bmatrix}
2 & 0 \\
0 & 3
\end{bmatrix}$
\task
$\begin{bmatrix}
4 & 0 & 0 \\
0 & 5 & 0 \\
0 & 0 & -1
\end{bmatrix}$
\task
$\begin{bmatrix}
-1 & 0 & 0 & 0 \\
0 & 2 & 0 & 0 \\
0 & 0 & 3 & 0 \\
0 & 0 & 0 & 0.1
\end{bmatrix}$
\end{tasks}
\end{exercise}
\exsol{%
a)~$\begin{bmatrix}
\nicefrac{1}{2} & 0 \\
0 & \nicefrac{1}{3}
\end{bmatrix}$
\quad
b)~$\begin{bmatrix}
\nicefrac{1}{4} & 0 & 0 \\
0 & \nicefrac{1}{5} & 0 \\
0 & 0 & -1
\end{bmatrix}$
\quad
c)~$\begin{bmatrix}
-1 & 0 & 0 & 0 \\
0 & \nicefrac{1}{2} & 0 & 0 \\
0 & 0 & \nicefrac{1}{3} & 0 \\
0 & 0 & 0 & 10
\end{bmatrix}$
}



%%%%%%%%%%%%%%%%%%%%%%%%%%%%%%%%%%%%%%%%%%%%%%%%%%%%%%%%%%%%%%%%%%%%%%%%%%%%%%

\sectionnewpage
\section{Eliminasi}
\label{elim:section}

%mbxINTROSUBSECTION

\sectionnotes{2--3 kuliah}

\subsection{Sistem persamaan linear}

Salah satu penerapan matriks adalah menyelesaikan sistem
persamaan linear\footnote{Walaupun mungkin kita memandangnya secara terbalik;
cukup sering kita menyelesaikan sistem persamaan linear
untuk mengetahui sesuatu tentang matriks, bukan sebaliknya.}.
Perhatikan sistem persamaan linear berikut
\begin{equation} \label{linalg:elim:eq}
\begin{aligned}
          2 x_1 +           2 x_2 +           2 x_3 & = 2 , \\
\phantom{9} x_1 + \phantom{9} x_2 +           3 x_3 & = 5 , \\
\phantom{9} x_1 +           4 x_2 + \phantom{9} x_3 & = 10 .
\end{aligned}
\end{equation}

Terdapat prosedur sistematis
yang disebut \emph{\myindex{eliminasi}} untuk menyelesaikan sistem semacam itu.
Dalam prosedur ini,
kita berusaha mengeliminasi setiap variabel dari semua persamaan kecuali satu.
Kita ingin berakhir dengan persamaan
seperti $x_3 = 2$, yang jawabannya dapat langsung kita baca.

Kita menuliskan suatu sistem persamaan linear sebagai persamaan matriks:
\begin{equation*}
A \vec{x} = \vec{b} .
\end{equation*}
Sistem \eqref{linalg:elim:eq} ditulis sebagai
\begin{equation*}
\underbrace{
\begin{bmatrix}
2 & 2 & 2 \\
1 & 1 & 3 \\
1 & 4 & 1
\end{bmatrix}
}_{A}
\underbrace{
\begin{bmatrix}
x_1 \\
x_2 \\
x_3
\end{bmatrix}
}_{\vec{x}}
=
\underbrace{
\begin{bmatrix}
2 \\
5 \\
10
\end{bmatrix}
}_{\vec{b}} .
\end{equation*}
Jika kita mengetahui invers $A$, pekerjaan kita selesai.
Kita cukup menyelesaikan
persamaan tersebut:
\begin{equation*}
\vec{x} = A^{-1} A \vec{x} = A^{-1} \vec{b} .
\end{equation*}
Namun, justru itulah sebagian dari masalahnya.
Kita belum mengetahui cara menghitung
invers untuk matriks yang lebih besar daripada $2 \times 2$.
Nanti kita akan melihat bahwa untuk menghitung invers, sebenarnya kita menyelesaikan
$A \vec{x} = \vec{b}$ untuk beberapa $\vec{b}$ yang berbeda. Dengan kata lain,
kita perlu melakukan eliminasi untuk mencari $A^{-1}$.
Selain itu, kita mungkin ingin menyelesaikan $A \vec{x} = \vec{b}$
jika $A$ tidak invertibel, atau mungkin bahkan tidak persegi.

\medskip

Mari kita kembali ke persamaan-persamaannya dan melihat cara memanipulasinya.
Terdapat beberapa operasi yang dapat kita lakukan pada persamaan tanpa mengubah
solusinya. Pertama, mungkin suatu operasi yang tampak bodoh, kita dapat menukar
dua persamaan dalam \eqref{linalg:elim:eq}:
\begin{equation*}
\begin{aligned}
\phantom{9} x_1 + \phantom{9} x_2 +           3 x_3 & = 5 , \\
          2 x_1 +           2 x_2 +           2 x_3 & = 2 , \\
\phantom{9} x_1 +           4 x_2 + \phantom{9} x_3 & = 10 .
\end{aligned}
\end{equation*}
Jelas persamaan-persamaan baru ini memiliki solusi $x_1,x_2,x_3$ yang sama.
Operasi kedua adalah mengalikan suatu persamaan dengan bilangan tak nol. Sebagai
contoh, kita mengalikan persamaan ketiga dalam \eqref{linalg:elim:eq}
dengan 3:
\begin{equation*}
\begin{aligned}
          2 x_1 + \phantom{9}  2 x_2 + 2 x_3 & = 2 , \\
\phantom{9} x_1 + \phantom{99}   x_2 + 3 x_3 & = 5 , \\
          3 x_1 +             12 x_2 + 3 x_3 & = 30 .
\end{aligned}
\end{equation*}
Terakhir, kita dapat menambahkan kelipatan suatu persamaan ke persamaan lain.
Misalnya, kita menambahkan 3 kali persamaan ketiga dalam \eqref{linalg:elim:eq}
ke persamaan kedua:
\begin{equation*}
\begin{aligned}
\phantom{(1+3)} 2 x_1 + \phantom{(1+12)}  2 x_2 + \phantom{(3+3)} 2 x_3 & = 2 , \\
\phantom{2} (1+3) x_1 + \phantom{2}(1+12)   x_2 + \phantom{2} (3+3) x_3 & = 5+30 , \\
\phantom{2 (1+3)} x_1 + \phantom{(1+12)}  4 x_2 + \phantom{(3+3) 2} x_3 & = 10 .
\end{aligned}
\end{equation*}
$x_1,x_2,x_3$ yang sama seharusnya tetap menjadi solusi
persamaan-persamaan baru.
Ini hanyalah contoh; kita belum semakin dekat pada solusinya.
Kita harus melakukan ketiga operasi ini dengan cara yang lebih logis,
tetapi ternyata
ketiga operasi ini cukup untuk menyelesaikan setiap persamaan linear.

Hal pertama adalah menuliskan persamaan secara lebih ringkas. Diberikan
\begin{equation*}
A \vec{x} = \vec{b} ,
\end{equation*}
kita menuliskan apa yang disebut \emph{\myindex{matriks diperbesar}}
\begin{equation*}
[ A ~|~ \vec{b} ] ,
\end{equation*}
dengan garis vertikal hanyalah penanda agar kita mengetahui tempat
\myquote{ruas kanan} persamaan dimulai.
Untuk sistem \eqref{linalg:elim:eq}, matriks diperbesarnya adalah
\begin{equation*}
\left[
\begin{array}{ccc|c}
2 & 2 & 2 & 2 \\
1 & 1 & 3 & 5 \\
1 & 4 & 1 & 10
\end{array}
\right] .
\end{equation*}
Seluruh proses eliminasi, yang akan kita uraikan,
sering diterapkan pada sebarang jenis matriks,
bukan hanya
matriks diperbesar.
Cukup pandang matriks tersebut sebagai matriks $3 \times 4$
\begin{equation*}
\begin{bmatrix}
2 & 2 & 2 & 2 \\
1 & 1 & 3 & 5 \\
1 & 4 & 1 & 10
\end{bmatrix} .
\end{equation*}

\subsection{Bentuk eselon baris dan operasi elementer}

Kita menerapkan ketiga operasi di atas pada matriks. Kita menyebutnya
\emph{\myindex{operasi elementer}} atau
\emph{\myindex{operasi baris elementer}}\index{operasi baris}.
Dengan menerjemahkan operasi-operasi tersebut ke dalam konteks matriks,
operasinya menjadi:
\begin{enumerate}[(i)]
\item Tukarkan dua baris.
\item Kalikan suatu baris dengan bilangan tak nol.
\item Tambahkan kelipatan suatu baris ke baris lain.
\end{enumerate}
\pagebreak[2]
Kita menjalankan operasi-operasi ini sampai
mencapai keadaan yang jawabannya mudah dibaca,
atau sampai mencapai kontradiksi yang menunjukkan tidak adanya solusi.

Secara lebih khusus, kita menjalankan operasi-operasi tersebut sampai memperoleh apa yang disebut
\emph{\myindex{bentuk eselon baris}}\index{bentuk eselon}.
Mari kita sebut
entri tak nol pertama (dari kiri) dalam setiap baris sebagai
\emph{\myindex{entri utama}}.
Suatu matriks berada dalam \emph{bentuk eselon baris} jika
syarat-syarat berikut dipenuhi:
\begin{enumerate}[(i)]
\item Entri utama pada sebarang baris terletak tepat di sebelah kanan
entri utama pada baris di atasnya.
\item Setiap baris nol berada di bawah semua baris tak nol.
\item Semua entri utama adalah 1.
\end{enumerate}
Suatu matriks berada dalam \emph{\myindex{bentuk eselon baris tereduksi}}
jika, lebih lanjut, syarat berikut dipenuhi.
\begin{enumerate}[(i),resume]
\item Semua entri di atas suatu entri utama bernilai nol.
\end{enumerate}

Perhatikan bahwa definisi ini berlaku untuk matriks berukuran sebarang.

\begin{example}
Matriks-matriks berikut berada dalam bentuk eselon baris. Entri-entri
utamanya ditandai:
\begin{equation*}
\begin{bmatrix}
\mybxsm{1} & 2 & 9 & 3 \\
0 & 0 & \mybxsm{1} & 5 \\
0 & 0 & 0 & \mybxsm{1}
\end{bmatrix}
\qquad
\begin{bmatrix}
\mybxsm{1} & -1 & -3  \\
0 & \mybxsm{1} & 5  \\
0 & 0 & \mybxsm{1}
\end{bmatrix}
\qquad
\begin{bmatrix}
\mybxsm{1} & 2 & 1 \\
0 & \mybxsm{1} & 2 \\
0 & 0 & 0
\end{bmatrix}
\qquad
\begin{bmatrix}
0 & \mybxsm{1} & -5 & 2 \\
0 & 0 & 0 & \mybxsm{1} \\
0 & 0 & 0 & 0
\end{bmatrix}
\end{equation*}
Tidak satu pun matriks di atas berada dalam bentuk eselon baris
\emph{tereduksi}. Sebagai
contoh, dalam matriks pertama, entri-entri di atas entri utama kedua dan ketiga
tidak bernilai nol; entri-entri itu adalah 9, 3, dan 5.
Matriks-matriks berikut berada dalam bentuk eselon baris tereduksi. Entri-entri
utamanya ditandai:
\begin{equation*}
\begin{bmatrix}
\mybxsm{1} & 3 & 0 & 8 \\
0 & 0 & \mybxsm{1} & 6 \\
0 & 0 & 0 & 0
\end{bmatrix}
\qquad
\begin{bmatrix}
\mybxsm{1} & 0 & 2 &  0  \\
0 & \mybxsm{1} & 3 & 0  \\
0 & 0 & 0 & \mybxsm{1}
\end{bmatrix}
\qquad
\begin{bmatrix}
\mybxsm{1} & 0 & 3 \\
0 & \mybxsm{1} & -2 \\
0 & 0 & 0
\end{bmatrix}
\qquad
\begin{bmatrix}
0 & \mybxsm{1} & 2 & 0 \\
0 & 0 & 0 & \mybxsm{1} \\
0 & 0 & 0 & 0
\end{bmatrix}
\end{equation*}
\end{example}

Prosedur yang akan kita uraikan untuk mencari bentuk eselon baris tereduksi
dari suatu matriks
disebut \emph{\myindex{eliminasi Gauss--Jordan}}.
Bagian pertamanya, yang menghasilkan bentuk eselon baris, disebut
\emph{\myindex{eliminasi Gauss}} atau
\emph{\myindex{reduksi baris}}. Untuk beberapa masalah, bentuk eselon
baris sudah memadai, dan melakukan bagian pertama ini saja memerlukan sedikit lebih sedikit pekerjaan.

Untuk mencapai bentuk eselon baris, kita bekerja secara sistematis.
Kita bergerak kolom demi kolom, dimulai dari kolom pertama. Kita mencari
entri paling atas pada kolom pertama yang tidak nol, dan menyebutnya
\emph{\myindex{pivot}}. Jika tidak terdapat entri tak nol,
kita beralih ke kolom berikutnya. Kita menukar baris untuk menempatkan baris dengan
pivot sebagai baris pertama. Kita membagi baris pertama
dengan pivot agar entri pivot menjadi 1. Sekarang perhatikan semua baris di bawahnya
dan kurangkan kelipatan yang tepat dari baris pivot sehingga semua entri
di bawah pivot menjadi nol.

Setelah prosedur ini, kita melupakan bahwa kita memiliki baris pertama (baris itu kini sudah tetap),
dan melupakan kolom yang memuat pivot serta semua kolom nol yang mendahuluinya.
Di bawah baris pivot, semua entri pada kolom-kolom ini hanyalah
nol. Kemudian kita berfokus pada matriks yang lebih kecil dan mengulangi langkah-langkah
di atas.

Cara terbaik untuk menunjukkannya adalah dengan contoh, jadi mari kita kembali ke
contoh dari awal bagian ini.
Kita mempertahankan garis vertikal dalam matriks,
walaupun prosedur ini berlaku pada sebarang matriks, bukan hanya
matriks diperbesar.
Kita mulai dengan kolom pertama dan menemukan pivot, yang dalam
kasus ini adalah entri pertama pada kolom pertama.
\begin{equation*}
\left[
\begin{array}{ccc|c}
\mybxsm{2} & 2 & 2 & 2 \\
1 & 1 & 3 & 5 \\
1 & 4 & 1 & 10
\end{array}
\right]
\end{equation*}
Kita mengalikan baris pertama dengan
$\nicefrac{1}{2}$.
\begin{equation*}
\left[
\begin{array}{ccc|c}
\mybxsm{1} & 1 & 1 & 1 \\
1 & 1 & 3 & 5 \\
1 & 4 & 1 & 10
\end{array}
\right]
\end{equation*}
Kita mengurangkan baris pertama dari baris kedua dan ketiga (dua operasi
elementer).
\begin{equation*}
\left[
\begin{array}{ccc|c}
1 & 1 & 1 & 1 \\
0 & 0 & 2 & 4 \\
0 & 3 & 0 & 9
\end{array}
\right]
\end{equation*}
Untuk sementara, kita sudah selesai dengan kolom pertama dan baris pertama. Kita hampir
berpura-pura bahwa matriks tersebut tidak memiliki kolom pertama dan baris pertama.
\begin{equation*}
\left[
\begin{array}{ccc|c}
* & * & * & * \\
* & 0 & 2 & 4 \\
* & 3 & 0 & 9
\end{array}
\right]
\end{equation*}
Baiklah\@, perhatikan kolom kedua, dan perhatikan bahwa sekarang pivot berada pada
baris ketiga.
\begin{equation*}
\left[
\begin{array}{ccc|c}
1 & 1 & 1 & 1 \\
0 & 0 & 2 & 4 \\
0 & \mybxsm{3} & 0 & 9
\end{array}
\right]
\end{equation*}
Kita menukar baris.
\begin{equation*}
\left[
\begin{array}{ccc|c}
1 & 1 & 1 & 1 \\
0 & \mybxsm{3} & 0 & 9 \\
0 & 0 & 2 & 4
\end{array}
\right]
\end{equation*}
Lalu kita membagi baris pivot dengan 3.
\begin{equation*}
\left[
\begin{array}{ccc|c}
1 & 1 & 1 & 1 \\
0 & \mybxsm{1} & 0 & 3 \\
0 & 0 & 2 & 4
\end{array}
\right]
\end{equation*}
Kita tidak perlu mengurangkan apa pun karena semua yang berada di bawah pivot sudah
nol. Kita melanjutkan. Sekali lagi kita mulai mengabaikan baris kedua dan kolom
kedua, lalu berfokus pada
\begin{equation*}
\left[
\begin{array}{ccc|c}
* & * & * & * \\
* & * & * & * \\
* & * & 2 & 4
\end{array}
\right] .
\end{equation*}
Kita menemukan pivot, lalu membagi baris tersebut dengan 2:
\begin{equation*}
\left[
\begin{array}{ccc|c}
1 & 1 & 1 & 1 \\
0 & 1 & 0 & 3 \\
0 & 0 & \mybxsm{2} & 4
\end{array}
\right]
\qquad \to \qquad
\left[
\begin{array}{ccc|c}
1 & 1 & 1 & 1 \\
0 & 1 & 0 & 3 \\
0 & 0 & 1 & 2
\end{array}
\right] .
\end{equation*}
Matriks tersebut kini berada dalam bentuk eselon baris.

Persamaan yang bersesuaian dengan baris terakhir adalah $x_3 = 2$.
Kita mengetahui $x_3$ dan dapat
menyubstitusikannya ke dalam dua persamaan pertama untuk memperoleh persamaan bagi
$x_1$ dan $x_2$. Kemudian kita dapat melakukan hal yang sama dengan $x_2$, sampai kita menyelesaikan
ketiga variabel. Prosedur ini disebut
\emph{\myindex{substitusi balik}}, dan kita dapat melakukannya melalui operasi
elementer.
Kita mulai dari pivot paling bawah (entri utama dalam bentuk
eselon baris) dan mengurangkan kelipatan yang tepat dari baris di atasnya untuk
membuat semua entri di atas pivot ini menjadi nol.
Kemudian kita beralih ke pivot berikutnya, dan seterusnya.
Setelah selesai, kita akan memperoleh matriks dalam bentuk eselon baris tereduksi.

Kita lanjutkan contoh kita.
Kurangkan baris terakhir dari baris pertama untuk memperoleh
\begin{equation*}
\left[
\begin{array}{ccc|c}
1 & 1 & 0 & -1 \\
0 & 1 & 0 & 3 \\
0 & 0 & 1 & 2
\end{array}
\right] .
\end{equation*}
Entri di atas pivot pada baris kedua sudah nol.
Jadi, kita beralih ke pivot berikutnya, yaitu pivot pada baris kedua. Kita mengurangkan
baris ini dari baris teratas untuk memperoleh
\begin{equation*}
\left[
\begin{array}{ccc|c}
1 & 0 & 0 & -4 \\
0 & 1 & 0 & 3 \\
0 & 0 & 1 & 2
\end{array}
\right] .
\end{equation*}
Matriks tersebut berada dalam bentuk eselon baris tereduksi.

Jika sekarang kita menuliskan persamaan untuk $x_1,x_2,x_3$, kita memperoleh
\begin{equation*}
x_1 = -4, \qquad x_2 = 3, \qquad x_3 = 2 .
\end{equation*}
Dengan kata lain, kita telah menyelesaikan sistem tersebut.

\subsection{Solusi tak tunggal dan sistem tak konsisten}

Mungkin saja
solusi suatu sistem persamaan linear
tidak tunggal, atau tidak ada solusi. Andaikan sejenak
bahwa bentuk eselon baris yang kita peroleh adalah
\begin{equation*}
\left[
\begin{array}{ccc|c}
1 & 2 & 3 & 4 \\
0 & 0 & 1 & 3 \\
0 & 0 & 0 & 1
\end{array}
\right] .
\end{equation*}
Maka kita memiliki persamaan $0=1$ yang berasal dari baris terakhir. Hal itu
mustahil,
dan persamaan-persamaan tersebut disebut \emph{\myindex{tak konsisten}}. Tidak terdapat solusi
untuk $A \vec{x} = \vec{b}$.


Di sisi lain, jika kita memperoleh bentuk eselon baris
\begin{equation*}
\left[
\begin{array}{ccc|c}
1 & 2 & 3 & 4 \\
0 & 0 & 1 & 3 \\
0 & 0 & 0 & 0
\end{array}
\right] ,
\end{equation*}
maka tidak ada masalah dalam mencari solusi. Bahkan, kita akan menemukan terlalu
banyak solusi. Mari kita lanjutkan dengan substitusi balik (mengurangkan 3 kali baris ketiga
dari baris pertama) untuk mencari bentuk eselon baris tereduksi, dan mari kita tandai
pivot-pivotnya.
\begin{equation*}
\left[
\begin{array}{ccc|c}
\mybxsm{1} & 2 & 0 & -5 \\
0 & 0 & \mybxsm{1} & 3 \\
0 & 0 & 0 & 0
\end{array}
\right]
\end{equation*}
Baris terakhir semuanya nol; baris itu hanya menyatakan $0=0$ dan kita mengabaikannya.
Dua persamaan yang tersisa adalah
\begin{equation*}
x_1 + 2 x_2 = -5 , \qquad
x_3 = 3 .
\end{equation*}
Mari kita selesaikan variabel-variabel yang bersesuaian dengan
pivot, yakni $x_1$ dan $x_3$, karena terdapat pivot pada kolom pertama
dan kolom ketiga:
\begin{align*}
& x_1 = - 2 x_2 -5 , \\
& x_3 = 3 .
\end{align*}
Variabel $x_2$ dapat bernilai apa pun yang Anda inginkan, dan kita tetap memperoleh solusi.
$x_2$ disebut \emph{\myindex{variabel bebas}}.
Terdapat tak hingga banyak solusi, satu untuk setiap pilihan $x_2$.
Jika kita memilih $x_2=0$,
maka $x_1 = -5$ dan $x_3 = 3$ memberikan suatu solusi. Namun, kita juga memperoleh solusi
dengan memilih, misalnya, $x_2 = 1$, yang dalam hal ini $x_1 = -7$ dan $x_3 = 3$,
atau dengan memilih $x_2 = -5$, yang dalam hal ini $x_1 = 5$ dan $x_3 = 3$.

\medskip

Gagasan umumnya adalah bahwa
jika suatu baris semuanya nol pada kolom-kolom yang bersesuaian dengan
variabel, tetapi memiliki entri tak nol pada kolom yang bersesuaian dengan
ruas kanan $\vec{b}$, maka sistem tersebut tak konsisten dan tidak memiliki solusi.
Dengan kata lain, sistem itu tak konsisten jika Anda menemukan pivot pada
sebelah kanan garis vertikal yang digambar dalam matriks diperbesar. Jika tidak, sistem
tersebut \emph{\myindex{konsisten}}, dan setidaknya terdapat satu solusi.
\pagebreak[2]

Andaikan sistem tersebut konsisten (setidaknya terdapat satu solusi):
\begin{enumerate}[(i)]
\item Jika setiap kolom yang bersesuaian dengan suatu variabel memiliki elemen pivot,
maka solusinya tunggal.
\item Jika terdapat kolom-kolom yang bersesuaian dengan variabel tanpa pivot,
maka variabel-variabel tersebut adalah \emph{variabel bebas} yang dapat dipilih
secara sebarang, dan terdapat tak hingga banyak solusi.
\end{enumerate}

\medskip

Ketika $\vec{b} = \vec{0}$, kita memiliki apa yang disebut
\emph{\myindex{persamaan matriks homogen}}
\begin{equation*}
A \vec{x} = \vec{0} .
\end{equation*}
Tidak perlu menuliskan matriks diperbesar dalam
kasus ini. Karena operasi elementer tidak mengubah kolom nol,
kolom tersebut selalu tetap menjadi kolom nol. Lebih lanjut, $A \vec{x} = \vec{0}$ selalu memiliki
setidaknya satu solusi, yakni $\vec{x} = \vec{0}$. Sistem semacam itu
selalu konsisten. Sistem tersebut mungkin memiliki solusi lain: Jika kita menemukan
variabel bebas apa pun, kita memperoleh tak hingga banyak solusi.

Himpunan solusi $A \vec{x} = \vec{0}$ cukup sering muncul,
sehingga orang memberinya nama. Himpunan itu disebut
\emph{\myindex{ruang nol}} atau
\emph{\myindex{kernel}} dari $A$.
Salah satu tempat kernel muncul adalah invertibilitas suatu matriks persegi $A$.
Jika kernel $A$ memuat suatu vektor tak nol, maka kernel tersebut memuat
tak hingga banyak vektor (terdapat variabel bebas). Namun, kemudian
mustahil membalik $\vec{0}$, karena tak hingga banyak vektor dipetakan ke
$\vec{0}$, sehingga tidak ada vektor tunggal yang dibawa $A$ ke $\vec{0}$.
Jadi, jika kernel tak trivial, yaitu jika terdapat vektor tak nol apa pun dalam
kernel,
dengan kata lain, jika terdapat variabel bebas apa pun, atau dengan kata lain lagi,
jika bentuk eselon baris $A$ memiliki kolom tanpa pivot,
maka $A$ tidak invertibel. Kita akan kembali pada gagasan ini nanti.

\subsection{Kebebasan linear dan peringkat}

Jika baris-baris suatu matriks bersesuaian dengan persamaan, mungkin baik untuk mengetahui
berapa banyak persamaan yang sebenarnya kita perlukan untuk memperoleh himpunan solusi yang sama.
Secara serupa, jika kita menemukan sejumlah solusi bagi persamaan linear
$A \vec{x} = \vec{0}$, kita dapat bertanya apakah kita telah menemukan cukup banyak
sehingga semua solusi lain dapat dibentuk dari himpunan yang diberikan.
Konsep yang kita inginkan adalah kebebasan linear.
Konsep yang sama berguna untuk persamaan diferensial, misalnya
dalam \chapterref{ho:chapter}.

Diberikan vektor baris atau kolom $\vec{y}_1, \vec{y}_2, \ldots, \vec{y}_n$,
suatu \emph{\myindex{kombinasi linear}} adalah ungkapan berbentuk
\begin{equation*}
\alpha_1 \vec{y}_1 +
\alpha_2 \vec{y}_2 +
\cdots +
\alpha_n \vec{y}_n ,
\end{equation*}
dengan $\alpha_1, \alpha_2, \ldots, \alpha_n$ semuanya skalar.
Sebagai contoh,
$3 \vec{y}_1 + \vec{y}_2 - 5 \vec{y}_3$ adalah kombinasi linear
dari $\vec{y}_1$, $\vec{y}_2$, dan $\vec{y}_3$.

Kita telah melihat kombinasi linear sebelumnya. Ungkapan
\begin{equation*}
A \vec{x}
\end{equation*}
adalah kombinasi linear dari kolom-kolom $A$, sedangkan
\begin{equation*}
\vec{x}^T A = (A^T \vec{x})^T
\avoidbreak
\end{equation*}
adalah kombinasi linear dari baris-baris $A$.

Cara kombinasi linear muncul dalam kajian kita tentang persamaan
diferensial serupa dengan perhitungan berikut. Andaikan
$\vec{x}_1$, $\vec{x}_2$, \ldots, $\vec{x}_n$ adalah solusi
dari $A \vec{x}_1 = \vec{0}$,
$A \vec{x}_2 = \vec{0}$, \ldots,
$A \vec{x}_n = \vec{0}$.
Maka kombinasi linear
\begin{equation*}
\vec{y} = \alpha_1 \vec{x}_1 +
\alpha_2 \vec{x}_2 +
\cdots +
\alpha_n \vec{x}_n
\end{equation*}
adalah solusi dari $A \vec{y} = \vec{0}$:
\begin{multline*}
A \vec{y} =
A (\alpha_1 \vec{x}_1 +
\alpha_2 \vec{x}_2 +
\cdots +
\alpha_n \vec{x}_n )
=
\\
=
\alpha_1 A \vec{x}_1 +
\alpha_2 A \vec{x}_2 +
\cdots +
\alpha_n A \vec{x}_n
=
\alpha_1 \vec{0} +
\alpha_2 \vec{0} +
\cdots +
\alpha_n \vec{0} = \vec{0} .
\end{multline*}

Jadi jika kita menemukan cukup banyak solusi, kita telah memperoleh semua solusi. Pertanyaannya,
kapan kita telah menemukan cukup banyak solusi?

Kita mengatakan vektor-vektor $\vec{y}_1$, $\vec{y}_2$, \ldots, $\vec{y}_n$
\emph{\myindex{bebas linear}} jika satu-satunya solusi dari
\begin{equation*}
\alpha_1 \vec{x}_1 +
\alpha_2 \vec{x}_2 +
\cdots +
\alpha_n \vec{x}_n
=
\vec{0}
\end{equation*}
adalah $\alpha_1 = \alpha_2 = \cdots = \alpha_n = 0$.
Jika tidak, kita mengatakan vektor-vektor tersebut \emph{\myindex{bergantung linear}}.

Sebagai contoh, vektor-vektor
$\left[ \begin{smallmatrix} 1 \\ 2 \end{smallmatrix} \right]$
dan
$\left[ \begin{smallmatrix} 0 \\ 1 \end{smallmatrix} \right]$
bebas linear. Mari kita coba:
\begin{equation*}
\alpha_1
\begin{bmatrix} 1 \\ 2 \end{bmatrix}
+
\alpha_2
\begin{bmatrix} 0 \\ 1 \end{bmatrix}
=
\begin{bmatrix} \alpha_1 \\ 2 \alpha_1 + \alpha_2 \end{bmatrix}
=
\vec{0} =
\begin{bmatrix} 0 \\ 0 \end{bmatrix} .
\end{equation*}
Jadi $\alpha_1 = 0$, dan kemudian jelas bahwa $\alpha_2 = 0$ juga. Dengan
kata lain, kedua vektor itu bebas linear.

Jika suatu himpunan vektor bergantung linear, yakni, beberapa $\alpha_j$
tidak nol, maka kita dapat menyatakan satu vektor dalam vektor-vektor lainnya.
Misalkan $\alpha_1 \neq 0$. Karena
$\alpha_1 \vec{x}_1 +
\alpha_2 \vec{x}_2 +
\cdots +
\alpha_n \vec{x}_n
=
\vec{0}$, maka
\begin{equation*}
\vec{x}_1
=
\frac{-\alpha_2}{\alpha_1}
\vec{x}_2 -
\frac{-\alpha_3}{\alpha_1}
\vec{x}_3 +
\cdots +
\frac{-\alpha_n}{\alpha_1}
\vec{x}_n .
\end{equation*}
Sebagai contoh,
\begin{equation*}
2
\begin{bmatrix} 1 \\ 2 \\ 3 \end{bmatrix}
-4
\begin{bmatrix} 1 \\ 1 \\ 1 \end{bmatrix}
+
2 \begin{bmatrix} 1 \\ 0 \\ -1 \end{bmatrix}
=
\begin{bmatrix} 0 \\ 0 \\ 0 \end{bmatrix} ,
\end{equation*}
sehingga
\begin{equation*}
\begin{bmatrix} 1 \\ 2 \\ 3 \end{bmatrix}
=
2
\begin{bmatrix} 1 \\ 1 \\ 1 \end{bmatrix}
-
\begin{bmatrix} 1 \\ 0 \\ -1 \end{bmatrix} .
\end{equation*}

Anda mungkin telah memperhatikan bahwa menentukan $\alpha_j$ tersebut hanyalah menyelesaikan
persamaan linear, sehingga mungkin tidak mengherankan bahwa untuk memeriksa
apakah suatu himpunan vektor bebas linear kita menggunakan reduksi baris.

Diberikan suatu himpunan vektor, kita mungkin bukan hanya ingin mengetahui apakah
vektor-vektor itu bebas linear atau tidak; kita mungkin ingin menemukan suatu subhimpunan yang bebas
linear. Atau mungkin kita ingin menemukan vektor-vektor lain yang
menghasilkan kombinasi linear yang sama dan bebas linear.
Cara menentukannya ialah membentuk suatu matriks dari vektor-vektor kita. Jika kita
mempunyai vektor baris, kita menjadikannya baris-baris suatu matriks. Jika kita mempunyai
vektor kolom, kita menjadikannya kolom-kolom suatu matriks.
Himpunan semua kombinasi linear dari suatu himpunan vektor disebut
\emph{\myindex{rentang}} vektor-vektor tersebut.
\begin{equation*}
\operatorname{span} \bigl\{ \vec{x}_1, \vec{x}_2 , \ldots , \vec{x}_n \bigr\}
=
\bigl\{
\text{Himpunan semua kombinasi linear dari
$\vec{x}_1, \vec{x}_2 , \ldots , \vec{x}_n$}
\bigr\} .
\end{equation*}


Diberikan suatu matriks $A$, jumlah maksimum baris yang bebas linear disebut
\emph{\myindex{peringkat}} dari $A$, dan kita menulis
\myquote{$\operatorname{rank} A$} untuk peringkat tersebut.
Sebagai contoh,
\begin{equation*}
\operatorname{rank}
\begin{bmatrix}
1 & 1 & 1 \\
2 & 2 & 2 \\
-1 & -1 & -1
\end{bmatrix}
=
1 .
\end{equation*}
Baris kedua dan ketiga
merupakan kelipatan dari baris pertama. Kita tidak dapat memilih lebih dari satu baris sambil
tetap mempunyai himpunan yang bebas linear. Namun, berapakah
\begin{equation*}
\operatorname{rank}
\begin{bmatrix}
1 & 2 & 3 \\
4 & 5 & 6 \\
7 & 8 & 9
\end{bmatrix} \quad = \quad ?
\end{equation*}
Pertanyaan itu tampaknya lebih sulit dijawab. Dua baris
pertama bebas linear (yang satu bukan kelipatan yang lain),
sehingga peringkatnya setidaknya
dua. Jika kita menyusun persamaan untuk $\alpha_1$, $\alpha_2$, dan
$\alpha_3$, kita akan memperoleh suatu sistem dengan takhingga banyak solusi. Salah satu
solusinya ialah
\begin{equation*}
\begin{bmatrix}
1 & 2 & 3
\end{bmatrix} -2
\begin{bmatrix}
4 & 5 & 6
\end{bmatrix} +
\begin{bmatrix}
7 & 8 & 9
\end{bmatrix} =
\begin{bmatrix}
0 & 0 & 0
\end{bmatrix} .
\end{equation*}
Jadi himpunan yang terdiri atas ketiga baris tersebut bergantung linear; peringkatnya tidak mungkin
3. Oleh karena itu, peringkatnya adalah 2.

Namun, bagaimana kita dapat melakukannya secara lebih sistematis? Kita cari bentuk eselon
barisnya!
\begin{equation*}
\text{Bentuk eselon baris dari}
\quad
\begin{bmatrix}
1 & 2 & 3 \\
4 & 5 & 6  \\
7 & 8 & 9
\end{bmatrix}
\quad
\text{adalah}
\quad
\begin{bmatrix}
1 & 2 & 3 \\
0 & 1 & 2  \\
0 & 0 & 0
\end{bmatrix} .
\end{equation*}
Operasi baris elementer tidak mengubah himpunan kombinasi linear dari
baris-baris tersebut (itulah salah satu alasan utama operasi ini didefinisikan demikian).
Dengan kata lain, rentang baris-baris dari $A$ sama
dengan rentang baris-baris dari bentuk eselon baris $A$.
Secara khusus, jumlah baris yang bebas linear tetap sama.
Dan pada bentuk eselon baris, semua baris taknol bebas linear.
Hal ini tidak sulit dilihat.
Pertimbangkan kedua baris taknol dalam contoh di atas.
Misalkan kita
mencoba menentukan $\alpha_1$ dan $\alpha_2$
dalam
\begin{equation*}
\alpha_1
\begin{bmatrix}
1 & 2 & 3
\end{bmatrix}
+
\alpha_2
\begin{bmatrix}
0 & 1 & 2
\end{bmatrix} =
\begin{bmatrix}
0 & 0 & 0
\end{bmatrix} .
\end{equation*}
Karena kolom pertama
dari matriks eselon baris bernilai nol kecuali pada baris pertama, diperoleh
$\alpha_1 = 0$. Dengan menggunakan $\alpha_1=0$ dan melihat kolom
kedua, kita memperoleh $\alpha_2=0$.
Kita hanya mempunyai dua baris taknol,
dan keduanya bebas linear, jadi peringkat matriks tersebut adalah 2.

Rentang baris disebut \emph{\myindex{ruang baris}}.
Ruang baris dari $A$ dan dari bentuk eselon baris $A$ adalah sama.
Dalam contoh tersebut,
\begin{equation*}
\begin{split}
\text{ruang baris dari }
\begin{bmatrix}
1 & 2 & 3 \\
4 & 5 & 6 \\
7 & 8 & 9
\end{bmatrix}
& =
\operatorname{span}
\left\{
\begin{bmatrix}
1 & 2 & 3
\end{bmatrix}
,
\begin{bmatrix}
4 & 5 & 6
\end{bmatrix}
,
\begin{bmatrix}
7 & 8 & 9
\end{bmatrix}
\right\}
\\
& =
\operatorname{span}
\left\{
\begin{bmatrix}
1 & 2 & 3
\end{bmatrix}
,
\begin{bmatrix}
0 & 1 & 2
\end{bmatrix}
\right\} .
\end{split}
\end{equation*}

\medskip

Serupa dengan ruang baris, rentang kolom disebut
\emph{\myindex{ruang kolom}}.
\begin{equation*}
\text{ruang kolom dari }
\begin{bmatrix}
1 & 2 & 3 \\
4 & 5 & 6 \\
7 & 8 & 9
\end{bmatrix}
=
\operatorname{span}
\left\{
\begin{bmatrix}
1 \\ 4 \\ 7
\end{bmatrix}
,
\begin{bmatrix}
2 \\ 5 \\ 8
\end{bmatrix}
,
\begin{bmatrix}
3 \\ 6 \\ 9
\end{bmatrix}
\right\} .
\end{equation*}

Jadi mungkin berguna pula untuk menemukan jumlah kolom yang bebas linear
dari $A$. Salah satu caranya ialah menemukan jumlah baris yang bebas linear
dari $A^T$. Fakta yang sangat berguna ialah bahwa jumlah
kolom yang bebas linear
selalu sama dengan jumlah baris yang bebas linear:

\begin{theorem}
$\operatorname{rank} A = \operatorname{rank} A^T$
\end{theorem}

Secara khusus, untuk menemukan suatu himpunan kolom yang bebas linear kita perlu
melihat tempat pivot berada. Jika Anda mengingat pembahasan di atas, ketika menyelesaikan $A \vec{x}
= \vec{0}$, kuncinya adalah menemukan pivot.
Setiap kolom nonpivot bersesuaian dengan
variabel bebas. Artinya, kita dapat menyatakan kolom-kolom nonpivot dalam
kolom-kolom pivot. Mari kita lihat suatu contoh. Pertama-tama kita mereduksi suatu
matriks sebarang:
\begin{equation*}
\begin{bmatrix}
1 & 2 & 3 & 4 \\
2 & 4 & 5 & 6 \\
3 & 6 & 7 & 8
\end{bmatrix} .
\end{equation*}
Kita menemukan suatu pivot dan mereduksi baris-baris di bawahnya:
\begin{equation*}
\begin{bmatrix}
\mybxsm{1} & 2 & 3 & 4 \\
2 & 4 & 5 & 6 \\
3 & 6 & 7 & 8
\end{bmatrix}
\to
\begin{bmatrix}
\mybxsm{1} & 2 & 3 & 4 \\
0 & 0 & -1 & -2 \\
3 & 6 & 7 & 8
\end{bmatrix}
\to
\begin{bmatrix}
\mybxsm{1} & 2 & 3 & 4 \\
0 & 0 & -1 & -2 \\
0 & 0 & -2 & -4
\end{bmatrix} .
\end{equation*}
Kita menemukan pivot berikutnya, menjadikannya satu, lalu mengulangi prosesnya:
\begin{equation*}
\begin{bmatrix}
\mybxsm{1} & 2 & 3 & 4 \\
0 & 0 & \mybxsm{-1} & -2 \\
0 & 0 & -2 & -4
\end{bmatrix}
\to
\begin{bmatrix}
\mybxsm{1} & 2 & 3 & 4 \\
0 & 0 & \mybxsm{1} & 2 \\
0 & 0 & -2 & -4
\end{bmatrix}
\to
\begin{bmatrix}
\mybxsm{1} & 2 & 3 & 4 \\
0 & 0 & \mybxsm{1} & 2 \\
0 & 0 & 0 & 0
\end{bmatrix} .
\end{equation*}
Matriks terakhir adalah bentuk eselon baris dari matriks tersebut.
Pertimbangkan pivot-pivot yang kita tandai.
Kolom pivot adalah kolom pertama dan ketiga.
Semua kolom lain bersesuaian dengan variabel bebas ketika menyelesaikan
$A \vec{x} = \vec{0}$, sehingga semua kolom lain dapat dinyatakan dalam kolom pertama dan
ketiga. Dengan kata lain,
\begin{equation*}
\text{ruang kolom dari }
\begin{bmatrix}
1 & 2 & 3 & 4 \\
2 & 4 & 5 & 6 \\
3 & 6 & 7 & 8
\end{bmatrix}
=
\operatorname{span}
\left\{
\begin{bmatrix}
1 \\
2 \\
3
\end{bmatrix}
,
\begin{bmatrix}
2 \\
4 \\
6
\end{bmatrix}
,
\begin{bmatrix}
3 \\
5 \\
7
\end{bmatrix}
,
\begin{bmatrix}
4 \\
6 \\
8
\end{bmatrix}
\right\}
=
\operatorname{span}
\left\{
\begin{bmatrix}
1 \\
2 \\
3
\end{bmatrix}
,
\begin{bmatrix}
3 \\
5 \\
7
\end{bmatrix}
\right\} .
\end{equation*}
Kita mungkin dapat menggunakan pasangan kolom lain untuk memperoleh rentang yang sama, tetapi kolom
pertama dan ketiga dijamin berhasil karena keduanya merupakan kolom pivot.

Pembahasan di atas dapat dikembangkan menjadi bukti teorema tersebut jika
kita menginginkannya.
Karena setiap baris taknol
dalam bentuk eselon baris memuat suatu pivot, peringkatnya adalah jumlah
pivot, yang sama dengan jumlah maksimum kolom yang bebas
linear.

\medskip

Gagasan itu juga berlaku dalam arah sebaliknya. Misalkan kita mempunyai sekumpulan vektor kolom
dan kita hanya perlu menemukan suatu himpunan yang bebas linear. Sebagai contoh, misalkan
kita mulai dengan vektor-vektor
\begin{equation*}
\vec{v}_1 =
\begin{bmatrix}
1 \\
2 \\
3
\end{bmatrix}
,
\quad
\vec{v}_2 =
\begin{bmatrix}
2 \\
4 \\
6
\end{bmatrix}
,
\quad
\vec{v}_3 =
\begin{bmatrix}
3 \\
5 \\
7
\end{bmatrix}
,
\quad
\vec{v}_4 =
\begin{bmatrix}
4 \\
6 \\
8
\end{bmatrix} .
\end{equation*}
Vektor-vektor ini tidak bebas linear seperti yang kita lihat di atas. Secara khusus,
rentang $\vec{v}_1$ dan $\vec{v}_3$ sama dengan
rentang keempat vektor tersebut. Jadi $\vec{v}_2$ dan $\vec{v}_4$
keduanya dapat ditulis sebagai kombinasi linear dari $\vec{v}_1$ dan $\vec{v}_3$.
Hal yang sering muncul dalam praktik ialah kita memperoleh suatu himpunan vektor
yang rentangnya merupakan himpunan solusi suatu persoalan. Namun mungkin kita memperoleh terlalu
banyak vektor, dan ingin menyederhanakannya. Sebagai contoh di atas, semua vektor dalam
rentang
$\vec{v}_1, \vec{v}_2, \vec{v}_3, \vec{v}_4$ dapat ditulis sebagai
$\alpha_1 \vec{v}_1 + \alpha_2 \vec{v}_2 + \alpha_3 \vec{v}_3 + \alpha_4
\vec{v}_4$ untuk beberapa bilangan $\alpha_1,\alpha_2,\alpha_3,\alpha_4$. Namun,
setiap vektor tersebut juga dapat ditulis sebagai
$a \vec{v}_1 + b \vec{v}_3$ untuk dua bilangan $a$ dan $b$. Harus
diakui bahwa bentuk itu jauh lebih sederhana. Lebih lanjut, bilangan $a$ dan $b$ ini
unik. Hal ini akan dibahas lebih lanjut dalam bagian berikutnya.

Untuk menemukan himpunan bebas linear ini, kita cukup mengambil vektor-vektor kita
dan membentuk matriks $[ \vec{v}_1 ~ \vec{v}_2 ~ \vec{v}_3 ~ \vec{v}_4 ]$,
yakni matriks
\begin{equation*}
\begin{bmatrix}
1 & 2 & 3 & 4 \\
2 & 4 & 5 & 6 \\
3 & 6 & 7 & 8
\end{bmatrix} .
\end{equation*}
Kita jalankan mesin reduksi baris, masukkan matriks ini ke dalamnya, temukan
kolom-kolom pivot, lalu pilih kolom-kolom tersebut. Dalam kasus ini, $\vec{v}_1$ dan
$\vec{v}_3$.

\subsection{Menghitung invers}

Jika matriks $A$ persegi dan terdapat solusi tunggal
$\vec{x}$ bagi $A \vec{x} = \vec{b}$ untuk setiap $\vec{b}$ (tidak ada variabel
bebas), maka $A$ dapat diinverskan.
Hal ini ekuivalen dengan matriks $n \times n$ berupa $A$ mempunyai peringkat $n$.

Secara khusus, jika $A \vec{x} = \vec{b}$ maka $\vec{x} = A^{-1} \vec{b}$.
Sekarang kita hanya perlu menghitung $A^{-1}$. Kita tentu dapat
melakukan eliminasi setiap kali ingin menemukan $A^{-1} \vec{b}$, tetapi itu
tidak masuk akal. Pemetaan $A^{-1}$ bersifat linear dan
karena itu diberikan oleh suatu matriks, dan kita telah melihat bahwa untuk menentukan matriks tersebut
kita hanya perlu menemukan ke mana $A^{-1}$ memetakan vektor-vektor basis baku
$\vec{e}_1$,
$\vec{e}_2$, \ldots,
$\vec{e}_n$.

Artinya, untuk menemukan kolom pertama $A^{-1}$, kita menyelesaikan
$A \vec{x} = \vec{e}_1$, karena kemudian $A^{-1} \vec{e}_1 = \vec{x}$.
Untuk menemukan kolom kedua $A^{-1}$, kita menyelesaikan
$A \vec{x} = \vec{e}_2$. Demikian seterusnya. Sebenarnya kita hanya perlu melakukan $n$
eliminasi. Namun, prosesnya dapat dibuat lebih mudah lagi.
Jika dipikirkan, eliminasinya sama untuk
semua yang berada di sisi kiri matriks diperbesar. Dengan melakukan
$n$ eliminasi secara terpisah, kita akan mengulangi sebagian besar perhitungan.
Sebaiknya kerjakan semuanya sekaligus.

Oleh karena itu, untuk menemukan invers $A$, kita menulis suatu matriks diperbesar berukuran $n
\times 2n$, yaitu $[ \,A ~|~ I\, ]$, dengan $I$ matriks identitas
yang kolom-kolomnya tepat merupakan vektor-vektor basis baku.
Kemudian kita melakukan reduksi baris hingga mencapai bentuk eselon baris
tereduksi. Jika $A$ dapat diinverskan, maka pivot dapat ditemukan di setiap kolom $A$,
sehingga
bentuk eselon baris tereduksi dari $[ \,A ~|~ I\, ]$
berbentuk $[ \,I ~|~ A^{-1}\, ]$.
Lalu kita tinggal membaca invers $A^{-1}$.
Jika Anda tidak menemukan pivot di setiap
satu dari $n$ kolom pertama matriks diperbesar tersebut, maka
$A$ tidak dapat diinverskan.

Hal ini paling baik dilihat melalui contoh. Misalkan kita ingin menginverskan matriks
\begin{equation*}
\begin{bmatrix}
1 & 2 & 3 \\
2 & 0 & 1 \\
3 & 1 & 0
\end{bmatrix} .
\end{equation*}
Kita menulis matriks diperbesar dan mulai mereduksinya:
\begin{align*}
& \left[
\begin{array}{ccc|ccc}
\mybxsm{1} & 2 & 3 & 1 & 0 & 0\\
2 & 0 & 1 & 0 & 1 & 0 \\
3 & 1 & 0 & 0 & 0 & 1
\end{array}
\right]
\to
& &
\left[
\begin{array}{ccc|ccc}
\mybxsm{1} & 2 & 3 & 1 & 0 & 0\\
0 & -4 & -5 & -2 & 1 & 0 \\
0 & -5 & -9 & -3 & 0 & 1
\end{array}
\right]
\to
\\
\to
& \left[
\begin{array}{ccc|ccc}
\mybxsm{1} & 2 & 3 & 1 & 0 & 0\\
0 & \mybxsm{1} & \nicefrac{5}{4} & \nicefrac{1}{2} & \nicefrac{-1}{4} & 0 \\
0 & -5 & -9 & -3 & 0 & 1
\end{array}
\right]
\to
& &
\left[
\begin{array}{ccc|ccc}
\mybxsm{1} & 2 & 3 & 1 & 0 & 0\\
0 & \mybxsm{1} & \nicefrac{5}{4} & \nicefrac{1}{2} & \nicefrac{-1}{4} & 0 \\
0 & 0 & \nicefrac{-11}{4} & \nicefrac{-1}{2} & \nicefrac{-5}{4} & 1
\end{array}
\right]
\to
\\
\to
& \left[
\begin{array}{ccc|ccc}
\mybxsm{1} & 2 & 3 & 1 & 0 & 0\\
0 & \mybxsm{1} & \nicefrac{5}{4} & \nicefrac{1}{2} & \nicefrac{-1}{4} & 0 \\
0 & 0 & \mybxsm{1} & \nicefrac{2}{11} & \nicefrac{5}{11} & \nicefrac{-4}{11}
\end{array}
\right]
\to
& &
\left[
\begin{array}{ccc|ccc}
\mybxsm{1} & 2 & 0 & \nicefrac{5}{11} & \nicefrac{-15}{11} & \nicefrac{12}{11} \\
0 & \mybxsm{1} & 0 & \nicefrac{3}{11} & \nicefrac{-9}{11} & \nicefrac{5}{11} \\
0 & 0 & \mybxsm{1} & \nicefrac{2}{11} & \nicefrac{5}{11} & \nicefrac{-4}{11}
\end{array}
\right]
\to
\\
\to
& \left[
\begin{array}{ccc|ccc}
\mybxsm{1} & 0 & 0 & \nicefrac{-1}{11} & \nicefrac{3}{11} & \nicefrac{2}{11} \\
0 & \mybxsm{1} & 0 & \nicefrac{3}{11} & \nicefrac{-9}{11} & \nicefrac{5}{11} \\
0 & 0 & \mybxsm{1} & \nicefrac{2}{11} & \nicefrac{5}{11} & \nicefrac{-4}{11}
\end{array}
\right] .
\end{align*}
Jadi
\begin{equation*}
{\begin{bmatrix}
1 & 2 & 3 \\
2 & 0 & 1 \\
3 & 1 & 0
\end{bmatrix}}^{-1}
=
\begin{bmatrix}
\nicefrac{-1}{11} & \nicefrac{3}{11} & \nicefrac{2}{11} \\
\nicefrac{3}{11} & \nicefrac{-9}{11} & \nicefrac{5}{11} \\
\nicefrac{2}{11} & \nicefrac{5}{11} & \nicefrac{-4}{11}
\end{bmatrix} .
\end{equation*}
Tidak terlalu buruk, bukan? Mungkin lebih sulit daripada menginverskan matriks $2 \times 2$
yang mempunyai rumus sederhana, tetapi tidak terlalu buruk. Dalam praktik, hal ini
dilakukan secara efisien oleh komputer.

\subsection{Latihan}

\begin{samepage}
\begin{exercise}
Hitung bentuk eselon baris tereduksi untuk matriks-matriks berikut:
\begin{tasks}(4)
\task
$\begin{bmatrix}
1 & 3 & 1 \\
0 & 1 & 1
\end{bmatrix}$
\task
$\begin{bmatrix}
3 & 3 \\
6 & -3
\end{bmatrix}$
\task
$\begin{bmatrix}
3 & 6 \\
-2 & -3
\end{bmatrix}$
\task
$\begin{bmatrix}
6 & 6 & 7 & 7 \\
1 & 1 & 0 & 1
\end{bmatrix}$
\task
$\begin{bmatrix}
9 & 3 & 0 & 2 \\
8 & 6 & 3 & 6 \\
7 & 9 & 7 & 9
\end{bmatrix}$
\task
$\begin{bmatrix}
2 & 1 & 3 & -3 \\
6 & 0 & 0 & -1 \\
-2 & 4 & 4 & 3
\end{bmatrix}$
\task
$\begin{bmatrix}
6 & 6 & 5 \\
0 & -2 & 2 \\
6 & 5 & 6
\end{bmatrix}$
\task
$\begin{bmatrix}
0 & 2 & 0 & -1 \\
6 & 6 & -3 & 3 \\
6 & 2 & -3 & 5
\end{bmatrix}$
\end{tasks}
\end{exercise}
\end{samepage}

\begin{exercise}
Hitung invers dari matriks-matriks yang diberikan
\begin{tasks}(3)
\task
$\begin{bmatrix}
1 & 0 & 0 \\
0 & 0 & 1 \\
0 & 1 & 0
\end{bmatrix}$
\task
$\begin{bmatrix}
1 & 1 & 1 \\
0 & 2 & 1 \\
0 & 0 & 1
\end{bmatrix}$
\task
$\begin{bmatrix}
1 & 2 & 3 \\
2 & 0 & 1 \\
0 & 2 & 1
\end{bmatrix}$
\end{tasks}
\end{exercise}

\begin{exercise}
Selesaikan (temukan semua solusi), atau tunjukkan bahwa tidak ada solusi
\begin{tasks}(2)
\task
$\begin{aligned}
 4x_1+3x_2 & = -2 \\
 -x_1+\phantom{3} x_2 & = 4
\end{aligned}$
\task
$\begin{aligned}
  x_1+5x_2+3x_3 & = 7 \\
 8x_1+7x_2+8x_3 & = 8 \\
 4x_1+8x_2+6x_3 & = 4
\end{aligned}$
\task
$\begin{aligned}
 4x_1+8x_2+2x_3 & = 3 \\
 -x_1-2x_2+3x_3 & = 1 \\
 4x_1+8x_2 \phantom{{}+3x_3} & = 2
\end{aligned}$
\task
$\begin{aligned}
  x+2y+3z & = 4 \\
2  x-\phantom{2} y+3z & = 1 \\
3  x+\phantom{2} y+6z & = 6
\end{aligned}$
\end{tasks}
\end{exercise}

\begin{exercise}
Dengan menghitung invers,
selesaikan sistem-sistem berikut untuk $\vec{x}$.
\begin{tasks}(2)
\task
$\begin{bmatrix}
4 & 1 \\
-1 & 3
\end{bmatrix} \vec{x} =
\begin{bmatrix} 13 \\ 26 \end{bmatrix}$
\task
$\begin{bmatrix}
3 & 3 \\
3 & 4
\end{bmatrix} \vec{x} =
\begin{bmatrix} 2 \\ -1 \end{bmatrix}$
\end{tasks}
\end{exercise}

\begin{exercise} \label{exercise:rankmatrix}
Hitung peringkat matriks-matriks yang diberikan
\begin{tasks}(3)
\task
$\begin{bmatrix}
6 & 3 & 5 \\
1 & 4 & 1 \\
7 & 7 & 6
\end{bmatrix}$
\task
$\begin{bmatrix}
5 & -2 & -1 \\
3 & 0 & 6 \\
2 & 4 & 5
\end{bmatrix}$
\task
$\begin{bmatrix}
1 & 2 & 3 \\
-1 & -2 & -3 \\
2 & 4 & 6
\end{bmatrix}$
\end{tasks}
\end{exercise}

\begin{exercise}
Untuk matriks-matriks dalam \exerciseref{exercise:rankmatrix}, temukan
suatu himpunan bebas linear dari vektor baris yang merentang ruang baris
(vektor-vektor itu tidak harus berupa baris matriks tersebut).
\end{exercise}

\begin{exercise}
Untuk matriks-matriks dalam \exerciseref{exercise:rankmatrix}, temukan
suatu himpunan bebas linear dari kolom-kolom yang merentang ruang kolom.
Artinya, temukan kolom-kolom pivot matriks tersebut.
\end{exercise}

\begin{exercise}
Temukan suatu subhimpunan bebas linear dari vektor-vektor berikut yang mempunyai
rentang yang sama.
\begin{equation*}
\begin{bmatrix}
-1 \\ 1 \\ 2
\end{bmatrix}
, \quad
\begin{bmatrix}
2 \\ -2 \\ -4
\end{bmatrix}
, \quad
\begin{bmatrix}
-2 \\ 4 \\ 1
\end{bmatrix}
, \quad
\begin{bmatrix}
-1 \\ 3 \\ -2
\end{bmatrix}
\end{equation*}
\end{exercise}

\setcounter{exercise}{100}

\begin{exercise}
Hitung bentuk eselon baris tereduksi untuk matriks-matriks berikut:
\begin{tasks}(4)
\task
$\begin{bmatrix}
1 & 0 & 1 \\
0 & 1 & 0
\end{bmatrix}$
\task
$\begin{bmatrix}
1 & 2 \\
3 & 4
\end{bmatrix}$
\task
$\begin{bmatrix}
1 & 1 \\
-2 & -2
\end{bmatrix}$
\task
$\begin{bmatrix}
1 & -3 & 1 \\
4 & 6 & -2 \\
-2 & 6 & -2
\end{bmatrix}$
\task
$\begin{bmatrix}
2 & 2 & 5 & 2 \\
1 & -2 & 4 & -1 \\
0 & 3 & 1 & -2
\end{bmatrix}$
\task
$\begin{bmatrix}
-2 & 6 & 4 & 3 \\
6 & 0 & -3 & 0 \\
4 & 2 & -1 & 1
\end{bmatrix}$
\task
$\begin{bmatrix}
0 & 0 & 0 & 0 \\
0 & 0 & 0 & 0
\end{bmatrix}$
\task
$\begin{bmatrix}
1 & 2 & 3 & 3 \\
1 & 2 & 3 & 5
\end{bmatrix}$
\end{tasks}
\end{exercise}
\exsol{%
a)~$\begin{bmatrix}
1 & 0 & 1 \\
0 & 1 & 0
\end{bmatrix}$
\quad b)~$\begin{bmatrix}
1 & 0 \\
0 & 1
\end{bmatrix}$
\quad c)~$\begin{bmatrix}
1 & 1 \\
0 & 0
\end{bmatrix}$
\quad d)~$\begin{bmatrix}
1 & 0 & 0 \\
0 & 1 & -1/3 \\
0 & 0 & 0
\end{bmatrix}$
\quad e)~$\begin{bmatrix}
1 & 0 & 0 & 77/15 \\
0 & 1 & 0 & -2/15 \\
0 & 0 & 1 & -8/5
\end{bmatrix}$
\quad f)~$\begin{bmatrix}
1 & 0 & -1/2 & 0 \\
0 & 1 & 1/2 & 1/2 \\
0 & 0 & 0 & 0
\end{bmatrix}$
\quad g)~$\begin{bmatrix}
0 & 0 & 0 & 0 \\
0 & 0 & 0 & 0
\end{bmatrix}$
\quad h)~$\begin{bmatrix}
1 & 2 & 3 & 0 \\
0 & 0 & 0 & 1
\end{bmatrix}$
}

\begin{exercise}
Hitung invers dari matriks-matriks yang diberikan
\begin{tasks}(3)
\task
$\begin{bmatrix}
0 & 1 & 0 \\
-1 & 0 & 0 \\
0 & 0 & 1
\end{bmatrix}$
\task
$\begin{bmatrix}
1 & 1 & 1 \\
1 & 1 & 0 \\
1 & 0 & 0
\end{bmatrix}$
\task
$\begin{bmatrix}
2 & 4 & 0 \\
2 & 2 & 3 \\
2 & 4 & 1
\end{bmatrix}$
\end{tasks}
\end{exercise}
\exsol{%
a)~$\begin{bmatrix}
0 & -1 & 0 \\
1 & 0 & 0 \\
0 & 0 & 1
\end{bmatrix}$
\quad
b)~$\begin{bmatrix}
0 & 0 & 1 \\
0 & 1 & -1 \\
1 & -1 & 0
\end{bmatrix}$
\quad
c)~$\begin{bmatrix}
\nicefrac{5}{2} & 1 & -3 \\
-1 & \nicefrac{-1}{2} & \nicefrac{3}{2} \\
-1 & 0 & 1
\end{bmatrix}$
}

\begin{exercise}
Selesaikan (temukan semua solusi), atau tunjukkan bahwa tidak ada solusi
\begin{tasks}(2)
\task
$\begin{aligned}
4x_1+3x_2 & = -1 \\
5x_1+6x_2 & = 4
\end{aligned}$
\task
$\begin{aligned}
 5x+6y+5z & = 7 \\
 6x+8y+6z & = -1 \\
 5x+2y+5z & = 2
\end{aligned}$
\task
$\begin{aligned}
a+\phantom{5}b+\phantom{6}c & = -1 \\
a+5b+6c & = -1 \\
-2a+5b+6c & = 8
\end{aligned}$
\task
$\begin{aligned}
-2 x_1+2x_2+8x_3 & = 6 \\
x_2+\phantom{8}x_3 & = 2 \\
x_1+4x_2+\phantom{8}x_3 & = 7
\end{aligned}$
\end{tasks}
\end{exercise}
\exsol{%
a) $x_1=-2$, $x_2 = \nicefrac{7}{3}$
\quad
b)~tidak ada solusi
\quad
c)~$a = -3$, $b=10$, $c=-8$
\quad
d)~$x_3$ bebas, $x_1 = -1+3x_3$, $x_2 = 2-x_3$
}

\begin{exercise}
Dengan menghitung invers,
selesaikan sistem-sistem berikut untuk $\vec{x}$.
\begin{tasks}(2)
\task
$\begin{bmatrix}
-1 & 1 \\
3 & 3
\end{bmatrix} \vec{x} =
\begin{bmatrix} 4 \\ 6 \end{bmatrix}$
\task
$\begin{bmatrix}
2 & 7 \\
1 & 6
\end{bmatrix} \vec{x} =
\begin{bmatrix} 1 \\ 3 \end{bmatrix}$
\end{tasks}
\end{exercise}
\exsol{%
a)~$\begin{bmatrix} -1 \\ 3 \end{bmatrix}$ \quad
b)~$\begin{bmatrix} -3 \\ 1 \end{bmatrix}$
}

\begin{exercise} \label{exercise:rankmatrixans}
Hitung peringkat matriks-matriks yang diberikan
\begin{tasks}(3)
\task
$\begin{bmatrix}
7 & -1 & 6 \\
7 & 7 & 7 \\
7 & 6 & 2
\end{bmatrix}$
\task
$\begin{bmatrix}
1 & 1 & 1 \\
1 & 1 & 1 \\
2 & 2 & 2
\end{bmatrix}$
\task
$\begin{bmatrix}
0 & 3 & -1 \\
6 & 3 & 1 \\
4 & 7 & -1
\end{bmatrix}$
\end{tasks}
\end{exercise}
\exsol{%
a) 3 \quad b) 1 \quad c) 2
}

\begin{exercise}
Untuk matriks-matriks dalam \exerciseref{exercise:rankmatrixans}, temukan
suatu himpunan bebas linear dari vektor baris yang merentang ruang baris
(vektor-vektor itu tidak harus berupa baris matriks tersebut).
\end{exercise}
\exsol{%
a)~$\begin{bmatrix} 1 & 0 & 0\end{bmatrix}$,
$\begin{bmatrix} 0 & 1 & 0\end{bmatrix}$,
$\begin{bmatrix} 0 & 0 & 1\end{bmatrix}$
\quad
b)~$\begin{bmatrix} 1 & 1 & 1\end{bmatrix}$
\quad
c)~$\begin{bmatrix} 1 & 0 & \nicefrac{1}{3}\end{bmatrix}$,
$\begin{bmatrix} 0 & 1 & \nicefrac{-1}{3}\end{bmatrix}$
}

\begin{exercise}
Untuk matriks-matriks dalam \exerciseref{exercise:rankmatrixans}, temukan
suatu himpunan bebas linear dari kolom-kolom yang merentang ruang kolom.
Artinya, temukan kolom-kolom pivot matriks tersebut.
\end{exercise}
\exsol{%
a)~$\begin{bmatrix} 7 \\ 7 \\ 7\end{bmatrix}$,
$\begin{bmatrix} -1 \\ 7 \\ 6\end{bmatrix}$,
$\begin{bmatrix} 7 \\ 6 \\ 2\end{bmatrix}$
\quad
b)~$\begin{bmatrix} 1 \\ 1 \\ 2\end{bmatrix}$
\quad
c)~$\begin{bmatrix} 0 \\ 6 \\ 4\end{bmatrix}$,
$\begin{bmatrix} 3 \\ 3 \\ 7\end{bmatrix}$
}

\begin{exercise}
Temukan suatu subhimpunan bebas linear dari vektor-vektor berikut yang mempunyai
rentang yang sama.
\begin{equation*}
\begin{bmatrix}
0 \\ 0 \\ 0
\end{bmatrix}
, \quad
\begin{bmatrix}
3 \\ 1 \\ -5
\end{bmatrix}
, \quad
\begin{bmatrix}
0 \\ 3 \\ -1
\end{bmatrix}
, \quad
\begin{bmatrix}
-3 \\ 2 \\ 4
\end{bmatrix}
\end{equation*}
\end{exercise}
\exsol{%
$\begin{bmatrix}
3 \\ 1 \\ -5
\end{bmatrix}
,
\begin{bmatrix}
0 \\ 3 \\ -1
\end{bmatrix}$
}



%%%%%%%%%%%%%%%%%%%%%%%%%%%%%%%%%%%%%%%%%%%%%%%%%%%%%%%%%%%%%%%%%%%%%%%%%%%%%%

\sectionnewpage
\section{Subruang, dimensi, dan kernel}
\label{subspaces:section}

%mbxINTROSUBSECTION

\sectionnotes{1 kuliah}

\subsection{Subruang, basis, dan dimensi}

Kita sering mendapati diri sedang memperhatikan himpunan
solusi suatu persamaan linear $L\vec{x} = \vec{0}$ untuk suatu matriks $L$.
Artinya, kita tertarik pada kernel $L$.
Himpunan semua solusi tersebut mempunyai struktur yang baik: Himpunan itu tampak dan
berperilaku sangat mirip dengan suatu ruang Euklides ${\mathbb R}^k$.

Kita mengatakan bahwa suatu himpunan $S$ dari vektor-vektor di ${\mathbb R}^n$ adalah
\emph{\myindex{subruang}} jika
kapan pun $\vec{x}$ dan $\vec{y}$ merupakan anggota $S$ dan
$\alpha$ suatu skalar, maka
\begin{equation*}
\vec{x} + \vec{y}, \qquad \text{dan} \qquad \alpha \vec{x}
\end{equation*}
juga merupakan anggota $S$. Artinya, kita dapat menjumlahkan dan mengalikan dengan skalar
dan tetap berada di $S$. Jadi setiap kombinasi linear vektor-vektor dari
$S$ tetap berada di $S$. Itulah hakikat suatu subruang. Subruang adalah suatu subhimpunan
tempat kita dapat mengambil kombinasi linear dan tetap berakhir di dalam subhimpunan tersebut.
Akibatnya, rentang dari sejumlah vektor secara otomatis merupakan suatu subruang.

\begin{example} \label{example:simplesubspaces}
\begin{enumerate}[(i)]
\item
Jika kita mengambil $S = {\mathbb R}^n$, maka $S$ ini merupakan subruang dari
${\mathbb R}^n$. Menjumlahkan dua vektor mana pun dalam ${\mathbb R}^n$ menghasilkan vektor dalam
${\mathbb R}^n$, demikian pula perkalian dengan skalar.
\item
Himpunan $S' = \{ \vec{0} \}$, yakni,
himpunan yang hanya terdiri atas vektor nol, juga
merupakan subruang dari ${\mathbb R}^n$. Hanya ada satu vektor dalam
subruang ini, jadi kita hanya perlu memeriksa definisi untuk vektor itu, dan semuanya terpenuhi:
$\vec{0}+\vec{0} = \vec{0}$ dan $\alpha \vec{0} = \vec{0}$.
\item
Himpunan $S''$ dari semua vektor berbentuk
$(a,a)$ untuk setiap bilangan real $a$, misalnya $(1,1)$, $(3,3)$, atau $(-0.5,-0.5)$,
merupakan subruang dari ${\mathbb R}^2$. Menjumlahkan dua vektor semacam itu, misalnya
$(1,1)+(3,3) = (4,4)$, kembali menghasilkan vektor dengan bentuk yang sama; demikian pula
perkalian dengan suatu skalar, misalnya $8(1,1) = (8,8)$.
\end{enumerate}
\end{example}

Jika $S$ suatu subruang dan kita dapat menemukan $k$ vektor bebas linear dalam $S$,
\begin{equation*}
\vec{v}_1, \vec{v}_2, \ldots, \vec{v}_k ,
\end{equation*}
sedemikian sehingga setiap vektor lain dalam $S$ merupakan kombinasi linear dari $\vec{v}_1,
\vec{v}_2,\ldots, \vec{v}_k$,
maka himpunan
$\{ \vec{v}_1, \vec{v}_2, \ldots, \vec{v}_k \}$ disebut suatu
\emph{\myindex{basis}} dari $S$. Dengan kata lain, $S$
adalah rentang dari
$\{ \vec{v}_1, \vec{v}_2, \ldots, \vec{v}_k \}$
dan tidak ada subhimpunan yang lebih kecil dari vektor-vektor ini yang merentang $S$.
Kita mengatakan bahwa $S$ berdimensi $k$,
dan kita menulis
\begin{equation*}
\dim S = k .
\end{equation*}
Kita mempunyai teorema berikut.

\begin{theorem}
\pagebreak[2]
Jika $S \subset {\mathbb R}^n$ suatu subruang dan $S$ bukan subruang
trivial $\{ \vec{0} \}$, maka terdapat suatu bilangan bulat positif tunggal
$k$ (dimensinya) dan suatu basis (yang tidak tunggal)
$\{ \vec{v}_1, \vec{v}_2, \ldots, \vec{v}_k \}$, sedemikian sehingga setiap
$\vec{w}$ dalam $S$ dapat direpresentasikan secara tunggal oleh
\begin{equation*}
\vec{w} =
\alpha_1 \vec{v}_1 +
\alpha_2 \vec{v}_2 +
\cdots
+
\alpha_k \vec{v}_k ,
\avoidbreak
\end{equation*}
untuk beberapa skalar $\alpha_1$, $\alpha_2$, \ldots, $\alpha_k$.
Dengan \myquote{direpresentasikan secara tunggal} kita bermaksud bahwa
skalar-skalar $\alpha_1$, $\alpha_2$, \ldots, $\alpha_k$ ini tunggal.
\end{theorem}

Sebagaimana suatu vektor dalam ${\mathbb R}^k$ direpresentasikan oleh suatu $k$-tupel
bilangan, demikian pula suatu vektor dalam subruang berdimensi $k$ dari ${\mathbb R}^n$
direpresentasikan oleh suatu $k$-tupel bilangan yang tunggal, setidaknya setelah kita menetapkan suatu basis.
Basis yang berbeda akan memberikan $k$-tupel bilangan berbeda untuk vektor yang sama.

Perlu kita tegaskan kembali bahwa meskipun $k$ tunggal
(suatu subruang tidak mungkin mempunyai dua dimensi berbeda---setiap
basis mempunyai jumlah vektor yang sama),
himpunan vektor basis sama sekali tidak tunggal. Terdapat banyak
basis berbeda untuk setiap subruang tertentu. Menemukan basis yang tepat bagi
suatu subruang merupakan bagian besar dari apa yang dikerjakan dalam aljabar linear. Bahkan,
itulah yang banyak kita kerjakan dalam
persamaan diferensial linear, meskipun pada pandangan pertama
mungkin tidak tampak demikian.

\begin{example}
\begin{enumerate}[(i)]
\item
Basis baku
\begin{equation*}
\vec{e}_1, \vec{e}_2, \ldots, \vec{e}_n ,
\end{equation*}
merupakan suatu basis dari ${\mathbb R}^n$ (karena itu namanya).
Jadi seperti yang diharapkan,
\begin{equation*}
\dim {\mathbb R}^n = n .
\end{equation*}
\item
Kedua himpunan $\bigl\{ (1,0) , (0,1) \bigr\}$ dan
$\bigl\{ (1,1), (1,-1) \bigr\}$ merupakan
basis dari ${\mathbb{R}}^2$. Suatu vektor $(3,7)$ dapat
ditulis sebagai
$(3,7)=3 (1,0) + 7 (0,1)$
(dan tidak dengan cara lain menggunakan basis pertama),
dan dapat ditulis sebagai
$(3,7)=5 (1,1) - 2(1,-1)$
(dan tidak dengan cara lain menggunakan basis kedua).
\item
Subruang $\{ \vec{0} \}$ berdimensi $0$.
Basisnya hanyalah himpunan vektor kosong.
\item
Subruang $S''$ dari \exampleref{example:simplesubspaces}, yakni himpunan
vektor $(a,a)$, berdimensi~1. Salah satu basis yang mungkin hanyalah
$\bigl\{ (1,1) \bigr\}$, yaitu vektor tunggal
$(1,1)$: Setiap vektor dalam $S''$ dapat direpresentasikan oleh $a (1,1) =
(a,a)$. Demikian pula, basis lain yang mungkin adalah
$\bigl\{ (-1,-1) \bigr\}$. Maka
vektor $(a,a)$ akan direpresentasikan sebagai $(-a) (-1,-1)$.
\item
Himpunan $\bigl\{ (1,1) , (-1,-1) \bigr\}$ memang
merentang $S''$, tetapi bukan basis karena bukan himpunan bebas linear.
Akibatnya, suatu vektor seperti $(2,2)$ dapat ditulis dengan beberapa
cara menggunakan himpunan vektor ini:
$(2,2) = 2(1,1) + 0 (-1,-1)$,
atau
$(2,2) = 5(1,1) + 3 (-1,-1)$, dan seterusnya.
\end{enumerate}
\end{example}

Ruang baris dan ruang kolom suatu matriks juga merupakan contoh
subruang,
karena keduanya diberikan sebagai rentang vektor.
Kita dapat menggunakan
apa yang kita ketahui tentang peringkat, ruang baris, dan ruang kolom
dari bagian sebelumnya untuk menemukan suatu basis.

\begin{example}
Dalam bagian sebelumnya, kita mempertimbangkan matriks
\begin{equation*}
A =
\begin{bmatrix}
1 & 2 & 3 & 4 \\
2 & 4 & 5 & 6 \\
3 & 6 & 7 & 8
\end{bmatrix} .
\end{equation*}
Dengan menggunakan reduksi baris untuk menemukan kolom-kolom pivot, kita memperoleh
\begin{equation*}
\text{ruang kolom dari $A$} \left(
\begin{bmatrix}
1 & 2 & 3 & 4 \\
2 & 4 & 5 & 6 \\
3 & 6 & 7 & 8
\end{bmatrix}
\right)
=
\operatorname{span}
\left\{
\begin{bmatrix}
1 \\
2 \\
3
\end{bmatrix}
,
\begin{bmatrix}
3 \\
5 \\
7
\end{bmatrix}
\right\} .
\end{equation*}
Yang kita lakukan ialah menemukan suatu basis dari ruang kolom.
Basis ini mempunyai dua elemen, sehingga ruang kolom dari $A$ berdimensi dua.
Perhatikan bahwa peringkat $A$ adalah dua.
\end{example}

Kita akan mengikuti prosedur yang sama jika ingin menemukan suatu basis dari
subruang $X$ yang direntang oleh
\begin{equation*}
\begin{bmatrix}
1 \\
2 \\
3
\end{bmatrix}
,
\begin{bmatrix}
2 \\
4 \\
6
\end{bmatrix}
,
\begin{bmatrix}
3 \\
5 \\
7
\end{bmatrix}
,
\begin{bmatrix}
4 \\
6 \\
8
\end{bmatrix}
.
\end{equation*}
Kita cukup membentuk matriks $A$ dengan vektor-vektor ini sebagai kolom
dan mengulangi perhitungan di atas. Subruang $X$ kemudian merupakan ruang kolom dari
$A$.

\begin{example}
Pertimbangkan matriks
\begin{equation*}
L =
\begin{bmatrix}
{1} & 2 & 0 & 0 & 3 \\
0 & 0 & {1} & 0 & 4 \\
0 & 0 & 0 & {1} & 5
\end{bmatrix} .
\end{equation*}
Matriks tersebut sudah berada dalam bentuk eselon baris tereduksi.
Matriks itu berperingkat 3.
Ruang kolom adalah rentang dari kolom-kolom pivot.
Ruang tersebut adalah ruang berdimensi 3
\begin{equation*}
\text{ruang kolom dari $L$} =
\operatorname{span} \left\{
\begin{bmatrix}
1 \\
0 \\
0
\end{bmatrix}
,
\begin{bmatrix}
0 \\
1 \\
0
\end{bmatrix}
,
\begin{bmatrix}
0 \\
0 \\
1
\end{bmatrix}
\right\}
= {\mathbb{R}}^3 .
\end{equation*}
Ruang baris adalah ruang berdimensi 3
\begin{equation*}
\text{ruang baris dari $L$} =
\operatorname{span} \left\{
\begin{bmatrix}
1 & 2 & 0 & 0 & 3
\end{bmatrix}
,
\begin{bmatrix}
0 & 0 & 1 & 0 & 4
\end{bmatrix}
,
\begin{bmatrix}
0 & 0 & 0 & 1 & 5
\end{bmatrix}
\right\} .
\end{equation*}
Karena vektor-vektor ini mempunyai 5 komponen, kita memandang ruang baris dari $L$
sebagai suatu subruang dari ${\mathbb{R}}^5$.
\end{example}

Cara dimensi-dimensi tersebut muncul dalam contoh bukanlah
kebetulan. Karena jumlah vektor yang perlu kita ambil
selalu sama dengan jumlah pivot, dan jumlah pivot
adalah peringkat, kita memperoleh hasil berikut.

\begin{theorem}[Peringkat]
Dimensi ruang kolom dan dimensi ruang baris
dari suatu matriks $A$ sama-sama setara dengan peringkat $A$.
\end{theorem}

\subsection{Kernel}

Himpunan solusi dari suatu persamaan linear $L\vec{x} = \vec{0}$,
yaitu kernel dari $L$, merupakan suatu subruang:
Jika $\vec{x}$ dan $\vec{y}$ adalah solusi,
maka
\begin{equation*}
L(\vec{x}+\vec{y}) =
L\vec{x}+L\vec{y} =
\vec{0}+\vec{0} = \vec{0} ,
\qquad \text{dan} \qquad
L(\alpha \vec{x}) =
\alpha L \vec{x} =
\alpha \vec{0} = \vec{0}.
\end{equation*}
Jadi $\vec{x}+\vec{y}$ dan $\alpha \vec{x}$ merupakan solusi.
Dimensi kernel disebut \emph{\myindex{nulitas}} dari
matriks tersebut.

Gagasan yang sama mengatur solusi persamaan diferensial
linear. Kita mencoba mendeskripsikan kernel dari suatu operator diferensial
linear, dan karena kernel itu merupakan subruang, kita mencari suatu basis dari
kernel ini. Sebagian besar buku ini dikhususkan untuk menemukan basis-basis semacam itu.

Kernel suatu matriks sama dengan kernel dari bentuk eselon baris
tereduksinya. Untuk matriks dalam bentuk eselon baris tereduksi, kernel cukup mudah
ditemukan. Jika kita mengambil hasil kali suatu matriks $L$ dan suatu vektor $\vec{x}$,
maka setiap entri dalam
$\vec{x}$ bersesuaian dengan suatu kolom dari $L$, yakni kolom yang dikalikan oleh entri tersebut.
Untuk menemukan kernel,
pilih suatu
kolom nonpivot dan buat suatu vektor yang mempunyai $-1$ pada entri
yang bersesuaian dengan kolom nonpivot ini serta nol pada semua entri lain
yang bersesuaian dengan kolom-kolom nonpivot lainnya.
Kemudian, untuk semua entri
yang bersesuaian dengan kolom pivot, masukkan tepat nilai pada
baris yang bersesuaian dari kolom nonpivot agar vektor itu menjadi
solusi dari $L \vec{x} = \vec{0}$.
Prosedur ini paling baik dipahami melalui contoh.

\begin{example}
Pertimbangkan
\begin{equation*}
L =
\begin{bmatrix}
\mybxsm{1} & 2 & 0 & 0 & 3 \\
0 & 0 & \mybxsm{1} & 0 & 4 \\
0 & 0 & 0 & \mybxsm{1} & 5
\end{bmatrix} .
\end{equation*}
Matriks ini berada dalam bentuk eselon baris tereduksi; pivot-pivotnya ditandai.
Terdapat dua kolom nonpivot, sehingga kernel berdimensi 2, yakni,
kernel itu merupakan rentang dari 2 vektor. Mari kita temukan vektor pertama.
Kita melihat kolom nonpivot pertama, yaitu kolom ke-$2$,
dan menempatkan $-1$ pada
entri ke-$2$ dari vektor kita. Kita menempatkan $0$ pada entri ke-$5$
karena kolom ke-$5$ juga merupakan kolom nonpivot:
\begin{equation*}
\begin{bmatrix}
? \\ -1 \\ ? \\ ? \\ 0
\end{bmatrix} .
\end{equation*}
Mari kita isi sisanya. Ketika vektor ini mengenai baris pertama dari $L$, kita memperoleh
$-2$ ditambah $1$ kali apa pun yang menggantikan tanda tanya pertama. Kita
ingin memperoleh
nol ketika mengenai baris pertama,
jadi ganti tanda tanya pertama dengan $2$.
Untuk memperoleh nol ketika mengenai baris kedua dan ketiga, cukup dengan mengganti
tanda tanya yang tersisa dengan $0$. Sebenarnya kita sedang mengisikan kolom nonpivot
ke dalam entri-entri yang tersisa. Mari kita periksa sambil menandai bilangan mana ditempatkan
di mana:
\begin{equation*}
\begin{bmatrix}
1 & \mybxsm{2} & 0 & 0 & 3 \\
0 & \mybxsm{0} & 1 & 0 & 4 \\
0 & \mybxsm{0} & 0 & 1 & 5
\end{bmatrix}
\begin{bmatrix}
\mybxsm{2} \\ -1 \\ \mybxsm{0} \\ \mybxsm{0} \\ 0
\end{bmatrix}
=
\begin{bmatrix}
0 \\ 0 \\ 0
\end{bmatrix}
.
\end{equation*}
Berhasil! Bagaimana dengan vektor kedua? Kita mulai dengan
\begin{equation*}
\begin{bmatrix}
? \\ 0 \\ ? \\ ? \\ -1 .
\end{bmatrix}
\end{equation*}
Kita mengganti tanda tanya pertama dengan 3, tanda tanya kedua dengan 4, dan
tanda tanya ketiga dengan 5. Mari kita periksa, dengan menandai unsur-unsurnya seperti sebelumnya,
\begin{equation*}
\begin{bmatrix}
1 & 2 & 0 & 0 & \mybxsm{3} \\
0 & 0 & 1 & 0 & \mybxsm{4} \\
0 & 0 & 0 & 1 & \mybxsm{5}
\end{bmatrix}
\begin{bmatrix}
\mybxsm{3} \\ 0 \\ \mybxsm{4} \\ \mybxsm{5} \\ -1
\end{bmatrix}
=
\begin{bmatrix}
0 \\ 0 \\ 0
\end{bmatrix}
.
\end{equation*}
Terdapat dua kolom nonpivot, jadi kita hanya memerlukan dua vektor.
Kita telah menemukan suatu basis dari kernel. Jadi,
\begin{equation*}
\text{kernel dari $L$} =
\operatorname{span} \left\{
\begin{bmatrix}
2 \\ -1 \\ 0 \\ 0 \\ 0
\end{bmatrix}
,
\begin{bmatrix}
3 \\ 0 \\ 4 \\ 5 \\ -1
\end{bmatrix}
\right\}
\end{equation*}
\end{example}

Dalam menemukan suatu basis dari kernel, yang kita lakukan ialah menyatakan semua
solusi dari
$L \vec{x} = \vec{0}$ sebagai kombinasi linear dari beberapa vektor yang diberikan.

\pagebreak[2]
Prosedur untuk menemukan suatu basis dari kernel suatu matriks $L$:
\begin{enumerate}[(i)]
\item Temukan bentuk eselon baris tereduksi dari $L$.
\item Tuliskan suatu basis dari kernel seperti di atas, satu vektor untuk setiap
kolom nonpivot.
\end{enumerate}


Peringkat suatu matriks adalah dimensi ruang kolom, dan ruang itu
merupakan rentang dari kolom-kolom pivot, sedangkan kernel merupakan rentang dari vektor-vektor tertentu,
satu untuk setiap kolom nonpivot. Jadi jumlah kedua bilangan itu harus sama dengan jumlah
kolom.

\begin{theorem}[Peringkat--Nulitas]
Jika suatu matriks $A$ mempunyai $n$ kolom, peringkat $r$, dan nulitas $k$ (dimensi
kernel), maka
\begin{equation*}
n = r+k .
\end{equation*}
\end{theorem}

Teorema tersebut sangat berguna dalam penerapan. Teorema itu memungkinkan kita menghitung
peringkat $r$ jika nulitas $k$ diketahui dan sebaliknya, tanpa melakukan
pekerjaan tambahan.
Salah satu contoh penerapannya ialah versi sederhana dari apa yang disebut
\emph{\myindex{alternatif Fredholm}}. Hasil serupa berlaku untuk
persamaan diferensial. Pertimbangkan
\begin{equation*}
A \vec{x} = \vec{b} ,
\end{equation*}
dengan $A$ suatu matriks persegi $n \times n$.
Kemudian terdapat dua kemungkinan yang saling mengecualikan:
\begin{enumerate}[(i)]
\item
Terdapat suatu solusi taknol $\vec{x}$ bagi $A \vec{x} = \vec{0}$.
\item
Persamaan $A \vec{x} = \vec{b}$ mempunyai suatu solusi tunggal $\vec{x}$ untuk setiap
$\vec{b}$.
\end{enumerate}
Bagaimana teorema Peringkat--Nulitas berperan di sini? Jika $A$ mempunyai
suatu solusi taknol $\vec{x}$ bagi $A \vec{x} = \vec{0}$, maka nulitas $k$
positif. Namun, peringkat $r = n-k$ harus lebih kecil dari $n$.
Artinya, ruang kolom dari $A$ berdimensi kurang dari $n$, sehingga ruang itu
merupakan subruang yang tidak mencakup seluruh ${\mathbb{R}}^n$.
Jadi ${\mathbb{R}}^n$ harus
memuat suatu vektor $\vec{b}$ yang tidak berada dalam ruang kolom dari $A$. Bahkan, sebagian besar
vektor dalam ${\mathbb{R}}^n$ tidak berada dalam ruang kolom dari $A$.



\subsection{Latihan}

\begin{exercise}
Untuk himpunan-himpunan vektor berikut, temukan suatu basis bagi subruang yang direntang oleh
vektor-vektor tersebut, dan temukan dimensi subruangnya.
\begin{tasks}(3)
\task
$
\begin{bmatrix}
1 \\ 1 \\ 1
\end{bmatrix}
, \quad
\begin{bmatrix}
-1 \\ -1 \\ -1
\end{bmatrix}
$
\task
$
\begin{bmatrix}
1 \\ 0 \\ 5
\end{bmatrix}
, \quad
\begin{bmatrix}
0 \\ 1 \\ 0
\end{bmatrix}
, \quad
\begin{bmatrix}
0 \\ -1 \\ 0
\end{bmatrix}
$
\task
$
\begin{bmatrix}
-4 \\ -3 \\ 5
\end{bmatrix}
, \quad
\begin{bmatrix}
2 \\ 3 \\ 3
\end{bmatrix}
, \quad
\begin{bmatrix}
2 \\ 0 \\ 2
\end{bmatrix}
$
\task
$
\begin{bmatrix}
1 \\ 3 \\ 0
\end{bmatrix}
, \quad
\begin{bmatrix}
0 \\ 2 \\ 2
\end{bmatrix}
, \quad
\begin{bmatrix}
-1 \\ -1 \\ 2
\end{bmatrix}
$
\task
$
\begin{bmatrix}
1 \\ 3
\end{bmatrix}
, \quad
\begin{bmatrix}
0 \\ 2
\end{bmatrix}
, \quad
\begin{bmatrix}
-1 \\ -1
\end{bmatrix}
$
\task
$
\begin{bmatrix}
3 \\ 1 \\ 3
\end{bmatrix}
, \quad
\begin{bmatrix}
2 \\ 4 \\ -4
\end{bmatrix}
, \quad
\begin{bmatrix}
-5 \\ -5 \\ -2
\end{bmatrix}
$
\end{tasks}
\end{exercise}

\begin{exercise}
\pagebreak[2]
Untuk matriks-matriks berikut, temukan suatu basis bagi kernel (ruang nol).
\begin{tasks}(4)
\task
$\begin{bmatrix}
1 & 1 & 1 \\
1 & 1 & 5 \\
1 & 1 & -4
\end{bmatrix}$
\task
$\begin{bmatrix}
2 & -1 & -3 \\
4 & 0 & -4 \\
-1 & 1 & 2
\end{bmatrix}$
\task
$\begin{bmatrix}
-4 & 4 & 4 \\
-1 & 1 & 1 \\
-5 & 5 & 5
\end{bmatrix}$
\task
$\begin{bmatrix}
-2 & 1 & 1 & 1 \\
-4 & 2 & 2 & 2 \\
1 & 0 & 4 & 3
\end{bmatrix}$
\end{tasks}
\end{exercise}

\begin{exercise}
Misalkan suatu matriks $5 \times 5$ berupa $A$ mempunyai peringkat 3. Berapakah nulitasnya?
\end{exercise}

\begin{exercise}
Misalkan $X$ adalah himpunan semua vektor dari ${\mathbb{R}}^3$ yang
komponen ketiganya nol. Apakah $X$ suatu subruang? Jika ya, temukan suatu basis
dan dimensinya.
\end{exercise}

\begin{exercise}
Pertimbangkan suatu matriks persegi $A$, dan misalkan $\vec{x}$ suatu
vektor taknol sedemikian sehingga $A \vec{x} = \vec{0}$. Apa yang dinyatakan alternatif Fredholm
tentang keterinversan $A$?
\end{exercise}

\begin{exercise}
Pertimbangkan
\begin{equation*}
M =
\begin{bmatrix}
1 & 2 & 3 \\
2 & ? & ? \\
-1 & ? & ?
\end{bmatrix} .
\end{equation*}
Jika nulitas matriks ini adalah 2, isilah tanda tanya tersebut. Petunjuk: Berapakah
peringkatnya?
\end{exercise}

\setcounter{exercise}{100}

\begin{exercise}
Untuk himpunan-himpunan vektor berikut, temukan suatu basis bagi subruang yang direntang oleh
vektor-vektor tersebut, dan temukan dimensi subruangnya.
\begin{tasks}(3)
\task
$
\begin{bmatrix}
1 \\ 2
\end{bmatrix}
, \quad
\begin{bmatrix}
1 \\ 1
\end{bmatrix}
$
\task
$
\begin{bmatrix}
1 \\ 1 \\ 1
\end{bmatrix}
, \quad
\begin{bmatrix}
2 \\ 2 \\ 2
\end{bmatrix}
, \quad
\begin{bmatrix}
1 \\ 1 \\ 2
\end{bmatrix}
$
\task
$
\begin{bmatrix}
5 \\ 3 \\ 1
\end{bmatrix}
, \quad
\begin{bmatrix}
5 \\ -1 \\ 5
\end{bmatrix}
, \quad
\begin{bmatrix}
-1 \\ 3 \\ -4
\end{bmatrix}
$
\task
$
\begin{bmatrix}
2 \\ 2 \\ 4
\end{bmatrix}
, \quad
\begin{bmatrix}
2 \\ 2 \\ 3
\end{bmatrix}
, \quad
\begin{bmatrix}
4 \\ 4 \\ -3
\end{bmatrix}
$
\task
$
\begin{bmatrix}
1 \\ 0
\end{bmatrix}
, \quad
\begin{bmatrix}
2 \\ 0
\end{bmatrix}
, \quad
\begin{bmatrix}
3 \\ 0
\end{bmatrix}
$
\task
$
\begin{bmatrix}
1 \\ 0 \\ 0
\end{bmatrix}
, \quad
\begin{bmatrix}
2 \\ 0 \\ 0
\end{bmatrix}
, \quad
\begin{bmatrix}
0 \\ 1 \\ 2
\end{bmatrix}
$
\end{tasks}
\end{exercise}
\exsol{%
a)~$\begin{bmatrix}
1 \\ 2
\end{bmatrix}
,
\begin{bmatrix}
1 \\ 1
\end{bmatrix}$ berdimensi 2,
\quad
b)~$
\begin{bmatrix}
1 \\ 1 \\ 1
\end{bmatrix}
,
\begin{bmatrix}
1 \\ 1 \\ 2
\end{bmatrix}$ berdimensi 2,
\quad
c)~$
\begin{bmatrix}
5 \\ 3 \\ 1
\end{bmatrix}
,
\begin{bmatrix}
5 \\ -1 \\ 5
\end{bmatrix}
,
\begin{bmatrix}
-1 \\ 3 \\ -4
\end{bmatrix}
$ berdimensi 3,
\quad
d)~$
\begin{bmatrix}
2 \\ 2 \\ 4
\end{bmatrix}
,
\begin{bmatrix}
2 \\ 2 \\ 3
\end{bmatrix}
$ berdimensi 2,
\quad
e)~$\begin{bmatrix}
1 \\ 0
\end{bmatrix}$ berdimensi 1,
\quad
f)~$\begin{bmatrix}
1 \\ 0 \\ 0
\end{bmatrix}
,
\begin{bmatrix}
0 \\ 1 \\ 2
\end{bmatrix}
$ berdimensi~2
}

\begin{exercise}
Untuk matriks-matriks berikut, temukan suatu basis bagi kernel (ruang nol).
\begin{tasks}(4)
\task
$\begin{bmatrix}
2 & 6 & 1 & 9 \\
1 & 3 & 2 & 9 \\
3 & 9 & 0 & 9
\end{bmatrix}$
\task
$
\begin{bmatrix}
2 & -2 & -5 \\
-1 & 1 & 5 \\
-5 & 5 & -3
\end{bmatrix}$
\task
$
\begin{bmatrix}
1 & -5 & -4 \\
2 & 3 & 5 \\
-3 & 5 & 2
\end{bmatrix}$
\task
$
\begin{bmatrix}
0 & 4 & 4 \\
0 & 1 & 1 \\
0 & 5 & 5
\end{bmatrix}$
\end{tasks}
\end{exercise}
\exsol{%
a)~$\begin{bmatrix} 3 \\ -1 \\ 0 \\ 0 \end{bmatrix}$, $\begin{bmatrix} 3 \\
0 \\ 3 \\ -1 \end{bmatrix}$
\quad
b)~$\begin{bmatrix} -1 \\ -1 \\ 0 \end{bmatrix}$
\quad
c)~$\begin{bmatrix} 1 \\ 1 \\ -1 \end{bmatrix}$
\quad
d)~$\begin{bmatrix} -1 \\ 0 \\ 0 \end{bmatrix}$, $\begin{bmatrix} 0 \\ 1 \\ -1 \end{bmatrix}$
}

\begin{exercise}
Misalkan ruang kolom dari suatu matriks $9 \times 5$ berupa $A$ berdimensi 3. Temukan
\begin{tasks}(2)
\task
Peringkat $A$.
\task
Nulitas $A$.
\task
Dimensi ruang baris dari $A$.
\task
Dimensi ruang nol dari $A$.
\task
Ukuran subhimpunan maksimum dari
baris-baris $A$ yang bebas linear.
\end{tasks}
\end{exercise}
\exsol{%
a)~3 \quad b)~2 \quad c)~3 \quad d)~2 \quad e)~3
}

%%%%%%%%%%%%%%%%%%%%%%%%%%%%%%%%%%%%%%%%%%%%%%%%%%%%%%%%%%%%%%%%%%%%%%%%%%%%%%

\sectionnewpage
\section{Hasil kali dalam dan proyeksi}
\label{innerproduct:section}

%mbxINTROSUBSECTION

\sectionnotes{1--2 kuliah}

\subsection{Hasil kali dalam dan ortogonalitas}

Untuk melakukan geometri dasar, kita memerlukan panjang dan sudut.
Kita telah melihat panjang Euklides, jadi mari kita tentukan sudut.
Sebagian besar, kita memperhatikan
sudut siku-siku\footnote{Ketika Euklides mendefinisikan sudut dalam
\emph{Elements}, satu-satunya sudut yang benar-benar ia definisikan adalah sudut siku-siku.}.

Diberikan dua vektor (kolom) di ${\mathbb{R}}^n$,
kita mendefinisikan
\emph{\myindex{hasil kali dalam}} (baku)\index{hasil kali dalam baku} sebagai
hasil kali titik:
\begin{equation*}
\langle \vec{x} , \vec{y} \rangle =
\vec{x} \cdot \vec{y}
=
\vec{y}^T \vec{x}
=
x_1 y_1 + x_2 y_2 + \cdots + x_n y_n
=
\sum_{i=1}^n x_i y_i .
\end{equation*}
Mengapa tampaknya kita memberikan notasi baru untuk hasil kali titik
($\langle \vec{x} , \vec{y} \rangle$ alih-alih hanya
$\vec{x} \cdot \vec{y}$)?
Karena terdapat hasil kali dalam lain yang mungkin, yang bukan hasil kali
titik, meskipun kita tidak akan membahas hasil kali lain di sini.
Hasil kali dalam bahkan dapat didefinisikan pada ruang fungsi,
seperti yang kita lakukan dalam
\chapterref{FS:chapter}:
\begin{equation*}
\langle f(t) , g(t) \rangle =
\int_{a}^{b}
f(t) g(t) \, dt .
\end{equation*}
Namun kita telah menyimpang dari pokok bahasan.

Hasil kali dalam memenuhi aturan-aturan berikut:
\begin{enumerate}[(i)]
\item $\langle \vec{x} , \vec{x} \rangle \geq 0$, dan
$\langle \vec{x} , \vec{x} \rangle = 0$ jika dan hanya jika $\vec{x} = 0$,
\item $\langle \vec{x} , \vec{y} \rangle = \langle \vec{y} , \vec{x}
\rangle$,
\item $\langle a\vec{x} , \vec{y} \rangle =
\langle \vec{x} , a\vec{y} \rangle =
a \langle \vec{x} , \vec{y} \rangle$,
\item $\langle \vec{x} +  \vec{y} , \vec{z} \rangle =
\langle \vec{x} , \vec{z} \rangle +
\langle \vec{y} , \vec{z} \rangle$ dan
$\langle \vec{x}, \vec{y} + \vec{z} \rangle =
\langle \vec{x} , \vec{y} \rangle +
\langle \vec{x} , \vec{z} \rangle$.
\end{enumerate}
Apa pun yang memenuhi sifat-sifat di atas dapat disebut hasil kali
dalam, meskipun dalam bagian ini kita memperhatikan
hasil kali dalam baku di ${\mathbb{R}}^n$.

\medskip

Hasil kali dalam baku menghasilkan panjang Euklides:
\begin{equation*}
\lVert{\vec{x}}\rVert = \sqrt{\langle \vec{x}, \vec{x} \rangle}
= \sqrt{x_1^2 + x_2^2 + \cdots + x_n^2} .
\end{equation*}
Bagaimana hasil kali dalam menghasilkan sudut?
Anda mungkin mengingat suatu rumus
untuk hasil kali dalam baku (hasil kali titik)
dari kalkulus multivariabel dalam
dua atau tiga dimensi,
dalam suku-suku
sudut $\theta$ antara kedua vektor:
\begin{equation*}
\langle \vec{x}, \vec{y} \rangle
=
\lVert{\vec{x}}\rVert \lVert{\vec{y}}\rVert \cos \theta.
\avoidbreak
\end{equation*}
Artinya, $\theta$ adalah sudut yang dibentuk oleh $\vec{x}$ dan $\vec{y}$
ketika keduanya berpangkal pada titik yang sama.

Di ${\mathbb{R}}^n$ (dimensi berapa pun), kita cukup mengatakan bahwa $\theta$
dari rumus tersebut adalah sudutnya.
Hal ini masuk akal karena dua vektor mana pun yang berpangkal di titik asal
terletak dalam suatu bidang berdimensi 2 (subruang),
dan rumus tersebut berlaku dalam 2 dimensi.
Sebenarnya, dengan cara ini kita bahkan dapat membicarakan sudut antara fungsi, dan
kita melakukannya dalam \chapterref{FS:chapter}, ketika membicarakan fungsi ortogonal
(fungsi yang saling membentuk sudut siku-siku).

Untuk menghitung sudut, kita selesaikan untuk $\cos \theta$:
\begin{equation*}
\cos \theta
=
\frac{\langle \vec{x}, \vec{y} \rangle}{\lVert{\vec{x}}\rVert
\lVert{\vec{y}}\rVert} .
\end{equation*}
Sudut-sudut kita selalu dinyatakan dalam radian.
Kita menghitung kosinus sudut tersebut,
yang sebenarnya merupakan hasil terbaik
yang dapat kita peroleh. Diberikan dua vektor dengan sudut $\theta$, kita dapat menyatakan sudutnya sebagai
$-\theta$, $2\pi-\theta$, dan seterusnya.
Lihat \figurevref{vec-angle:fig}.
Untungnya,
$\cos \theta = \cos (-\theta) = \cos(2\pi - \theta)$.
Jika kita menentukan $\theta$ menggunakan kosinus invers $\cos^{-1}$,
kita cukup menetapkan bahwa $0 \leq \theta \leq \pi$.

\begin{myfig}
\capstart
\inputpdft{vec-angle}{%
Diagram dengan dua vektor x dan y yang berawal dari titik yang sama,
x mengarah ke atas dan ke kanan, sedangkan y mengarah ke bawah dan ke kanan.
Sudut dari y ke x berlawanan arah jarum jam ditandai theta.
Sudut dari x ke y searah jarum jam ditandai minus theta.
Sudut dari x ke y berlawanan arah jarum jam ditandai 2 pi minus theta.}
\caption{Sudut antara vektor-vektor.\label{vec-angle:fig}}
\end{myfig}

\begin{example}
Mari kita hitung sudut antara vektor $(3,0)$ dan $(1,1)$ di
bidang.
Hitung
\begin{equation*}
\cos \theta =
\frac{\bigl\langle (3,0) , (1,1) \bigr\rangle}{\lVert(3,0)\rVert \lVert(1,1)\rVert}
=
\frac{3 + 0}{3 \sqrt{2}} = \frac{1}{\sqrt{2}} .
\end{equation*}
Oleh karena itu, $\theta = \nicefrac{\pi}{4}$.
\end{example}

Seperti yang telah kita katakan, sudut paling penting adalah sudut siku-siku. Sudut siku-siku
adalah $\nicefrac{\pi}{2}$ radian, dan $\cos (\nicefrac{\pi}{2}) = 0$, sehingga
rumusnya sangat mudah dalam kasus ini.
Kita mengatakan
vektor-vektor $\vec{x}$ dan $\vec{y}$
\emph{\myindex{ortogonal}} jika saling membentuk sudut siku-siku,
yakni, jika
\begin{equation*}
\langle \vec{x} , \vec{y} \rangle
=
0 .
\end{equation*}
Vektor-vektor $(1,0,0,1)$ dan $(1,2,3,-1)$ ortogonal. Demikian pula
$(1,1)$ dan $(1,-1)$. Namun, $(1,1)$ dan $(1,2)$ tidak ortogonal karena
hasil kali dalamnya adalah $3$, bukan 0.

\subsection{Proyeksi ortogonal}

Penerapan khas aljabar linear ialah mengambil suatu persoalan sulit,
menuliskan semuanya dalam basis yang tepat, dan dalam basis baru ini persoalan tersebut
menjadi sederhana. Basis yang sangat berguna ialah basis ortogonal,
yaitu basis yang semua vektor basisnya ortogonal. Ketika menggambar suatu
sistem koordinat dalam dua atau tiga dimensi, kita hampir selalu menggambar sumbu-sumbunya
saling ortogonal.

Dengan memperumum konsep ini pada fungsi,
dalam \chapterref{FS:chapter} sangat berguna untuk menyatakan suatu
fungsi menggunakan basis ortogonal tertentu, yaitu deret Fourier.

Untuk menyatakan satu vektor dalam suatu basis ortogonal, pertama-tama kita perlu
\emph{memproyeksikan} satu vektor pada vektor lain.
Diberikan suatu vektor taknol $\vec{v}$, kita mendefinisikan
\emph{\myindex{proyeksi ortogonal}}\index{proyeksi!ortogonal}
dari $\vec{w}$ pada $\vec{v}$ sebagai
\begin{equation*}
\operatorname{proj}_{\vec{v}}(\vec{w})
=
\left(
\frac{\langle \vec{w} , \vec{v} \rangle}{ \langle \vec{v} , \vec{v} \rangle}
\right)
\vec{v} .
\end{equation*}
Untuk gagasan geometrisnya, lihat \figurevref{vec-orthoproj:fig}. Artinya, kita
menemukan \myquote{bayangan $\vec{w}$} pada garis yang direntang oleh $\vec{v}$ jika
arah sinar matahari
tepat tegak lurus terhadap garis tersebut.
Cara lain untuk memikirkannya ialah bahwa ujung anak panah dari
$\operatorname{proj}_{\vec{v}}(\vec{w})$ merupakan titik pada garis
yang direntang oleh $\vec{v}$ yang paling dekat dengan ujung anak panah $\vec{w}$.
Dalam jarak Euklides,
$\vec{u} = \operatorname{proj}_{\vec{v}}(\vec{w})$ meminimumkan
jarak
$\lVert \vec{w} - \vec{u} \rVert$ di antara semua vektor $\vec{u}$ yang merupakan
kelipatan dari $\vec{v}$.
Karena itu, proyeksi ini sering muncul dalam matematika
terapan dalam berbagai konteks ketika kita tidak dapat menyelesaikan suatu persoalan
secara eksak: Kita tidak selalu dapat menyelesaikan
\myquote{Temukan $\vec{w}$ sebagai kelipatan dari
$\vec{v}$,}
tetapi $\operatorname{proj}_{\vec{v}}(\vec{w})$ merupakan \myquote{solusi} terbaik.

\begin{myfig}
\capstart
\inputpdft{vec-orthoproj}{%
Diagram yang menunjukkan vektor v berwarna abu-abu muda dan vektor w yang membentuk sudut
theta dengan v serta berpangkal pada titik yang sama, ditandai dengan garis tebal gelap. Suatu
segitiga siku-siku dengan sudut siku-siku pada garis yang diberikan oleh vektor v
ditampilkan, dan vektor lebih kecil yang bersesuaian dalam arah v serta ditandai
sebagai proj sub v dari w ditampilkan dengan garis tebal gelap.}
\caption{Proyeksi ortogonal.\label{vec-orthoproj:fig}}
\end{myfig}

Rumus tersebut mengikuti dari trigonometri dasar. Jika $\theta$ lancip, panjang
$\operatorname{proj}_{\vec{v}}(\vec{w})$ seharusnya
sebesar $\cos \theta$ kali panjang $\vec{w}$, yakni $(\cos
\theta)\lVert\vec{w}\rVert$.
Kita mengambil vektor satuan dalam arah $\vec{v}$, yakni
$\frac{\vec{v}}{\lVert \vec{v} \rVert}$, dan mengalikannya dengan panjang
proyeksi. Dengan kata lain,
\begin{equation*}
\operatorname{proj}_{\vec{v}}(\vec{w})
=
(\cos \theta) \lVert \vec{w} \rVert
\frac{\vec{v}}{\lVert \vec{v} \rVert}
=
\frac{(\cos \theta) \lVert \vec{w} \rVert \lVert \vec{v} \rVert}{
{\lVert \vec{v} \rVert}^2
}
\vec{v}
=
\frac{\langle \vec{w}, \vec{v} \rangle}{
\langle \vec{v}, \vec{v} \rangle
}
\vec{v} .
\end{equation*}
Perhitungan ini berlaku untuk $\theta$ tumpul ($\cos \theta < 0$)
karena kita juga perlu membalik arah.

\begin{example}
Misalkan kita ingin memproyeksikan vektor $(3,2,1)$ pada vektor $(1,2,3)$.
Hitung
\begin{equation*}
\begin{split}
\operatorname{proj}_{(1,2,3)} \bigl( (3,2,1) \bigr)
=
\frac{\langle (3,2,1) , (1,2,3) \rangle}{\langle (1,2,3) , (1,2,3) \rangle}
(1,2,3)
& =
\frac{3 \cdot 1 + 2 \cdot 2 + 1 \cdot 3}{ 1 \cdot 1 + 2 \cdot 2 + 3 \cdot 3}
(1,2,3)
\\
& =
\frac{10}{14}
(1,2,3)
=
\left(\frac{5}{7},\frac{10}{7},\frac{15}{7}\right) .
\end{split}
\end{equation*}

Kita periksa kembali bahwa proyeksinya ortogonal. Artinya,
$\vec{w}-\operatorname{proj}_{\vec{v}}(\vec{w})$ seharusnya ortogonal terhadap
$\vec{v}$.
Lihat sudut siku-siku dalam \figurevref{vec-orthoproj:fig}. Dengan kata lain,
\begin{equation*}
(3,2,1) - \operatorname{proj}_{(1,2,3)} \bigl( (3,2,1) \bigr)
=
\left(3-\frac{5}{7},2-\frac{10}{7},1-\frac{15}{7}\right)
=
\left(\frac{16}{7},\frac{4}{7},\frac{-8}{7}\right)
\end{equation*}
seharusnya ortogonal terhadap $(1,2,3)$. Hasil kali dalamnya harus nol:
\begin{equation*}
\left\langle
\left(\frac{16}{7},\frac{4}{7},\frac{-8}{7}\right)
,
(1,2,3)
\right\rangle
=
\frac{16}{7} \cdot 1 + \frac{4}{7} \cdot 2 -\frac{8}{7} \cdot 3
=
0 .
\end{equation*}
\end{example}

\subsection{Basis ortogonal}

Seperti yang telah kita katakan, suatu basis $\vec{v}_1,\vec{v}_2,\ldots,\vec{v}_n$
adalah \emph{\myindex{basis ortogonal}} jika semua vektor dalam
basis tersebut saling ortogonal, yakni, jika
\begin{equation*}
\langle \vec{v}_j , \vec{v}_k \rangle = 0
\end{equation*}
untuk semua pilihan $j$ dan $k$ dengan $j \neq k$ (suatu vektor taknol tidak mungkin
ortogonal terhadap dirinya sendiri).
Suatu basis selanjutnya disebut \emph{\myindex{basis ortonormal}} jika semua
vektor dalam basis tersebut juga merupakan vektor satuan, yakni, jika semua vektor
bermagnitudo 1.
Sebagai contoh, basis baku $\{ (1,0,0), (0,1,0), (0,0,1) \}$ merupakan
basis ortonormal dari ${\mathbb{R}}^3$:
Setiap pasangan ortogonal, dan setiap vektor bermagnitudo
satu.

Alasan kita tertarik pada basis ortogonal (atau ortonormal) ialah
basis tersebut membuat representasi suatu vektor (atau suatu proyeksi pada
subruang) dalam basis itu menjadi sangat sederhana. Rumus proyeksi ortogonal
pada suatu vektor memberikan koefisiennya kepada kita. Dalam
\chapterref{FS:chapter}, kita menggunakan gagasan yang sama dengan menemukan
basis ortogonal yang tepat bagi himpunan solusi suatu persamaan diferensial.
Kemudian kita dapat menemukan sebarang solusi khusus hanya dengan menerapkan
rumus proyeksi ortogonal, yang hanya terdiri atas beberapa hasil kali dalam.

Mari kita kembali ke aljabar linear. Pertimbangkan suatu subruang $S$
dan suatu basis ortogonal
$\vec{v}_1, \vec{v}_2, \ldots, \vec{v}_n$ bagi $S$.
Kita ingin
menyatakan suatu vektor $\vec{x}$ dalam basis ini.
Jika $\vec{x}$ tidak berada dalam rentang
basis tersebut (ketika $\vec{x}$ tidak berada dalam $S$),
tentu saja hal itu tidak mungkin,
tetapi rumus berikut
setidaknya memberi kita proyeksi ortogonal pada subruang $S$,
atau dengan kata lain, pendekatan terbaik
bagi $\vec{x}$ dalam subruang tersebut---vektor dalam $S$ yang paling dekat dengan $\vec{x}$.

Pertama misalkan $\vec{x}$ berada dalam rentang. Maka vektor itu merupakan jumlah dari
proyeksi-proyeksi ortogonal:
\begin{equation*}
\vec{x} =
\operatorname{proj}_{\vec{v}_1} ( \vec{x} )
+
\operatorname{proj}_{\vec{v}_2} ( \vec{x} )
+
\cdots
+
\operatorname{proj}_{\vec{v}_n} ( \vec{x} )
=
\frac{\langle \vec{x}, \vec{v}_1 \rangle}{
\langle \vec{v}_1, \vec{v}_1 \rangle
}
\vec{v}_1
+
\frac{\langle \vec{x}, \vec{v}_2 \rangle}{
\langle \vec{v}_2, \vec{v}_2 \rangle
}
\vec{v}_2
+
\cdots
+
\frac{\langle \vec{x}, \vec{v}_n \rangle}{
\langle \vec{v}_n, \vec{v}_n \rangle
}
\vec{v}_n .
\end{equation*}
Dengan kata lain, jika kita ingin menulis
$\vec{x} =
a_1 \vec{v}_1 +
a_2 \vec{v}_2 + \cdots +
a_n \vec{v}_n$, maka
\begin{equation*}
a_1 =
\frac{\langle \vec{x}, \vec{v}_1 \rangle}{
\langle \vec{v}_1, \vec{v}_1 \rangle
} , \quad
a_2 =
\frac{\langle \vec{x}, \vec{v}_2 \rangle}{
\langle \vec{v}_2, \vec{v}_2 \rangle
} , \quad \ldots , \quad
a_n =
\frac{\langle \vec{x}, \vec{v}_n \rangle}{
\langle \vec{v}_n, \vec{v}_n \rangle
} .
\end{equation*}

Cara lain untuk menurunkan rumus ini ialah bekerja secara terbalik. Misalkan
$\vec{x} =
a_1 \vec{v}_1 +
a_2 \vec{v}_2 + \cdots +
a_n \vec{v}_n$. Ambil hasil kali dalam dengan $\vec{v}_j$, dan
gunakan sifat-sifat hasil kali dalam:
\begin{equation*}
\begin{split}
\langle \vec{x} , \vec{v}_j \rangle
& =
\langle a_1 \vec{v}_1 +
a_2 \vec{v}_2 + \cdots +
a_n \vec{v}_n , \vec{v}_j \rangle
\\
& =
a_1 \langle \vec{v}_1 , \vec{v}_j \rangle +
a_2 \langle \vec{v}_2 , \vec{v}_j \rangle +
\cdots +
a_n \langle \vec{v}_n , \vec{v}_j \rangle .
\end{split}
\end{equation*}
Karena basisnya ortogonal,
$\langle \vec{v}_k , \vec{v}_j \rangle = 0$ kapan pun
$k \neq j$. Artinya, hanya satu suku, yaitu suku ke-$j$,
di ruas kanan yang taknol dan kita memperoleh
\begin{equation*}
\langle \vec{x} , \vec{v}_j \rangle
=
a_j \langle \vec{v}_j , \vec{v}_j \rangle .
\end{equation*}
Dengan menyelesaikan untuk $a_j$, kita memperoleh $a_j =
\frac{\langle \vec{x}, \vec{v}_j \rangle}{
\langle \vec{v}_j, \vec{v}_j \rangle
}$ seperti sebelumnya.

\begin{example}
Vektor-vektor $(1,1)$ dan $(1,-1)$ membentuk suatu basis ortogonal dari ${\mathbb{R}}^2$.
Misalkan kita ingin merepresentasikan $(3,4)$ dalam basis ini,
yakni, kita ingin menemukan $a_1$ dan $a_2$ sedemikian sehingga
\begin{equation*}
(3,4) = a_1 (1,1) + a_2 (1,-1) .
\end{equation*}
Kita hitung:
\begin{equation*}
a_1 =
\frac{\langle (3,4), (1,1) \rangle}{
\langle (1,1), (1,1) \rangle
}
=
\frac{7}{2}, \qquad
a_2 =
\frac{\langle (3,4), (1,-1) \rangle}{
\langle (1,-1), (1,-1) \rangle
}
=
\frac{-1}{2} .
\end{equation*}
Jadi
\begin{equation*}
(3,4) = \frac{7}{2} (1,1) + \frac{-1}{2} (1,-1) .
\end{equation*}
\end{example}

Jika basisnya ortonormal alih-alih sekadar ortogonal,
maka semua penyebutnya bernilai satu.
Mudah untuk menjadikan suatu basis ortonormal---bagi semua vektor dengan ukurannya.
Jika Anda ingin menguraikan banyak vektor,
mungkin lebih baik menemukan suatu basis ortonormal.
Dalam contoh di atas, basis ortonormal yang kita buat adalah
\begin{equation*}
\left( \frac{1}{\sqrt{2}} , \frac{1}{\sqrt{2}} \right) , \quad
\left( \frac{1}{\sqrt{2}} , \frac{-1}{\sqrt{2}} \right) .
\end{equation*}
Maka perhitungannya akan menjadi
\begin{equation*}
\begin{split}
(3,4)
& =
\left\langle
(3,4)
,
\left( \frac{1}{\sqrt{2}} , \frac{1}{\sqrt{2}} \right)
\right\rangle
\left( \frac{1}{\sqrt{2}} , \frac{1}{\sqrt{2}} \right)
+
\left\langle
(3,4)
,
\left( \frac{1}{\sqrt{2}} , \frac{-1}{\sqrt{2}} \right)
\right\rangle
\left( \frac{1}{\sqrt{2}} , \frac{-1}{\sqrt{2}} \right)
\\
& =
\frac{7}{\sqrt{2}}
\left( \frac{1}{\sqrt{2}} , \frac{1}{\sqrt{2}} \right)
+
\frac{-1}{\sqrt{2}}
\left( \frac{1}{\sqrt{2}} , \frac{-1}{\sqrt{2}} \right) .
\end{split}
\end{equation*}

Mungkin contoh ini tidak begitu mengagumkan, tetapi jika diberikan
vektor-vektor di ${\mathbb{R}}^{20}$ alih-alih ${\mathbb{R}}^2$,
tentu kita jauh lebih suka melakukan 20 hasil kali dalam
(atau 40 jika kita tidak mempunyai basis ortonormal) daripada
menyelesaikan suatu sistem berisi dua puluh persamaan dalam dua puluh peubah
dengan menggunakan reduksi baris pada matriks $20 \times 21$.

Seperti yang dikatakan di atas, rumus tersebut tetap berlaku meskipun $\vec{x}$ tidak berada dalam
subruang, meskipun kemudian rumus itu tidak menghasilkan vektor $\vec{x}$
melainkan proyeksinya. Lebih konkret, misalkan $S$ suatu subruang
yang merupakan rentang dari $\vec{v}_1,\vec{v}_2,\ldots,\vec{v}_n$ dan $\vec{x}$
sebarang vektor. Misalkan $\operatorname{proj}_{S}(\vec{x})$ adalah vektor dalam $S$
yang paling dekat dengan $\vec{x}$. Maka
\begin{equation*}
\operatorname{proj}_{S}(\vec{x}) =
\frac{\langle \vec{x}, \vec{v}_1 \rangle}{
\langle \vec{v}_1, \vec{v}_1 \rangle
}
\vec{v}_1
+
\frac{\langle \vec{x}, \vec{v}_2 \rangle}{
\langle \vec{v}_2, \vec{v}_2 \rangle
}
\vec{v}_2
+
\cdots
+
\frac{\langle \vec{x}, \vec{v}_n \rangle}{
\langle \vec{v}_n, \vec{v}_n \rangle
}
\vec{v}_n .
\end{equation*}

Tentu saja, jika $\vec{x}$ berada dalam $S$, maka $\operatorname{proj}_{S}(\vec{x}) =
\vec{x}$, karena vektor dalam $S$ yang paling dekat dengan $\vec{x}$ adalah $\vec{x}$ itu sendiri.
Namun kegunaan sebenarnya diperoleh ketika $\vec{x}$ tidak berada dalam $S$.
Dalam banyak bidang matematika terapan, kita tidak dapat menemukan solusi eksak suatu persoalan,
tetapi kita mencoba
menemukan solusi terbaik dari suatu subhimpunan kecil (subruang). Jumlah parsial dari
deret Fourier dalam \chapterref{FS:chapter} merupakan salah satu contoh. Contoh lainnya
ialah pendekatan kuadrat terkecil untuk mencocokkan suatu kurva dengan data. Contoh lain lagi
diberikan oleh
metode numerik yang paling umum digunakan untuk menyelesaikan persamaan diferensial
parsial, yaitu metode elemen hingga.

\begin{example}
Vektor-vektor $(1,2,3)$ dan $(3,0,-1)$ ortogonal, sehingga keduanya merupakan
suatu basis ortogonal dari suatu subruang $S$:
\begin{equation*}
S =
\operatorname{span} \bigl\{ (1,2,3), (3,0,-1) \bigr\} .
\end{equation*}
Mari kita temukan vektor dalam $S$ yang paling dekat dengan $(2,1,0)$. Artinya,
mari kita temukan $\operatorname{proj}_{S}\bigl((2,1,0)\bigr)$.
\begin{equation*}
\begin{split}
\operatorname{proj}_{S}\bigl((2,1,0)\bigr)
& =
\frac{\langle (2,1,0), (1,2,3) \rangle}{
\langle (1,2,3), (1,2,3) \rangle
}
(1,2,3)
+
\frac{\langle (2,1,0), (3,0,-1) \rangle}{
\langle (3,0,-1), (3,0,-1) \rangle
}
(3,0,-1)
\\
& =
\frac{2}{7}
(1,2,3)
+
\frac{3}{5}
(3,0,-1)
\\
&=
\left( \frac{73}{35} , \frac{4}{7} , \frac{9}{35} \right) .
\end{split}
\end{equation*}
\end{example}

\subsection{Proses Gram--Schmidt}

Sebelum meninggalkan basis ortogonal, mari kita catat suatu prosedur untuk membuat
basis ortogonal dari basis sebarang. Mungkin tidak sulit menemukan
basis ortogonal bagi suatu subruang berdimensi 2, tetapi bagi subruang berdimensi 20,
tugas itu tampak berat. Untungnya, proyeksi ortogonal dapat
digunakan untuk \myquote{memproyeksikan keluar} bagian-bagian vektor yang membuatnya
tidak ortogonal. Prosedur ini disebut \emph{\myindex{proses Gram--Schmidt}}.

Kita mulai dengan suatu basis vektor $\vec{v}_1,\vec{v}_2, \ldots,
\vec{v}_n$. Kita membangun suatu basis ortogonal $\vec{w}_1, \vec{w}_2,
\ldots, \vec{w}_n$ sebagai berikut.
\begin{align*}
\vec{w}_1 & = \vec{v}_1 , \displaybreak[0]\\
\vec{w}_2 & = \vec{v}_2
- \operatorname{proj}_{\vec{w}_1}(\vec{v}_2) , \displaybreak[0]\\
\vec{w}_3 & = \vec{v}_3
- \operatorname{proj}_{\vec{w}_1}(\vec{v}_3)
- \operatorname{proj}_{\vec{w}_2}(\vec{v}_3) , \displaybreak[0]\\
\vec{w}_4 & = \vec{v}_4
- \operatorname{proj}_{\vec{w}_1}(\vec{v}_4)
- \operatorname{proj}_{\vec{w}_2}(\vec{v}_4)
- \operatorname{proj}_{\vec{w}_3}(\vec{v}_4) , \\
& \vdots \\
\vec{w}_n & = \vec{v}_n
- \operatorname{proj}_{\vec{w}_1}(\vec{v}_n)
- \operatorname{proj}_{\vec{w}_2}(\vec{v}_n)
- \cdots
- \operatorname{proj}_{\vec{w}_{n-1}}(\vec{v}_n) .
\end{align*}
Yang kita lakukan pada langkah ke-$k$ ialah mengambil $\vec{v}_k$ dan
mengurangi proyeksi $\vec{v}_k$ pada subruang yang direntang oleh
$\vec{w}_1,\vec{w}_2,\ldots,\vec{w}_{k-1}$.

\begin{example}
Pertimbangkan vektor $(1,2,-1)$ dan $(0,5,-2)$, dan sebut $S$ sebagai rentang
kedua vektor tersebut. Mari kita temukan suatu basis ortogonal dari $S$
melalui proses Gram--Schmidt:
\begin{align*}
\vec{w}_1 & = (1,2,-1) , \\
\vec{w}_2 & = (0,5,-2)
- \operatorname{proj}_{(1,2,-1)}\bigl((0,5,-2)\bigr)
\\
& =
(0,5,-2) -
\frac{\langle (0,5,-2), (1,2,-1) \rangle}{
\langle (1,2,-1), (1,2,-1) \rangle
}
(1,2,-1)
=
(0,5,-2) -
2
(1,2,-1)
=
(-2,1,0) .
\end{align*}
Jadi $(1,2,-1)$ dan $(-2,1,0)$ merentang $S$ dan ortogonal. Mari kita periksa:
$\langle (1,2,-1) , (-2,1,0) \rangle = 0$.

Misalkan kita ingin menemukan suatu basis ortonormal, bukan sekadar basis ortogonal.
Kita cukup menjadikan
vektor-vektor tersebut vektor satuan dengan membaginya dengan magnitudonya. Kedua vektor
yang membentuk basis ortonormal dari $S$ adalah:
\begin{equation*}
\frac{1}{\sqrt{6}} (1,2,-1) = \left(
\frac{1}{\sqrt{6}},
\frac{2}{\sqrt{6}},
\frac{-1}{\sqrt{6}}
\right) ,
\qquad
\frac{1}{\sqrt{5}} (-2,1,0) = \left(
\frac{-2}{\sqrt{5}},
\frac{1}{\sqrt{5}},
0
\right) .
\end{equation*}
\end{example}

\subsection{Latihan}

\begin{exercise}
Temukan $s$ yang membuat vektor-vektor berikut ortogonal:
$(1,2,3)$, $(1,1,s)$.
\end{exercise}

\begin{exercise}
Temukan sudut $\theta$ antara
$(1,3,1)$ dan $(2,1,-1)$.
\end{exercise}

\begin{exercise}
Diberikan bahwa $\langle \vec{v} , \vec{w} \rangle = 3$ dan
$\langle \vec{v} , \vec{u} \rangle = -1$, hitung
\begin{tasks}(3)
\task $\langle \vec{u} , 2 \vec{v} \rangle$
\task $\langle \vec{v} , 2 \vec{w} + 3 \vec{u} \rangle$
\task $\langle \vec{w} + 3 \vec{u}, \vec{v} \rangle$
\end{tasks}
\end{exercise}

\begin{exercise}
Misalkan $\vec{v} = (1,1,-1)$. Temukan
\begin{tasks}(3)
\task $\operatorname{proj}_{\vec{v}}\bigl( (1,0,0) \bigr)$
\task $\operatorname{proj}_{\vec{v}}\bigl( (1,2,3) \bigr)$
\task $\operatorname{proj}_{\vec{v}}\bigl( (1,-1,0) \bigr)$
\end{tasks}
\end{exercise}

\begin{exercise}
\pagebreak[2]
Pertimbangkan vektor-vektor $(1,2,3)$, $(-3,0,1)$, $(1,-5,3)$.
\begin{tasks}(2)
\task Periksa bahwa vektor-vektor tersebut bebas linear dan karenanya membentuk suatu basis dari ${\mathbb{R}}^3$.
\task Periksa bahwa vektor-vektor tersebut saling ortogonal, dan karenanya
merupakan suatu basis ortogonal.
\task Representasikan $(1,1,1)$ sebagai kombinasi linear dari basis ini.
\task Jadikan basis tersebut ortonormal.
\end{tasks}
\end{exercise}

\begin{exercise}
Misalkan $S$ subruang yang direntang oleh
$(1,3,-1)$, $(1,1,1)$. Temukan suatu basis ortogonal dari $S$
dengan proses Gram--Schmidt.
\end{exercise}

\begin{exercise}
Dimulai dengan $(1,2,3)$, $(1,1,1)$, $(2,2,0)$, ikuti proses Gram--Schmidt
untuk menemukan suatu basis ortogonal dari ${\mathbb{R}}^3$.
\end{exercise}

\begin{exercise}
Temukan suatu basis ortogonal dari ${\mathbb{R}}^3$ sedemikian sehingga $(3,1,-2)$
merupakan salah satu vektornya. Petunjuk: Pertama temukan dua vektor tambahan untuk membentuk suatu
himpunan bebas linear.
\end{exercise}

\begin{exercise}
Dengan menggunakan kosinus dan sinus dari $\theta$, temukan suatu vektor satuan $\vec{u}$
di ${\mathbb{R}}^2$ yang
membentuk sudut $\theta$ dengan $\vec{\imath} = (1,0)$. Berapakah
$\langle \vec{\imath} , \vec{u} \rangle$?
\end{exercise}

\setcounter{exercise}{100}

\begin{exercise}
Temukan $s$ yang membuat vektor-vektor berikut ortogonal:
$(1,1,1)$, $(1,s,1)$.
\end{exercise}
\exsol{%
$s=-2$
}

\begin{exercise}
Temukan sudut $\theta$ antara
$(1,2,3)$ dan $(1,1,1)$.
\end{exercise}
\exsol{%
$\theta \approx 0.3876$
}

\begin{exercise}
Diberikan bahwa $\langle \vec{v} , \vec{w} \rangle = 1$,
$\langle \vec{v} , \vec{u} \rangle = -1$, dan
$\lVert \vec{v} \rVert = 3$, hitung
\begin{tasks}(3)
\task $\langle 3 \vec{u} , 5 \vec{v} \rangle$
\task $\langle \vec{v} , 2 \vec{w} + 3 \vec{u} \rangle$
\task $\langle \vec{w} + 3 \vec{v}, \vec{v} \rangle$
\end{tasks}
\end{exercise}
\exsol{%
a)~-15 \quad
b)~-1 \quad
c)~28
}

\begin{exercise}
Misalkan $\vec{v} = (1,0,-1)$. Temukan
\begin{tasks}(3)
\task $\operatorname{proj}_{\vec{v}}\bigl( (0,2,1) \bigr)$
\task $\operatorname{proj}_{\vec{v}}\bigl( (1,0,1) \bigr)$
\task $\operatorname{proj}_{\vec{v}}\bigl( (4,-1,0) \bigr)$
\end{tasks}
\end{exercise}
\exsol{%
a)~$(\nicefrac{-1}{2},0,\nicefrac{1}{2})$ \quad
b)~$(0,0,0)$ \quad
c)~$(2,0,-2)$
}

\begin{exercise}
Vektor-vektor $(1,1,-1)$, $(2,-1,1)$, $(0,1,1)$ membentuk suatu basis ortogonal.
Representasikan vektor-vektor berikut dalam basis ini:
\begin{tasks}(3)
\task $(-1,4,0)$
\task $(4,-1,3)$
\task $(-2,6,-2)$
\end{tasks}
\end{exercise}
\exsol{%
a)~$(1,1,-1)-(2,-1,1)+2(0,1,1)$ \quad
b)~$2(2,-1,1)+(0,1,1)$ \quad
c)~$2(1,1,-1)-2(2,-1,1)+2(0,1,1)$
}

\begin{exercise}
Misalkan $S$ subruang yang direntang oleh
$(2,-1,1)$, $(2,2,2)$. Temukan suatu basis ortogonal dari $S$
dengan proses Gram--Schmidt.
\end{exercise}
\exsol{%
$(2,-1,1)$, $(\nicefrac{2}{3},\nicefrac{8}{3},\nicefrac{4}{3})$
}

\begin{exercise}
Dimulai dengan $(1,1,-1)$, $(2,3,-1)$, $(1,-1,1)$, ikuti proses Gram--Schmidt
untuk menemukan suatu basis ortogonal dari ${\mathbb{R}}^3$.
\end{exercise}
\exsol{%
$(1,1,-1)$, $(0,1,1)$, $(\nicefrac{4}{3},\nicefrac{-2}{3},\nicefrac{2}{3})$
}

%%%%%%%%%%%%%%%%%%%%%%%%%%%%%%%%%%%%%%%%%%%%%%%%%%%%%%%%%%%%%%%%%%%%%%%%%%%%%%

\sectionnewpage
\section{Determinan}
\label{det:section}

\sectionnotes{1 kuliah}

Untuk matriks persegi kita mendefinisikan suatu besaran berguna yang disebut
\emph{\myindex{determinan}}. Definisikan
determinan suatu matriks $1 \times 1$ sebagai nilai satu-satunya entrinya,
\begin{equation*}
\det \left(
\begin{bmatrix}
a
\end{bmatrix}
\right)
\overset{\text{def}}{=}
a .
\end{equation*}
Untuk suatu matriks $2 \times 2$, definisikan
\begin{equation*}
\det \left(
\begin{bmatrix}
a & b \\
c & d
\end{bmatrix}
\right)
\overset{\text{def}}{=}
ad-bc .
\end{equation*}

Sebelum mendefinisikan
determinan untuk matriks yang lebih besar, kita catat
makna determinan.
Suatu matriks $n \times n$
memberikan suatu pemetaan dari ruang Euklides berdimensi $n$, yaitu ${\mathbb{R}}^n$, ke
dirinya sendiri.
Jadi suatu matriks $2 \times 2$ berupa $A$ merupakan pemetaan dari
bidang ke dirinya sendiri. Determinan
$A$ adalah faktor perubahan luas objek.
Jika kita mengambil persegi satuan (persegi dengan sisi 1) pada bidang, maka
$A$ memetakan persegi itu ke suatu jajaran genjang seluas $\lvert\det(A)\rvert$. Tanda
dari $\det(A)$ menunjukkan perubahan orientasi (negatif jika sumbu-sumbunya terbalik). Sebagai
contoh, misalkan
\begin{equation*}
A =
\begin{bmatrix}
1 & 1 \\
-1 & 1
\end{bmatrix} .
\end{equation*}
Maka $\det(A) = 1+1 = 2$.
Mari kita lihat ke mana $A$ memetakan persegi satuan---persegi dengan titik-titik sudut
$(0,0)$, $(1,0)$, $(1,1)$, dan $(0,1)$.
Titik $(0,0)$ dipetakan
ke $(0,0)$.
\begin{equation*}
\begin{bmatrix}
1 & 1 \\
-1 & 1
\end{bmatrix}
\begin{bmatrix}
1 \\ 0
\end{bmatrix} =
\begin{bmatrix}
1 \\
-1
\end{bmatrix}
,
\qquad
\begin{bmatrix}
1 & 1 \\
-1 & 1
\end{bmatrix}
\begin{bmatrix}
1 \\ 1
\end{bmatrix} =
\begin{bmatrix}
2 \\
0
\end{bmatrix}
,
\qquad
\begin{bmatrix}
1 & 1 \\
-1 & 1
\end{bmatrix}
\begin{bmatrix}
0 \\ 1
\end{bmatrix} =
\begin{bmatrix}
1 \\
1
\end{bmatrix}
.
\end{equation*}
Citra persegi tersebut adalah persegi lain dengan titik-titik sudut $(0,0)$, $(1,-1)$,
$(2,0)$, dan $(1,1)$. Persegi
citra itu mempunyai
sisi sepanjang $\sqrt{2}$, sehingga luasnya 2. Lihat
\figurevref{linalg-imagesquare:fig}.

\begin{myfig}
\capstart
\inputpdft{linalg-imagesquare}{%
Diberikan dua diagram bidang. Pada diagram kiri,
terdapat persegi dengan titik-titik sudut (0,0), (1,0), (1,1), dan (0,1).
Pada diagram kanan, terdapat persegi dengan titik-titik sudut yang bersesuaian
(0,0), (1, minus 1), (2,0), dan (1,1), dengan kesesuaian tersebut ditunjukkan
oleh anak panah.}
\caption{Citra persegi satuan melalui pemetaan
$A$.\label{linalg-imagesquare:fig}}
\end{myfig}

Secara umum, citra suatu persegi merupakan suatu \myindex{jajaran genjang}.
Dalam geometri sekolah menengah, Anda mungkin telah melihat suatu rumus untuk
menghitung luas suatu \myindex{jajaran genjang}
dengan titik-titik sudut $(0,0)$, $(a,c)$, $(b,d)$,
dan $(a+b,c+d)$. Luasnya adalah
\begin{equation*}
\left\lvert \, \det \left(
\begin{bmatrix} a & b \\ c & d \end{bmatrix}
\right) \, \right\rvert
=
\lvert
a d - b c
\rvert
.
\end{equation*}
Garis-garis vertikal di atas berarti nilai mutlak.
Matriks $\left[ \begin{smallmatrix} a & b \\ c & d \end{smallmatrix}
\right]$
memetakan persegi satuan ke jajaran genjang yang diberikan.

\medskip

Ada beberapa cara untuk mendefinisikan determinan suatu matriks $n \times n$.
Mari kita gunakan sesuatu yang disebut \emph{\myindex{ekspansi kofaktor}}.
Kita mendefinisikan $A_{ij}$ sebagai
matriks $A$ dengan baris ke-$i$ dan kolom ke-$j$
dihapus. Sebagai contoh,
\begin{equation*}
\text{jika} \qquad
A =
\begin{bmatrix}
1 & 2 & 3 \\
4 & 5 & 6 \\
7 & 8 & 9
\end{bmatrix} ,
\qquad
\text{maka}
\qquad
A_{12} =
\begin{bmatrix}
4 & 6 \\
7 & 9
\end{bmatrix}
\qquad
\text{dan}
\qquad
A_{23} =
\begin{bmatrix}
1 & 2 \\
7 & 8
\end{bmatrix} .
\end{equation*}
Sekarang kita mendefinisikan determinan secara rekursif,
\begin{equation*}
\det (A)
\overset{\text{def}}{=}
\sum_{j=1}^n
{(-1)}^{1+j}
a_{1j} \det (A_{1j}) ,
\end{equation*}
atau dengan kata lain,
\begin{equation*}
\det (A) =
a_{11} \det (A_{11}) -
a_{12} \det (A_{12}) +
a_{13} \det (A_{13}) -
\cdots
\begin{cases}
+ a_{1n} \det (A_{1n}) & \text{jika } n \text{ ganjil,} \\
- a_{1n} \det (A_{1n}) & \text{jika } n \text{ genap.}
\end{cases}
\end{equation*}
Untuk suatu matriks $3 \times 3$,
kita memperoleh $\det (A) = a_{11} \det (A_{11}) -
a_{12} \det (A_{12}) + a_{13} \det (A_{13})$. Sebagai contoh,
\begin{equation*}
\begin{split}
\det \left(
\begin{bmatrix}
1 & 2 & 3 \\
4 & 5 & 6 \\
7 & 8 & 9
\end{bmatrix}
\right)
& =
1 \cdot
\det \left(
\begin{bmatrix}
5 & 6 \\
8 & 9
\end{bmatrix}
\right)
-
2 \cdot
\det \left(
\begin{bmatrix}
4 & 6 \\
7 & 9
\end{bmatrix}
\right)
+
3 \cdot
\det \left(
\begin{bmatrix}
4 & 5 \\
7 & 8
\end{bmatrix}
\right) \\
& =
1 (5 \cdot 9 - 6 \cdot 8)
-
2 (4 \cdot 9 - 6 \cdot 7)
+
3 (4 \cdot 8 - 5 \cdot 7)
= 0 .
\end{split}
\end{equation*}

Ternyata kita tidak harus menggunakan baris pertama.
Artinya, untuk setiap $i$,
\begin{equation*}
\det (A)
=
\sum_{j=1}^n
{(-1)}^{i+j}
a_{ij} \det (A_{ij}) .
\end{equation*}
Kadang-kadang berguna untuk menggunakan baris selain baris pertama. Dalam contoh berikut,
lebih nyaman melakukan ekspansi sepanjang baris kedua. Perhatikan bahwa untuk
baris kedua kita mulai dengan tanda negatif.
\begin{equation*}
\begin{split}
\det \left(
\begin{bmatrix}
1 & 2 & 3 \\
0 & 5 & 0 \\
7 & 8 & 9
\end{bmatrix}
\right)
& =
- 0 \cdot
\det \left(
\begin{bmatrix}
2 & 3 \\
8 & 9
\end{bmatrix}
\right)
+
5 \cdot
\det \left(
\begin{bmatrix}
1 & 3 \\
7 & 9
\end{bmatrix}
\right)
-
0 \cdot
\det \left(
\begin{bmatrix}
1 & 2 \\
7 & 8
\end{bmatrix}
\right) \\
& =
0
+
5 (1 \cdot 9 - 3 \cdot 7)
+
0
= -60 .
\end{split}
\end{equation*}
Mari kita periksa apakah hasilnya sungguh sama dengan ekspansi sepanjang baris pertama,
\begin{equation*}
\begin{split}
\det \left(
\begin{bmatrix}
1 & 2 & 3 \\
0 & 5 & 0 \\
7 & 8 & 9
\end{bmatrix}
\right)
& =
1 \cdot
\det \left(
\begin{bmatrix}
5 & 0 \\
8 & 9
\end{bmatrix}
\right)
-
2 \cdot
\det \left(
\begin{bmatrix}
0 & 0 \\
7 & 9
\end{bmatrix}
\right)
+
3 \cdot
\det \left(
\begin{bmatrix}
0 & 5 \\
7 & 8
\end{bmatrix}
\right) \\
& =
1 (5 \cdot 9 - 0 \cdot 8)
-
2 (0 \cdot 9 - 0 \cdot 7)
+
3 (0 \cdot 8 - 5 \cdot 7)
= -60 .
\end{split}
\end{equation*}



Dalam menghitung determinan,
kita secara bergantian menjumlahkan dan mengurangkan determinan submatriks
$A_{ij}$ yang dikalikan dengan $a_{ij}$ untuk suatu $i$ tetap dan semua $j$.
Bilangan ${(-1)}^{i+j}\det(A_{ij})$ disebut
\emph{kofaktor\index{kofaktor}}
dari matriks tersebut. Itulah sebabnya
metode perhitungan determinan ini disebut
\emph{ekspansi kofaktor}.

Demikian pula, kita tidak perlu melakukan ekspansi sepanjang suatu baris; kita dapat melakukan ekspansi
sepanjang suatu kolom. Untuk setiap $j$,
\begin{equation*}
\det (A)
=
\sum_{i=1}^n
{(-1)}^{i+j}
a_{ij} \det (A_{ij}) .
\end{equation*}
Fakta terkait ialah
\begin{equation*}
\det (A) = \det (A^T) .
\end{equation*}

\medskip

Suatu matriks disebut \emph{\myindex{segitiga atas}} jika semua elemen di bawah
diagonal utama bernilai 0. Sebagai contoh,
\begin{equation*}
\begin{bmatrix}
1 & 2 & 3 \\
0 & 5 & 6 \\
0 & 0 & 9
\end{bmatrix}
\end{equation*}
adalah segitiga atas. Demikian pula, suatu matriks \emph{\myindex{segitiga bawah}}
adalah matriks yang semua elemen di atas diagonalnya bernilai nol. Sebagai contoh,
\begin{equation*}
\begin{bmatrix}
1 & 0 & 0 \\
4 & 5 & 0 \\
7 & 8 & 9
\end{bmatrix} .
\end{equation*}

Determinan matriks segitiga sangat mudah dihitung.
Pertimbangkan matriks segitiga bawah. Jika kita melakukan ekspansi sepanjang
baris pertama, kita menemukan bahwa determinannya adalah 1 kali determinan
dari matriks segitiga bawah $\left[ \begin{smallmatrix} 5 & 0 \\ 8 & 9
\end{smallmatrix} \right]$. Jadi determinannya hanyalah
hasil kali entri-entri diagonal:
\begin{equation*}
\det \left(
\begin{bmatrix}
1 & 0 & 0 \\
4 & 5 & 0 \\
7 & 8 & 9
\end{bmatrix}
\right)
=
1 \cdot 5 \cdot 9 = 45 .
\end{equation*}
Demikian pula untuk matriks segitiga atas,
\begin{equation*}
\det \left(
\begin{bmatrix}
1 & 2 & 3 \\
0 & 5 & 6 \\
0 & 0 & 9
\end{bmatrix}
\right)
=
1 \cdot 5 \cdot 9 = 45 .
\end{equation*}
Secara umum, jika $A$ segitiga, maka
\begin{equation*}
\det (A) = a_{11} a_{22} \cdots a_{nn} .
\end{equation*}

Jika $A$ diagonal, maka $A$ juga segitiga (atas dan bawah), sehingga
rumus yang sama berlaku. Sebagai contoh,
\begin{equation*}
\det \left(
\begin{bmatrix}
2 & 0 & 0 \\
0 & 3 & 0 \\
0 & 0 & 5
\end{bmatrix}
\right)
=
2 \cdot 3 \cdot 5 = 30 .
\end{equation*}

Secara khusus, matriks identitas $I$ bersifat diagonal, dan semua entri diagonalnya
bernilai 1. Jadi,
\begin{equation*}
\det(I) = 1 .
\end{equation*}

\medskip

Determinan memberi tahu Anda bagaimana objek geometris berskala.
Jika $B$ menggandakan ukuran objek geometris dan $A$ melipatgandakannya tiga kali,
maka $AB$ (yang menerapkan $B$ pada suatu objek lalu menerapkan $A$) seharusnya membuat ukurannya
bertambah dengan faktor $6$. Hal ini berlaku secara umum:

\begin{theorem}
\begin{equation*}
\det(AB) = \det(A)\det(B) .
\end{equation*}
\end{theorem}

Sifat ini merupakan salah satu sifat paling berguna dan sering digunakan untuk
sungguh-sungguh menghitung determinan. Salah satu akibat yang sangat menarik ialah
apa artinya bagi keberadaan invers.
Ambil $A$ dan $B$ sebagai invers, yakni, $AB=I$. Maka
\begin{equation*}
\det(A)\det(B) = \det(AB) = \det(I) = 1 .
\end{equation*}
Baik $\det(A)$ maupun $\det(B)$ tidak mungkin nol.
Fakta ini merupakan sifat determinan yang sangat berguna, dan
sering digunakan dalam buku ini:

\begin{theorem}
Suatu matriks $n \times n$ berupa $A$ dapat diinverskan jika dan hanya jika $\det (A) \neq 0$.
\end{theorem}

Sesungguhnya, $\det(A^{-1}) \det(A) = 1$ menyatakan bahwa
\begin{equation*}
\det(A^{-1}) =
\frac{1}{\det(A)}.
\end{equation*}
Jadi kita mengetahui determinan dari $A^{-1}$
tanpa menghitung $A^{-1}$.

Mari kita kembali ke rumus invers suatu matriks $2 \times 2$:
\begin{equation*}
\begin{bmatrix}
a & b \\
c & d
\end{bmatrix}^{-1}
=
\frac{1}{ad-bc}
\begin{bmatrix}
d & -b \\
-c & a
\end{bmatrix} .
\end{equation*}
Perhatikan determinan matriks
$[\begin{smallmatrix}a&b\\c&d\end{smallmatrix}]$
di penyebut pecahan tersebut.
Rumus itu hanya berlaku jika determinannya taknol.
Jika tidak, kita membagi dengan nol.

\medskip

Notasi umum bagi determinan ialah sepasang garis
vertikal:
\begin{equation*}
\begin{vmatrix}
a & b \\
c & d
\end{vmatrix}
=
\det \left(
\begin{bmatrix}
a & b \\
c & d
\end{bmatrix}
\right) .
\end{equation*}
Secara pribadi, saya menganggap notasi ini membingungkan karena garis vertikal biasanya
menyatakan besaran positif, sedangkan determinan dapat bernilai negatif. Pikirkan pula
bagaimana menuliskan nilai mutlak suatu determinan.
Notasi ini tidak digunakan dalam buku ini.

\subsection{Latihan}

\begin{exercise}
Hitung determinan matriks-matriks berikut:
\begin{tasks}(4)
\task
$\begin{bmatrix}
3
\end{bmatrix}$
\task
$\begin{bmatrix}
1 & 3 \\
2 & 1
\end{bmatrix}$
\task
$\begin{bmatrix}
2 & 1 \\
4 & 2
\end{bmatrix}$
\task
$\begin{bmatrix}
1 & 2 & 3 \\
0 & 4 & 5 \\
0 & 0 & 6
\end{bmatrix}$
\task
$\begin{bmatrix}
2 & 1 & 0 \\
-2 & 7 & -3 \\
0 & 2 & 0
\end{bmatrix}$
\task
$\begin{bmatrix}
2 & 1 & 3 \\
8 & 6 & 3 \\
7 & 9 & 7
\end{bmatrix}$
\task
$\begin{bmatrix}
0 & 2 & 5 & 7 \\
0 & 0 & 2 & -3 \\
3 & 4 & 5 & 7 \\
0 & 0 & 2 & 4
\end{bmatrix}$
\task
$\begin{bmatrix}
0 &  1 &  2 &  0 \\
1 &  1 & -1 & 2 \\
1 &  1 &  2 & 1 \\
2 & -1 & -2 & 3
\end{bmatrix}$
\end{tasks}
\end{exercise}

\begin{exercise}
Untuk nilai $x$ mana matriks-matriks berikut singular (tidak dapat diinverskan)?
\begin{tasks}(4)
\task
$\begin{bmatrix}
2 & 3 \\
2 & x
\end{bmatrix}$
\task
$\begin{bmatrix}
2 & x \\
1 & 2
\end{bmatrix}$
\task
$\begin{bmatrix}
x & 1 \\
4 & x
\end{bmatrix}$
\task
$\begin{bmatrix}
x & 0 & 1 \\
1 & 4 & 2 \\
1 & 6 & 2
\end{bmatrix}$
\end{tasks}
\end{exercise}

\begin{exercise}
Hitung
\begin{equation*}
\det \left( \begin{bmatrix}
2 & 1 & 2 & 3 \\
0 & 8 & 6 & 5 \\
0 & 0 & 3 & 9 \\
0 & 0 & 0 & 1
\end{bmatrix}^{-1}
\right)
\end{equation*}
tanpa menghitung inversnya.
\end{exercise}

\begin{exercise}
Misalkan
\begin{equation*}
L = \begin{bmatrix}
1 & 0 & 0 & 0 \\
2 & 1 & 0 & 0 \\
7 & \pi & 1 & 0 \\
2^8 & 5 & -99 & 1
\end{bmatrix}
\qquad \text{dan} \qquad
U = \begin{bmatrix}
5 & 9 & 1 & -\sin(1) \\
0 & 1 & 88 & -1 \\
0 & 0 & 1 & 3 \\
0 & 0 & 0 & 1
\end{bmatrix} .
\end{equation*}
Misalkan $A = LU$. Hitung $\det(A)$ dengan cara sederhana, tanpa menghitung $A$.
Petunjuk: Pertama, baca langsung $\det(L)$ dan $\det(U)$.
\end{exercise}

\begin{exercise}
Pertimbangkan pemetaan linear dari ${\mathbb R}^2$ ke ${\mathbb R}^2$
yang diberikan oleh matriks
$A = \left[ \begin{smallmatrix}
1 & x \\
2 & 1
\end{smallmatrix} \right]$
untuk suatu bilangan $x$. Anda ingin membuat $A$ sedemikian sehingga menggandakan luas
setiap bangun geometris. Apa saja kemungkinan nilai $x$ (terdapat dua
jawaban)?
\end{exercise}

\begin{exercise}
Misalkan $A$ dan $S$ merupakan matriks $n \times n$, dan $S$ dapat diinverskan.
Misalkan $\det(A) = 3$. Hitung $\det(S^{-1}AS)$ dan
$\det(SAS^{-1})$. Berikan alasan bagi jawaban Anda menggunakan teorema-teorema dalam bagian ini.
\end{exercise}

\begin{exercise}
Misalkan $A$ suatu matriks $n \times n$ sedemikian sehingga $\det(A)=1$.
Hitung $\det(x A)$ dengan diberikan suatu bilangan $x$.\linebreak[2]
Petunjuk: Pertama coba hitung
$\det(xI)$, kemudian perhatikan bahwa $xA = (xI)A$.
\end{exercise}

\setcounter{exercise}{100}

\begin{exercise}
\pagebreak[2]
Hitung determinan matriks-matriks berikut:
\begin{tasks}(4)
\task
$\begin{bmatrix}
-2
\end{bmatrix}$
\task
$\begin{bmatrix}
2 & -2 \\
1 & 3
\end{bmatrix}$
\task
$\begin{bmatrix}
2 & 2 \\
2 & 2
\end{bmatrix}$
\task
$\begin{bmatrix}
2 & 9 & -11 \\
0 & -1 & 5 \\
0 & 0 & 3
\end{bmatrix}$
\task
$\begin{bmatrix}
2 & 1 & 0 \\
-2 & 7 & 3 \\
1 & 1 & 0
\end{bmatrix}$
\task
$\begin{bmatrix}
5 & 1 & 3 \\
4 & 1 & 1 \\
4 & 5 & 1
\end{bmatrix}$
\task
$\begin{bmatrix}
3 & 2 & 5 & 7 \\
0 & 0 & 2 & 0 \\
0 & 4 & 5 & 0 \\
2 & 1 & 2 & 4
\end{bmatrix}$
\task
$\begin{bmatrix}
 0 &  2 &  1 &  0 \\
 1 &  2 & -3 &  4 \\
 5 &  6 & -7 &  8 \\
 1 &  2 &  3 & -2
\end{bmatrix}$
\end{tasks}
\end{exercise}
\exsol{%
a)~$-2$
\quad b)~$8$
\quad c)~$0$
\quad d)~$-6$
\quad e)~$-3$
\quad f)~$28$
\quad g)~$16$
\quad h)~$-24$
}

\begin{exercise}
Untuk nilai $x$ mana matriks-matriks berikut singular (tidak dapat diinverskan)?
\begin{tasks}(4)
\task
$\begin{bmatrix}
1 & 3 \\
1 & x
\end{bmatrix}$
\task
$\begin{bmatrix}
3 & x \\
1 & 3
\end{bmatrix}$
\task
$\begin{bmatrix}
x & 3 \\
3 & x
\end{bmatrix}$
\task
$\begin{bmatrix}
x & 1 & 0 \\
1 & 4 & 0 \\
1 & 6 & 2
\end{bmatrix}$
\end{tasks}
\end{exercise}
\exsol{%
a)~$3$
\quad b)~$9$
\quad c)~$\pm 3$
\quad d)~$\nicefrac{1}{4}$
}

\begin{exercise}
Hitung
\begin{equation*}
\det \left( \begin{bmatrix}
3 & 4 & 7 & 12 \\
0 & -1 & 9 & -8 \\
0 & 0 & -2 & 4 \\
0 & 0 & 0 & 2
\end{bmatrix}^{-1}
\right)
\end{equation*}
tanpa menghitung inversnya.
\end{exercise}
\exsol{%
$\nicefrac{1}{12}$
}

\begin{exercise}[menantang]
Temukan semua $x$ yang membuat invers matriks
\begin{equation*}
\begin{bmatrix}
1 & 2 \\ 1 & x
\end{bmatrix}^{-1}
\end{equation*}
hanya mempunyai entri bilangan bulat (tanpa pecahan).
Perhatikan bahwa terdapat dua jawaban.
\end{exercise}
\exsol{%
$1$ dan $3$
}
